\documentclass[11pt]{article}
\usepackage{fontspec}
\setmainfont{Times New Roman}
\usepackage{amsmath,amssymb}
\usepackage{geometry}
\usepackage{microtype}
\usepackage{hyperref}
\usepackage{mathrsfs}
\newcommand{\srcfn}[2]{\footnote{#2}}
\providecommand{\tightlist}{\setlength{\itemsep}{0pt}\setlength{\parskip}{0pt}}
\allowdisplaybreaks

\geometry{a4paper,margin=2.1cm}
\setlength{\parindent}{1.1em}
\setlength{\parskip}{0pt}
\emergencystretch=10em
\sloppy
\hbadness=10000
\raggedbottom
\hypersetup{hidelinks}
\newcommand{\foreign}[1]{#1}
\newcommand{\tocline}[2]{\noindent #1\dotfill #2\par}
\newcommand{\tocsec}[3]{\noindent\hspace*{1.2em}\S\,#1.\quad #2\dotfill #3\par}

\begin{document}
\thispagestyle{empty}

\begin{center}
{\LARGE\bfseries Алгебра хиперкомплексных величин\par}
\vspace{3.5em}
{\large Лекција проф. E. \foreign{Noether}\par}
{\large Зимны семестр 1929/30\par}
\vspace{1.5em}
{\large Обработал проф. M. \foreign{Deuring}\par}
\vspace{4em}
{\Large\bfseries Содржанье\par}
\end{center}
\vspace{1em}

\begingroup
\small
\setlength{\parindent}{0pt}
\tocline{\textbf{Увод}}{3}
\medskip
\tocline{\textbf{Глава I: Представјеньа}}{3}
\tocsec{1}{Дефиниција директного представјеньа --- дефиниција реципрочного представјеньа}{3}
\tocsec{2}{Класы представјень}{4}
\tocsec{3}{Модулы представјень}{4}
\tocsec{4}{Связ меджу модулами представјень и представјеньами}{5}
\medskip
\tocline{\textbf{Глава II: Галоисова теорија комутативных поль}}{7}
\tocsec{5}{Разширјенье коефициентного польа хиперкомплексных системов}{7}
\tocsec{6}{Иредуцибилне представјеньа комутативных системов}{8}
\tocsec{7}{Изоморфизмы польа}{8}
\tocsec{8}{Разпадно полье}{9}
\tocsec{9}{Изоморфизмы од $Z_1$ на $Z_i$}{9}
\tocsec{10}{Галоисова група}{9}
\tocsec{11}{Главны теорем Галоисовој теорије}{9}
\tocsec{12}{Формално значенье компонентов $e_i$}{11}
\tocsec{13}{Обчи теорем о продолженьу}{12}
\medskip
\tocline{\textbf{Глава III: Абелове групы}}{13}
\tocsec{14}{Групово колцо}{13}
\tocsec{15}{Групове колца абеловых груп}{14}
\tocsec{16}{Релације карактеров}{15}
\tocsec{17}{Галоисова теорија абеловых груп}{15}
\medskip
\tocline{\textbf{Глава IV: Двустранно просте колца}}{18}
\tocsec{18}{Помочна лема}{18}
\tocsec{19}{Представјеньа двустранно простых хиперкомплексных системов в некомутативных телах, кторе разширјајут јих коефициентне польа}{19}
\tocsec{20}{Некомутативне тела}{22}
\tocsec{21}{Галоисова теорија некомутативных тел}{26}
\clearpage
\tocline{\textbf{Глава V: Факторне системы}}{29}
\tocsec{22}{Група тел с даным центром}{29}
\tocsec{23}{Факторне системы}{30}
\tocsec{24}{Множенье факторных системов}{33}
\tocsec{25}{Нормално представјенье $\mathfrak K_r$ с Галоисовыми максималными комутативными подпольами}{39}
\tocsec{26}{Множенье скреженых представјень}{43}
\medskip
\tocline{\textbf{Глава VI: Теорија скреженых продуктов}}{44}
\tocsec{27}{Представјеньа $\mathfrak R_r$ како скрежене продукты}{44}
\tocsec{28}{Продуктны теорем за факторне системы}{47}
\tocsec{29}{Теорем о главном роду в минималном}{50}
\tocsec{30}{Специализација на цикличне разпадне польа}{50}
\tocsec{31}{Применьеньа цикличного специалного случаја}{52}
\endgroup

\clearpage
\section*{Увод}\label{einleitung}

Та лекција обрабја обчу теорију представјень колц, ктора взникла с
једној страны из теорије представјень конечных груп, а с другој страны
из алгебраичној и аритметичној теорије хиперкомплексных числовых
системов. То обчејше разуменье теорије представјень имаје следне
предности над досегашньим приступом (\foreign{Frobenius},
\foreign{Schur}):

Понеже разматрјамо не саму групу, але групово колцо, а обчејше
хиперкомплексны систем или либовольно колцо, можно јест применити
теорију колц и идеалов и тым освободити теорију од многих нејасных
изчислень. То применьенье обчеј теорије колц и идеалов не само
упрошчаје теорију, але такоже води ју далье --- напр. в пытаньу о
одношеньах групы с јеј подгрупами --- и отвара јеј нове области
применьеньа. Помоцју обчеј теорије представјень можно дати ново, врло
елегантно основанье Галоисовој теорији и --- можда јешче важнејше ---
установити Галоисову теорију некомутативных тел.

Основне појмы теорије груп, колц и идеалов и обче теоремы из теорије
представјень можно најдти в делу господичны \foreign{Noether}
\foreign{Hyperkomplexe Größen und Darstellungstheorie},
\foreign{Mathematische Zeitschrift} 30, стр.~641. То дело, цитовано
како \foreign{M.Z.}, треба сматрјати дополньеньем к тому изложеньу.
Исты материал јест такоже в записах лекције господичны
\foreign{Noether} \foreign{Über Hyperkomplexe Größen und
Gruppencharaktere}, зимны семестр 1927/28.

Појмы представјенье и модул представјеньа сут тут кратко објасньене,
да бы ввести реципрочны модул представјеньа и реципрочна
представјеньа, кторе сут нове само формално.

\section*{Глава I. Представјеньа}\label{kapitel-i.-darstellungen}

$\mathfrak{o}$ нехај буде колцо --- колцо, кторе се представльа --- а
$\mathsf T$ друго колцо --- то, в ктором се представльа.

% BEGIN BOOK_S01
\subsection*{§ 1. Дефиниција премого представјеньа}\label{definition-der-direkten-darstellung}

\emph{Премо представјенье} \(n\)-того ступеньа од \(\mathfrak{o}\) в \(\mathsf T\) јест подколцо \(\mathfrak D\) колца \(\mathsf T_n\) матриц \(n\)-того ступеньа с елементами из \(\mathsf T\), на кторо јест \(\mathfrak{o}\) образовано колцовым хомоморфизмом. В знаках

\[
\mathfrak{o}\to\mathfrak D\subseteq\mathsf T_n.
\]

\subsection*{Дефиниција реципрочного представјеньа}\label{definition-der-reziproken-darstellung}

\emph{Реципрочно представјенье} \(n\)-того ступеньа од \(\mathfrak{o}\) в \(\mathsf T^*\) јест подколцо \(\mathfrak D^*\) колца \(\mathsf T_n^*\) матриц \(n\)-того ступеньа с елементами из \(\mathsf T^*\), на кторо јест \(\mathfrak{o}\) образовано реципрочным колцовым хомоморфизмом. Реципрочны колцовы хомоморфизм значи следујуче: ако елементам из \(\mathfrak{o}\), \(c,d\), одповедајут елементи из \(\mathfrak D^*\), \(c^*,d^*\), то елементу из \(\mathfrak{o}\), \(cd\), треба одповедати елемент из \(\mathfrak D^*\), \(d^*\cdot c^*\), а елементу из \(\mathfrak{o}\), \(c+d\), елемент из \(\mathfrak D^*\), \(c^*+d^*\). В знаках

\[
\mathfrak{o}\nrightarrow\mathfrak D^*\subseteq\mathsf T_n^*.
\]

Како пример појма реципрочного хомоморфизма: ако матрицам реда \(n\) над комутативным польем прирядити јих транспоноване матрице, то приряđенье јест реципрочны изоморфизм колц.

Обче можно к либовольному колцу \(\mathsf T\) конструовати реципрочно изоморфно колцо \(\mathsf T^*\) следујучим начином:

Всакому елементу из \(\mathsf T\), \(c\), прирядимо симбол \(c^*\) и дефинирујемо сложенье тых симболов формулом \(c^*+d^*=(c+d)^*\), а множенье формулом

\[
c^*\cdot d^*=(dc)^*.
\]

Множство симболов \(c^*\) јест колцо, реципрочно изоморфно с \(\mathsf T\), и означујемо је \(\mathsf T^*\).
% END BOOK_S01

% BEGIN BOOK_S02
\subsection*{§ 2. Класы представјень}\label{darstellungsklassen}

Премо представјенье \(\mathfrak D\) од \(\mathfrak{o}\) в \(\mathsf T\) можно јест претворити в изоморфно представјенье \(\mathfrak D'\) тако, же се оно по елементах трансформује регуларној матрицоју \(P\), т.~ј. елементу из \(\mathfrak{o}\), \(c\), вместо \(C\subseteq\mathfrak D\) приряди се \(P^{-1}CP\subseteq P^{-1}\mathfrak D P\).

Два представјеньа, кторе можно тако трансформовати једно в друго, называјут се \emph{еквивалентными}.

Вса представјеньа, еквивалентне једному представјеньу, творят \emph{класу представјень} (к. п.).

Тако само дефинујеме еквивалентност реципрочных представјень; два реципрочна представјеньа называјут се еквивалентными, ако једно можно достати из другого трансформацијоју регуларној матрицоју. Вса представјеньа, еквивалентне једному установјеному реципрочному представјеньу, творят \emph{класу реципрочных представјень}.
% END BOOK_S02

% BEGIN BOOK_S03
\subsection*{\texorpdfstring{\(\S\) 3. Модулы представјень}{\textbackslash С 3. Модулы представјень}}\label{s-3.-darstellungsmoduln}

Дефиниција \emph{премого модула представјеньа} (п. м. п.).

Премы модул представјеньа од \(\mathfrak{o}\) взходно к \(\mathsf T\) јест всакы двојны \(\mathfrak{o}\)-\(\mathsf T\)-модул \(\mathfrak M\) (т.~ј. адитивно записана Абелова група, ктора јест левым модулом взходно к \(\mathfrak{o}\), правым модулом взходно к \(\mathsf T\) и задовольава смјешаны закон асоциативности: \(c\in\mathfrak{o}\), \(m\in\mathfrak M\), \(\tau\in\mathsf T\); \((cm)\tau=c(m\tau)\)), кторы можно записати како прему суму конечного числа једночленных \(\mathsf T\)-модулов --- причем \(\tau=0\) следује из \(x_i\tau=0\) ---

\[
\mathfrak M=x_1\cdot\mathsf T+x_2\cdot\mathsf T+\cdots+x_n\cdot\mathsf T.
\]

Дефиниција \emph{реципрочного модула представјеньа} (р. м. п.).

Реципрочны модул представјеньа од \(\mathfrak{o}\) взходно к \(\mathsf T^*\) јест всакы левы \(\mathfrak{o}\)-модул и левы \(\mathsf T^*\)-модул \(\mathfrak M^*\), кторы задовольава закон комутативности (ако \(c\in\mathfrak{o}\), \(m\in\mathfrak M^*\), \(\tau\in\mathsf T^*\), то \(c(\tau^*m^*)=\tau^*(cm)\)) и кторы можно записати како прему суму конечного числа простых \(\mathsf T^*\)-модулов --- причем \(\tau^*=0\) следује из \(\tau^*\cdot x_i^*=0\) ---

\[
\mathfrak M^*=\mathsf T^*\cdot x_1+\cdots+\mathsf T^*\cdot x_n.
\]
% END BOOK_S03

% BEGIN BOOK_S04
\subsection*{§ 4. Связ меджу модулами представјень и представјеньами}\label{der-zusammenhang-zwischen-darstellungsmoduln-und-darstellungen}

\textbf{1.} \emph{Всакы модул представјеньа једнозначно води к класе представјень (п. м. п. к премој класе представјень, р. м. п. к реципрочној класе представјень).}

И то следујучим начином:

Множенье премого модула представјеньа \(\mathfrak M\) елементом из \(\mathfrak{o}\), \(c\), значи хомоморфизм од \(\mathfrak M\) на себе самого, кторы заради закона асоциативности \((cm)\tau=c(m\tau)\) допушча елементы из \(\mathsf T\) како операторы. Оно јест теды \emph{линеарна трансформација} од \(\mathfrak M\) в себе, т.~ј. приряđенье \(m\mapsto cm=m'\) имаје својство:

\[
\text{Ako } y_i\mapsto y_i'(=cy_i) \text{ jest, to takože }
\sum y_i\tau_i\mapsto \sum y_i'\tau_i\left(=c\sum y_i\tau_i\right).
\]

Дефинујмо суму \(\mathfrak T_1+\mathfrak T_2\) двох линеарных трансформациј

\[
\mathfrak T_1:\ m\mapsto m'(=c_1m)
\quad\text{i}\quad
\mathfrak T_2:\ m\mapsto m''(=c_2m)
\]

како трансформацију \(m\mapsto m'+m''(=(c_1+c_2)m)\), а јих продукт \(\mathfrak T_1\cdot\mathfrak T_2\) како сложену трансформацију \(m\mapsto (m'')'=c_1c_2m\). Тада множство различных линеарных трансформациј од \(\mathfrak M\), кторе стварјајут елементи из \(\mathfrak{o}\), јест колцово хомоморфны образ \(\overline{\mathfrak D}\) од \(\mathfrak{o}\).

Ныне изберимо либовольну базу \(x_1,\ldots,x_n\) модула \(\mathfrak M\). Линеарну трансформацију, створјену елементом \(c\), од \(\mathfrak M\) очивидно јест достаточно задати тым, како се трансформујут базове елементи \(x_i\), то јест познаньем матрице \((\gamma_{ij})=C\) система једначень

\[
x'_j=cx_j=\sum x_i\gamma_{ij}
\]

бо за друге елементы \(x=\sum x_j\tau_j\) просто јест \(x'=\sum x'_j\tau_j\). Те матрице \(C\) сут једнозначно приряđене различным линеарным трансформацијам (бо једне определьајут друге) и, како се види непосреднье, творят колцо изоморфно с \(\overline{\mathfrak D}\), теды колцово хомоморфны образ од \(\mathfrak{o}\), представјенье \(n\)-того ступньа од \(\mathfrak{o}\).

Изоморфизм меджу \(\overline{\mathfrak D}\) и представјеньем садржан јест в следујучих формулах: ако линеарној трансформацији, створјеној елементом \(c\in\mathfrak{o}\), одповедаје матрица \(C\), теды

\[
c(x_1,\ldots,x_n)=(x_1,\ldots,x_n)C
\]

и ако на једнакы начин \(d\in\mathfrak{o}\) и \(D\in\mathsf T_n\) одповедајут једно другому

\[
d(x_1,\ldots,x_n)=(x_1,\ldots,x_n)D
\]

то очивидно јест

\[
\begin{aligned}
(c+d)(x_1,\ldots,x_n)
&=(cx_1+dx_1,\ldots,cx_n+dx_n)\\
&=(x_1,\ldots,x_n)(C+D)
\end{aligned}
\]

и

\[
\begin{aligned}
cd(x_1,\ldots,x_n)
&=c(dx_1,\ldots,dx_n)=c((x_1,\ldots,x_n)D)\\
&=(cx_1,\ldots,cx_n)D=(x_1,\ldots,x_n)CD.
\end{aligned}
\]

Тако линеарны трансформације, створјене елементами из \(\mathfrak{o}\), од \(\mathfrak M\) в себе при јих реализацији посредством установјене базы \(x_\nu\) модула \(\mathfrak M\) давајут представјенье \(n\)-того ступеньа од \(\mathfrak{o}\) в \(\mathsf T\).

Ако \((y_1,\ldots,y_n)\) јест друга база модула \(\mathfrak M\), она настава из \(x_1,\ldots,x_n\) множеньем регуларној матрицоју

\[(y_1,\ldots,y_n)=(x_1,\ldots,x_n)P\]

Ако линеарнаја трансформација од \(\mathfrak M\) взходно к базе \(x_\nu\) јест реализована матрицоју \(C\), тада взходно к базе \(y_\nu\) она јест реализована матрицоју \(P^{-1}CP\), бо ако

\[c(x_1,\ldots,x_n)=(x_1,\ldots,x_n)C\]
то важи

\[
\begin{aligned}
c(y_1,\ldots,y_n)
&=c(x_1,\ldots,x_n)P=((x_1,\ldots,x_n)C)P\\
&=(x_1,\ldots,x_n)CP=((y_1,\ldots,y_n)P^{-1})CP\\
&=(y_1,\ldots,y_n)P^{-1}CP.
\end{aligned}
\]

Из тога непосреднье следује:

\emph{\(\mathfrak M\) даваје вса представјеньа једној класы представјень, ако за реализацију линеарных трансформациј од \(\mathfrak M\), створјеных елементами из \(\mathfrak{o}\), употребимо все базы модула \(\mathfrak M\).}

\begin{enumerate}
\def\labelenumi{\arabic{enumi}.}
\setcounter{enumi}{1}
\tightlist
\item
  \emph{Всака класа представјень јест створјена модулом представјеньа на начин указаны под 1.} (Знову се ограничујемо на преме представјеньа.)
\end{enumerate}

\emph{Доказ.} Изберемо установјено представјенье из класы; нехај оно буде задано хомоморфным приряđеньем \(c\mapsto C\). Ако \(n\) јест ступень представјеньа, то с \(n\) неизвестными \(x_1,\ldots,x_n\) образујемо правы \(\mathsf T\)-модул

\[\mathfrak{M} = x_1 \mathsf{T} + \dots + x_n \mathsf{T}\]

(с обычными правилами рачунаньа \(\sum x_i\tau_i+\sum x_i\overline{\tau}_i=\sum x_i(\tau_i+\overline{\tau}_i)\), \(\sum x_i\tau_i=\sum x_i\overline{\tau}_i\) тогда и само тогда, ако \(\tau_i=\overline{\tau}_i\) за все \(i\), \((\sum x_i\tau_i)\varrho=\sum x_i\tau_i\varrho\)), и чинимо га левым \(\mathfrak{o}\)-модулом установјеньем \(c\sum x_i\tau_i=\sum x_i'\tau_i\), ако \((x_1',\ldots,x_n')=(x_1,\ldots,x_n)C\), причем под \(C\) разумејемо матрицу представјеньа, приряđену елементу из \(\mathfrak{o}\), \(c\).

Тым установјеньем \(\mathfrak M\) јест учиньено двојным \(\mathfrak{o}\)-\(\mathsf T\)-модулом, бо прво из релациј хомоморфизма следује

\[
c+d\mapsto C+D \quad \text{i} \quad cd\mapsto CD
\]

\[
(c+d)x_i=cx_i+dx_i
\]

\[
cd\cdot x_i=c\cdot dx_i,\quad\text{nadalje jest}\quad
c\left(\sum x_i\tau_i+\sum x_i\overline{\tau}_i\right)
=c\sum x_i\tau_i+c\sum x_i\overline{\tau}_i,
\]

то сут својства левого модула; друго, очивидно јест

\[c(\sum x_i \, \tau_i \, \varrho) = (c \, \sum x_i \, \tau_i) \cdot \varrho \,,\]

то јест смјешаны закон асоциативности.

\(\mathfrak{M}\) јест теды модул представјеньа, а употребјено представјенье јест ним створјено на начин указаны под 1., ако за реализацију ньеговых линеарных трансформациј употребимо базу \(x_1, \ldots, x_n\). Друге базы давајут остале представјеньа те класы.

За реципрочна представјеньа и модуле представјень все иде једнако; достава се, же модул

\[\mathfrak{M}^* = \mathsf{T}^* x_1^* + \dots + \mathsf{T}^* x_n^*\]

стварја реципрочно представјенье од \(\mathfrak{o}\) в \(\mathsf T^*\) тако, же елементу из \(\mathfrak{o}\), \(c\), приряđује матрицу \(C\) в једначеньу

\[
c\begin{pmatrix}x_1\\ \vdots\\ x_n\end{pmatrix}
=\begin{pmatrix}cx_1\\ \vdots\\ cx_n\end{pmatrix}
=C\begin{pmatrix}x_1\\ \vdots\\ x_n\end{pmatrix}
\]

Различне базы знову давајут различна представјеньа једној класы.

\section*{Глава II. Галоисова теорија комутативных поль}\label{kapitel-ii.-galoissche-theorie-kommutativer-koerper}

Досега развијена теорија представјень јест достаточна за основанье Галоисовој теорији, ктора имаје предност над обычној теоријеју, же не потребује специалне системы творечих елементов ни теоремы о полиномах неизвестных.\srcfn{1}{За следујуче сравньај M. Z. §§ 15 --- 20. M. Z. § 21 садржује прве теоремы следујучего оддела о Галоисовој теорији. За дефиницију хиперкомплексных системов види M. Z. § 6.}
% END BOOK_S04

% BEGIN BOOK_S05
\subsection*{§ 5. Разширјенье коефициентов хиперкомплексных системов}\label{koeffizientenerweiterung-hyperkomplexer-systeme}

Формулујемо следујучи обчи теорем, чиј доказ оставјамо неизведеным, бо он јест скоро саморазумны. Нека \(\mathfrak{o}=c_1\mathsf P+\cdots+c_n\mathsf P\) буде хиперкомплексны систем над комутативным польем \(\mathsf P\). Нека \(\Omega\) буде разширјујуче колцо од \(\mathsf P\), чији центар садржује \(\mathsf P\). Дефинујемо \emph{разширјены систем} \(\mathfrak{o}_{\Omega}\) како множство всих формалных сум \(\sum c_i\omega_i\), \(\omega_i\in\Omega\), кторо најпрво установјеньами

\[
\begin{array}{l}
\sum c_i\omega_i=\sum c_i\omega_i' \quad\text{togda i samo togda, ako } \omega_i=\omega_i' \text{ za vse } i;\\
\sum c_i\omega_i+\sum c_i\omega_i'=\sum c_i(\omega_i+\omega_i');\\
\omega\sum c_i\omega_i=\sum c_i\omega\omega_i;\quad
\sum c_i\omega_i\cdot\omega=\sum c_i\omega_i\omega;\\
\sum c_i\omega_i=\sum \omega_i c_i
\end{array}
\]

поставаје \(\Omega\)-модулом, комутативно связаным с \(\Omega\), а дальным установјеньем

\[
\sum c_i\omega_i\cdot \sum c_i\omega_i
=\sum_i\sum_j c_ic_j\omega_i\omega_j
\]

--- причем \(c_ic_j\) јест уже дефинованы како елемент из \(\mathfrak{o}\) --- поставаје разширјујучим колцом од \(\mathfrak{o}\). \(\mathfrak{o}_{\Omega}\) имаје ранг \(n\) над \(\Omega\) и можно га в форме

\[
\mathfrak{o}_{\Omega}=c_1\Omega+\cdots+c_n\Omega
\]

записати.

\emph{\(\mathfrak{o}_{\Omega}\) јест дефинованы независно од базы \(c_1,\ldots,c_n\). Ако \(\Omega\) јест комутативно полье, то \(\mathfrak{o}_{\Omega}\) јест хиперкомплексны систем над \(\Omega\).}
% END BOOK_S05

% BEGIN BOOK_S06
\subsection*{§ 6. Иредуцибилне представјеньа комутативных системов}\label{die-irreduziblen-darstellungen-kommutativer-systeme}

Нека \(\mathfrak Z=z_1\mathsf P+\cdots+z_n\mathsf P\) буде комутативны систем над \(\mathsf P\). Хочемо установити вса иредуцибилна представјеньа од \(\mathfrak Z\) в алгебраично затворјеном разширјеном польу \(\Omega\) од \(\mathsf P\). Понеже всако представјенье од \(\mathfrak Z\) в \(\Omega\) води к представјеньу од \(\mathfrak Z_\Omega=z_1\Omega+\cdots+z_n\Omega\) --- бо представјеньа базовых елементов, уже садржаных в \(\mathsf P\), сут достаточне за установјенье целокупного представјеньа --- и обратно всако представјенье од \(\mathfrak Z_\Omega\) садржује представјенье од \(\mathfrak Z\), можно јест пытати такоже за представјеньа од \(\mathfrak Z_\Omega\). Она сут по M. Z. § 20 садржана в регуларном представјеньу од \(\mathfrak Z_\Omega\), т.~ј. в представјеньу од \(\mathfrak Z_\Omega\), кторо настава, ако сам \(\mathfrak Z_\Omega\) употребимо како модул представјеньа. Достаточно јест чак ограничити се на \(\mathfrak Z_\Omega/\mathfrak C\) како модул представјеньа, причем под \(\mathfrak C\) разумејемо радикал од \(\mathfrak Z_\Omega\). (види M. Z. § 19).

Але \(\mathfrak Z_\Omega/\mathfrak C\) јест полно редуцибилны систем, теды према сума поль конечного ступеньа над \(\Omega\). Понеже \(\Omega\) јест алгебраично затворјено, те суманды сут изоморфне с \(\Omega\).

\[
\mathfrak Z_\Omega/\mathfrak C=e_1\Omega+\cdots+e_t\Omega.
\]

Елементи \(e_\nu\) сут компоненты јединице. Теды важи:

\emph{\(\mathfrak Z_\Omega\) имаје точно \(t\) различных клас иредуцибилных представјень в \(\Omega\), все првого ступеньа. При том \(t\) јест ранг од \(\mathfrak Z_\Omega/\mathfrak C\), теды највыше равны \(n\).}

Всака из тых клас представјень садржује само једно представјенье, бо трансформација представјеньа првого ступеньа в комутативном польу не изменьаје представјенье: \(\pi^{-1}\cdot\alpha\cdot\pi=\alpha\).

Зато постојут точно \(t\) различне хомоморфизмы \(\Theta_1,\ldots,\Theta_t\) од \(\mathfrak Z_\Omega\) на подколца од \(\Omega\).
% END BOOK_S06

% BEGIN BOOK_S07
\subsection*{§ 7. Изоморфизмы польа}\label{die-isomorphismen-eines-koerpers}

Ныне нека \(\mathfrak Z\) буде полье. Всакы хомоморфизм \(\Theta_i\) од \(\mathfrak Z_\Omega\) на подколцо од \(\Omega\) мора изоморфно образовати \(\mathfrak Z\) на подполье \(Z_i\) од \(\Omega\). Бо или јест \(Z_i=0\) --- а то тут не јест случај, бо јединице од \(\mathfrak Z\) одповедаје јединица од \(\Omega\) --- или \(Z_i\) јест изоморфны образ од \(\mathfrak Z\). Зато:

\emph{Ако \(\mathfrak Z\) јест конечно разширјенье \(n\)-того ступеньа од \(\mathsf P\), то в \(\Omega\) постојут толико различне польа, изоморфне с \(\mathfrak Z\), \(Z_i\), коликы јест ранг \(t\) од \(\mathfrak Z_\Omega/\mathfrak C\). Число поль \(Z_i\) јест теды највыше равно \(n\).}

Ныне называјемо \(\mathfrak Z\) польем \emph{првого рода} над \(\mathsf P\), ако число поль \(Z_i\) јест равно \(n\), а польем \emph{другого рода}, ако то число јест менье од \(n\).

\emph{\(\mathfrak Z\) јест тогда и само тогда првого рода над \(\mathsf P\), ако \(\mathfrak Z_\Omega\) јест полно редуцибилно.}

Тут видимо предност тога основаньа Галоисовој теорији над обычным. В обычној Галоисовој теорији разматрајут се изоморфизмы једного из поль \(Z_i\) на друге. Тут се избегаје та несиметрија. За доказ, же полье не имаје выше коньугованых поль, неж колико указује јего ступень, обычно се употребје примитивны елемент. Тут тој факт просто следује из тога, же вса иредуцибилна представјеньа од \(\mathfrak Z_\Omega\) сут садржана в регуларном представјеньу.
% END BOOK_S07

% BEGIN BOOK_S08
\subsection*{§ 8. Разпадно полье}\label{der-zerfaellungskoerper}

Польа \(Z_i\) сут иредуцибилна представјеньа од \(\mathfrak Z\) в \(\Omega\). Ако теды полье \(\Gamma\) меджу \(\mathsf P\) и \(\Omega\), \(\mathsf P\subseteq\Gamma\subseteq\Omega\), садржује вса \(Z_i\), то уже \(\mathfrak Z_\Gamma/\mathfrak C\) мора се разпадати на \(t\) поль првого ступеньа \(\Gamma\). Обратно, да бы \(\mathfrak Z_\Gamma/\mathfrak C\) разпадало се на \(t\) поль, кторе сут тада иредуцибилными представјеньами првого ступеньа од \(\mathfrak Z_\Gamma\), польа \(Z_i\) морајут бити садржане в \(\Gamma\). Ако разпадным польем од \(\mathfrak Z\) называмо полье \(\Gamma\) с својством, же \(\mathfrak Z_\Gamma/\mathfrak C\) поставаје премоју сумоју \(t\) поль, то имајемо:

\emph{Галоисово полье поль, изоморфных с \(\mathfrak Z\), \(Z_i\), т.~ј. јих композитум, јест минимално разпадно полье од \(\mathfrak Z\).}
% END BOOK_S08

% BEGIN BOOK_S09
\subsection*{§ 9. Изоморфизмы од \(Z_1\) на \(Z_i\)}\label{die-isomorphismen-von-z1-auf-die-zi}

\(\Theta_i\), ако разматрамо само \(\mathfrak Z\), јест изоморфизм од \(\mathfrak Z\) на \(Z_i\). Теды \(S_i=\Theta_i\Theta_1^{-1}\) јест изоморфизм од \(Z_1\) на \(Z_i\). Понеже \(\Theta_i\) сут различне, такоже \(S_i\) сут различне. \(S_i\) сут изоморфизмы једного польа \(Z_1\) на јего коньугована польа, кторе се обычно разматрајут в Галоисовој теорији. Другых не имаје, бо не имаје другых \(\Theta_i\).
% END BOOK_S09

% BEGIN BOOK_S10
\subsection*{§ 10. Галоисова група}\label{die-galoissche-gruppe}

Ако \(\mathfrak Z\) јест Галоисово полье, теды ако польа \(Z_i\) совпадајут како множства елементов, то \(S_i\) творят групу всих изоморфизмов, кторе оставјајут \(\mathsf P\) неподвижным, од \(Z\) на себе. Бо всакы продукт \(S_iS_j\) јест автоморфизм од \(Z\), а понеже не имаје других осим \(S_i\), \(S_iS_j\) јест равны некому \(S_k\). Теды \(S_i\) творят групу всих автоморфизмов, кторе оставјајут всакы елемент из \(\mathsf P\) неподвижным, од \(Z\); то јест \emph{Галоисова група од \(Z\) взходно к \(\mathsf P\)}.
% END BOOK_S10

% BEGIN BOOK_S11
\subsection*{§ 11. Главны теорем Галоисовој теорије}\label{der-hauptsatz-der-galoisschen-theorie}

Ныне нека \(\mathfrak Z\) буде Галоисово разширјенье првого рода од \(\mathsf P\). Доказујемо главны теорем Галоисовој теорије:

\emph{Польа \(\Gamma\) меджу \(\mathsf P\) и \(\mathfrak Z\) можно јест једнозначно взајемно прирядити подгрупам \(\mathfrak H\) Галоисовој групы \(\mathfrak G\) од \(Z\) взходно к \(\mathsf P\) тако, же польу \(\Gamma\) приряđена група \(\mathfrak H\) состоји се из всих автоморфизмов од \(Z\), кторе оставјајут всакы елемент из \(\Gamma\) неподвижным, а \(\Gamma\) јест множство всих елементов из \(Z\), кторе оставајут неподвижне при всих автоморфизмах, садржаных в \(\mathfrak H\).}

Кратко то формулујемо тако:

\begin{enumerate}
\item \emph{Ако \(\mathfrak H\) јест инвариантна група од \(\Gamma\), то \(\Gamma\) јест инвариантно полье од \(\mathfrak H\).}
\item \emph{Ако \(\Gamma\) јест инвариантно полье од \(\mathfrak H\), то \(\mathfrak H\) јест инвариантна група од \(\Gamma\).}
\end{enumerate}

\paragraph{1. можно свести на помочну лему:}\label{laesst-sich-zurueckfuehren-auf-den-hilfssatz}

\emph{Ако \(\mathsf P \subseteq \Sigma \subseteq \mathsf T \subseteq \mathsf Z\) и \(\mathsf T\) имаје ступень \(t\) над \(\Sigma\), то всакы изоморфизм од \(\Sigma\) на неко коньуговано полье (нужно садржано в \(\mathsf Z\)) можно јест продолжити до \(t\) различных изоморфизмов од \(\mathsf T\) на коньугована польа.}

Ако бо предположимо, же \(\mathfrak H\) јест инвариантна група од \(\Sigma\), а \(\mathsf T\) јест инвариантно полье од \(\mathfrak H\), можно јест закльучити следујуче: Елементи из \(\mathfrak H\) сут продолженьа на цело \(\mathsf Z\) идентичного изоморфизма од \(\Sigma\). Ако те продолженьа применимо само на \(\mathsf T\), из помочној лемы следује, же она разпадајут се на \(t\) клас тако, же изоморфизмы исте класы сут идентичне на \(\mathsf T\), а изоморфизмы различных клас сут различне на \(\mathsf T\). Але понеже \(\mathsf T\) јест инвариантно полье од \(\mathfrak H\), мора бити \(t=1\), \(\mathsf T=\Sigma\), что было доказати.

За доказ помочној лемы врнемо се к польам \(\mathfrak S\), \(\mathfrak T\) меджу \(\mathsf P\) и \(\mathsf Z\), кторым сут \(\Sigma\) и \(\mathsf T\) приряđене посредством \(\Theta_1\). Тада треба доказати:

\emph{Всакы изоморфизм \(H_i\) од \(\mathfrak S\) на неко \(\Sigma_i\) можно јест на точно \(t\) различных начинов продолжити до изоморфизмов од \(\mathfrak T\) на подпольа \(\mathsf T_i\) од \(\mathsf Z\).}

Заради предположеньа, же \(\mathfrak Z\) јест првого рода над \(\mathsf P\), \(\mathfrak Z_\Omega/\mathfrak C\) јест идентично с \(\mathfrak Z_\Omega\), теды \(\mathfrak S_\Omega\) разпада се на толико поль, колико указује ньегов ступень \(s\)
\[
\mathfrak S_\Omega = E_1\Omega+\cdots+E_s\Omega .
\]
Модулы представјень \(E_\nu\Omega\) стварјајут \(s\) изоморфизмов од \(\mathfrak S\) на \(s\) подполь \(\Sigma_1,\ldots,\Sigma_s\). За најданье изоморфизмов од \(\mathfrak T\) потребујемо \(\mathfrak T_\Omega\). Понеже \(E_\nu\Omega\) сут идеалы од \(\mathfrak S_\Omega\), важи \(E_\nu\Omega=\mathfrak S_\Omega E_\nu\). Надалье очивидно јест
\[
\mathfrak T_\Omega=\mathfrak T\cdot\Omega=\mathfrak T\cdot\mathfrak S_\Omega .
\]
Зато јест
\[
\mathfrak T_\Omega
=E_1\mathfrak S_\Omega\cdot\mathfrak T+\cdots+E_s\mathfrak S_\Omega\cdot\mathfrak T
=E_1\mathfrak T_\Omega+\cdots+E_s\mathfrak T_\Omega .
\]
То јест разклад од \(\mathfrak T_\Omega\) на идеалы. \(E_\nu\cdot\mathfrak T_\Omega\) разпада се на \(t\) простых идеалов. Бо важе релације ступньев
\[
[\mathfrak T_\Omega\cdot E_\nu:\mathfrak S_\Omega E_\nu]
\leqq [\mathfrak T_\Omega:\mathfrak S_\Omega]=[\mathfrak T:\mathfrak S]=t,
\]
теды заради \(\mathfrak S_\Omega E_\nu=\Omega E_\nu\simeq\Omega\):
\[
[\mathfrak T_\Omega\cdot E_\nu:\Omega]
=[\mathfrak T_\Omega E_\nu:\mathfrak S_\Omega E_\nu]\leqq t,
\]
и понеже сума всих ступньев \([\mathfrak T_\Omega E_\nu:\Omega]\) мора бити равна \(s\cdot t\):
\[
[\mathfrak T_\Omega\cdot E_\nu:\Omega]=t,\qquad
E_\nu\mathfrak T_\Omega=e_\nu^{(1)}\cdot\Omega+\cdots+e_\nu^{(t)}\Omega .
\]
Всакы \(e_\nu^{(i)}\Omega\) даваје иредуцибилно представјенье од \(\mathfrak T_\Omega\). \emph{Покажимо јешче}, же вса та представјеньа, ако јих применимо само на \(\mathfrak S_\Omega\), совпадајут с представјеньем, створјеным посредством \(E_\nu\cdot\Omega\); тада јесмо готови. То але јест јасно, бо из
\[
s\cdot E_\nu=E_\nu\sigma_\nu\quad
(s\in\mathfrak S_\Omega\text{ imaje }\sigma_\nu\text{ kako predstavjenje})
\]
следује
\[
\sum s\cdot e_\nu^{(i)}=\sum e_\nu^{(i)}\sigma_\nu,
\]
теды заради једнозначности записа сумы
\[
s\cdot e_\nu^{(i)}=e_\nu^{(i)}\sigma_\nu,
\]
нове представјеньа такоже приряджајут елементу \(s\) елемент \(\sigma_\nu\).

\paragraph{2. Нека \(\mathsf T\) буде инвариантно полье од \(\mathfrak H\).}
Хочемо познати \(\mathfrak H\) како инвариантну групу од \(\mathsf T\); за то \emph{јест достаточно} показати елемент из \(\mathsf T\), кторы не оставаје неподвижным при неком автоморфизму, не належечем к \(\mathfrak H\).

Елементи \(S_i\) Галоисовој групы \(\mathfrak G\) од \(\mathsf Z\) сут автоморфизмы, кторе оставјајут \(\mathsf P\) неподвижным, од \(\mathsf Z\). Чинимо јих автоморфизмами од \(\mathfrak Z\) установјеньем
\[
S_i z=z,\quad \text{ako } z \text{ jest v } \mathfrak Z.
\]
Применьенье од \(S_i\) на
\[
\mathfrak Z_{\mathsf Z}=e_1\mathsf Z+\cdots+e_n\mathsf Z
\]
мора пременьивати преме суманды \(e_\nu\mathsf Z\), бо оне сут једнозначно определьене; в особливости мора бити
\[
S_i e_\nu=e_{\nu'},\qquad\text{i}\qquad S_i e_\nu\ne S_k e_\nu
\quad\text{za } i\ne k
\]
Ако \(S_1,\ldots,S_h\) сут елементи из \(\mathfrak H\), изберемо неко \(e_\nu\), рецимо \(e'\), и разматрамо все \(e_\nu\), кторе из ньега стварјајут \(S_i\), \(i\leqq h\):
\[
e_1=S_1e',\ldots,e_h=S_he'.
\]
Понеже, ако \(i\leqq h, j\leqq h\),
\[
S_i e_j=S_iS_j e'=S_ke'=e_k \quad \text{pri}\quad k\leqq h
\]
всакы \(S_i\) из \(\mathfrak H\) пременьује \(e_1,\ldots,e_h\) меджу собоју,
\[
E_1=e_1+\cdots+e_h
\]
јест теды при всих \(S_i\in\mathfrak H\) инвариантны елемент из \(\mathfrak Z\). Ако представимо \(E_1\) посредством базы од \(\mathfrak Z\), јего коефициенты морајут бити при \(\mathfrak H\) инвариантне елементи из \(\mathsf Z\); зато належет к \(\mathsf T\).

\(E_1\) не јест инвариантны при никаком елементу изван \(\mathfrak H\), \(S_i\). Бо за тако \(S_i\)
\[
S_i e_1=e_i+\cdots
\]
садржује меджу компонентами, кторе не наступајут меджу \(e_1,\ldots,e_h\), компоненту \(e_i\), але запис сумы од \(E_1\) јест једнозначны.
% END BOOK_S11

% BEGIN BOOK_S12
\subsection*{§ 12. Формално значенье компонентов \(e_i\)}\label{die-formale-bedeutung-der-komponenten-ei}

За основу изберемо установјену базу \(z_1,\ldots,z_n\) од \(\mathfrak Z\).
\[
\mathfrak Z=z_1\mathsf P+\cdots+z_n\mathsf P
\]
(Ныне о \(\mathfrak Z\) предпокладамо само, же оно јест првого рода). Изоморфизм \(\Theta_i\) образује систем \(z_1,\ldots,z_n\) на базу \(\zeta_1^{(i)},\ldots,\zeta_n^{(i)}\) од \(\mathsf Z_i\), и то релацијами
\[
z_\nu e_i=e_i\zeta_\nu^{(i)} .
\]

Обратно, компонента јединице \(e_j\) мора бити изразива посредством \(z_\nu\) с коефициентами из \(\mathsf Z\)
\[
e_j=\sum_\nu \varrho_\nu^{(j)}z_\nu .
\]
Теды имајемо
\[
e_j=e_j^2=\sum \varrho_\nu^{(j)}z_\nu e_j
=\sum \varrho_\nu^{(j)}\cdot\zeta_\nu^{(j)}\cdot e_j
\]
и
\[
0=e_je_k=\sum \varrho_\nu^{(j)}\cdot z_\nu\cdot e_k
=\sum \varrho_\nu^{(j)}\cdot\zeta_\nu^{(k)}\cdot e_k,
\]
теды:
\[
\sum_\nu \varrho_\nu^{(j)}\cdot\zeta_\nu^{(j)}=1;
\qquad
\sum_\nu \varrho_\nu^{(j)}\cdot\zeta_\nu^{(k)}=0,\quad
\text{ako } j\ne k.
\]
То значи: Матрице
\[
\begin{pmatrix}
\zeta_1^{(1)} & \cdots & \zeta_n^{(1)}\\
\cdots & \cdots & \cdots\\
\cdots & \cdots & \cdots\\
\cdots & \cdots & \cdots\\
\zeta_1^{(n)} & \cdots & \zeta_n^{(n)}
\end{pmatrix}
\quad\text{i}\quad
\begin{pmatrix}
\varrho_1^{(1)} & \cdots & \varrho_1^{(n)}\\
\cdots & \cdots & \cdots\\
\cdots & \cdots & \cdots\\
\cdots & \cdots & \cdots\\
\varrho_n^{(1)} & \cdots & \varrho_n^{(n)}
\end{pmatrix}
\]
сут взајемно обратне. То можно такоже изразити следујуче: \(\varrho_1^{(i)},\ldots,\varrho_n^{(i)}\) јест база комплементарна к \(\zeta_1^{(i)},\ldots,\zeta_n^{(i)}\). Она имаје ролу в теорији диференты числового польа.
% END BOOK_S12

% BEGIN BOOK_S13
\subsection*{§ 13. Обчи теорем о продолженьу}\label{der-allgemeine-fortsetzungssatz}

Теорем о продолженьу из \(H_i\), доказан на стр.~10, можно јест обчити.

Нека \(\mathfrak Z=z_1\mathsf P+\cdots+z_n\mathsf P\) буде комутативны, хиперкомплексны, полно редуцибилны систем. Нека \(\mathfrak S\) и \(\mathfrak T\) буду -- в тој ситуацији нужно такоже полно редуцибилне -- подсистемы од \(\mathfrak Z\):
\[
\mathsf P\subseteq\mathfrak S\subseteq\mathfrak T\subseteq\mathfrak Z .
\]
По M. Z. § 13 всакы из трох системов имаје јединицу. Але те три јединице никако не морајут бити једнаке. (Сигурно не сут једнаке, ако \(\mathfrak S\) и \(\mathfrak T\) сут идеалы од \(\mathfrak Z\), различне од \(\mathfrak Z\).) Нека \(E\) буде јединица од \(\mathfrak S\). Систем \(\mathfrak T\cdot E\) јест модул конечного ранга над \(\mathfrak S\):
\[
E\mathfrak T=a_1\mathfrak S+\cdots+a_t\mathfrak S
\]
(\(\mathfrak T\) обче не јест!).

\emph{Доказ.} \(\mathfrak S\) јест према сума поль: \(\mathfrak S=\mathfrak K_1+\cdots+\mathfrak K_r\). Тому одповедаје разклад од \(\mathfrak T E\) на прему суму. Всакы суманд од \(\mathfrak T E\) мора садржевати једно из поль \(\mathfrak K_i\), иначе бы \(E\) ануловало го, але \(E\) јест јединица од \(\mathfrak T E\). Зато \(\mathfrak T E\) јест према сума разширјујучих колц од \(\mathfrak K_i\), т.~ј. према сума формы
\[
\mathfrak T E=a_{11}\mathfrak K_1+\cdots+a_{1p_1}\mathfrak K_1+\cdots+a_{rp_r}\mathfrak K_r,
\]
или \(E\mathfrak T=a_{11}\mathfrak S+\cdots+a_{rp_r}\mathfrak S\), что было доказати.

\begin{center}
Обчи теорем о продолженьу.
\end{center}

\emph{Всако иредуцибилно представјенье од \(\mathfrak S\) в \(\Omega\), задано хомоморфизмом \(H_i\), можно јест на точно \(t\) различных начинов продолжити до представјень од \(\mathfrak T\).}

Доказ јест аналогны доказу на стр.~10--11. Имајемо
\[
\mathfrak S_\Omega=E_1\Omega+\cdots+E_s\Omega,\qquad
E_\nu\Omega=\mathfrak S_\Omega E_\nu,
\]
\[
\mathfrak T_\Omega\cdot E=\mathfrak T\cdot E\Omega
=\mathfrak T\mathfrak S\Omega=\mathfrak T\mathfrak S_\Omega,
\]
\[
\mathfrak T_\Omega\cdot E
=\mathfrak T(E_1\mathfrak S_\Omega+\cdots+E_s\mathfrak S_\Omega)
=E_1\mathfrak T_\Omega+\cdots+E_s\mathfrak T_\Omega,
\]
\[
[\mathfrak T_\Omega E_\nu:\Omega]
=[\mathfrak T_\Omega E_\nu:\Omega E_\nu]
=[\mathfrak T_\Omega E_\nu:E_\nu\mathfrak S_\Omega]
\leqq[\mathfrak T_\Omega\cdot E:\mathfrak S_\Omega]
=[\mathfrak T\cdot E:\mathfrak S]=t,
\]
\[
\sum[\mathfrak T_\Omega\cdot E_\nu:\Omega]=st
\]
теды
\[
[\mathfrak T_\Omega E_\nu:\Omega]=t.
\]
\(\mathfrak T_\Omega E_\nu\) разпада се теды на \(t\) идеалов, из кторых всакы како суманд од \(\mathfrak T_\Omega E\), теды такоже од \(\mathfrak T_\Omega\), даваје представјенье од \(\mathfrak T_\Omega\), кторо, како на стр.~10--11, оказује се продолженьем представјеньа заданого посредством \(H_\nu\). Остала иредуцибилна представјеньа од \(\mathfrak T_\Omega\) -- створјена сумандами, не садржаными в \(\mathfrak T_\Omega\cdot E_\nu\), од \(\mathfrak T_\Omega\) -- не могут бити продолженьа представјень, створјеных посредством \(H_\nu\), од \(\mathfrak S_\Omega\), бо те идеалы сут ануловане посредством \(\mathfrak T_\Omega\cdot E_\nu\), а тым выше посредством \(\mathfrak S_\Omega\).

\section*{Глава III. Абелове групы}\label{kapitel-iii.-abelsche-gruppen}
% END BOOK_S13

% BEGIN BOOK_S14
\subsection*{§ 14. Групово колцо}\label{der-gruppenring}

Нека \(\mathfrak G\) буде конечна група с елементами \(a_1,\ldots,a_h\), а \(\mathsf P\) комутативно полье. С груповым множеньем како множеньем образујемо хиперкомплексны систем
\[
\mathfrak o[\mathfrak G]=a_1\mathsf P+\cdots+a_h\mathsf P
\]
\emph{групово колцо од \(\mathfrak G\) в \(\mathsf P\)}. (M. Z. § 6, Спеисер, Тхеорие д. Группен енд. Орд. § 48).

Пытајемо за представјеньа од \(\mathfrak o[\mathfrak G]\) в \(\mathsf P\) или в алгебраично затворјеном польу \(\Omega\) над \(\mathsf P\). Како уже было споменуто, иредуцибилна представјеньа сут створјена простыми левыми идеалами од \(\mathfrak o[\mathfrak G]/\mathfrak C\). Зато нас -- како в Галоисовој теорији -- интересује пытанье, когда \(\mathfrak o[\mathfrak G]\) имаје радикал. О том важи обчи теорем:

\begin{center}
\emph{\(\mathfrak o[\mathfrak G]\) јест тогда и само тогда полно редуцибилно, ако ред \(h\) од \(\mathfrak G\) не јест деливы карактеристикоју \(p\) од \(\mathsf P\).}
\end{center}

Најпрво доказујемо:
\[
\text{Iz } h\equiv0(p) \text{ slěduje } \mathfrak C\ne 0.
\]
\[
\mathfrak a=\mathsf P(a_1+\cdots+a_h)
\]
јест двустранны идеал од \(\mathfrak o\). Бо јест
\[
\mathsf P\cdot(a_1+\cdots+a_h)a
=\mathsf P\cdot(a_1a+\cdots+a_ha)
=\mathsf P(a_1+\cdots+a_h)=\mathfrak a
\]
и
\[
a\cdot\mathsf P(a_1+\cdots+a_h)
=\mathsf P(aa_1+\cdots+aa_h)
=\mathsf P(a_1+\cdots+a_h)=\mathfrak a
\]

заради групового својства елементов \(a_i\). \(\mathfrak a\) очивидно не јест нула, але јест нилпотентны:
\[
\left(\sum_i a_i\right)^2=\sum_i\sum_j a_i a_j
=h\sum_j a_j=0.
\]

Обратно тврдженье:
\[
\text{Iz } h\not\equiv0(p) \text{ slěduje } \mathfrak C=0,
\]
ныне доказујемо само за Абелове групы. (В M. Z. § 26 та чест теорема јест доказана обче.)
% END BOOK_S14

% BEGIN BOOK_S15
\subsection*{§ 15. Групове колца Абеловых груп}\label{die-gruppenringe-abelscher-gruppen}

Под \(\mathfrak A=\{a_1,\ldots,a_h\}\) разумејемо Абелову групу. Тада \(\mathfrak o\) јест комутативно, зато можно јест применити обче теоремы о комутативных хиперкомплексных системах.

\emph{Ако \(h\) не јест деливы карактеристикоју \(p\) од \(\mathsf P\), то \(\mathfrak o[\mathfrak A]\) имаје точно \(h\) различных иредуцибилных представјень в \(\Omega\). (Она сут првого ступеньа, бо \(\mathfrak o\) јест комутативно.)}

Из тога следује, же \(\mathfrak o[\mathfrak A]/\mathfrak C\) имаје тојже ранг како \(\mathfrak o\), теды же \(\mathfrak C=0\); то јест втора чест предходного, оставјеного без доказа теорема за Абелове групы.

\emph{Доказ.} \(\mathfrak A\) јест премы продукт цикличных груп \(\mathfrak Z_i\) редов \(h_i\)
\[
\mathfrak A=\mathfrak Z_1\times\cdots\times\mathfrak Z_t,\qquad
h=h_1h_2\cdots h_t .
\]
Ако \(z_i\) јест творечи елемент од \(\mathfrak Z_i\), то јему одповедајучи елемент иредуцибилного представјеньа (првого ступеньа) мора бити \(h_i\)-ты корень из јединице \(\varepsilon_\nu^{(h_i)}\). С другој страны, ако всакому \(z_i\) прирядимо \(h_i\)-ты корень из јединице \(\varepsilon_\nu^{(h_i)}\), достава се иредуцибилно представјенье од \(\mathfrak o\) в \(\Omega\). Понеже ныне \(h\not\equiv0(p)\), а тым выше \(h_i\not\equiv0(p)\), постојут точно \(h_1h_2\cdots h_t=h\) различне представјеньа од \(\mathfrak o\), что было доказати.

(Такоже јест легко видети: ако \(h=0(p)\), теды \(h_i=0(p)\) за најменеј једно \(h_i\), постојут највыше \(h_i/p\) различне \(\varepsilon_\nu^{(h_i)}\), теды менье неж \(h\) представјень, зато \(\mathfrak o_\Omega\) мора садржевати ненуловы радикал. Але из тога безпосреднье не следује ничто о радикалу од \(\mathfrak o\).)

Те \(h\) различне хомоморфизмы од \(\mathfrak o_\Omega\) на подколца од \(\Omega\) означимо с \(\Theta_1,\ldots,\Theta_h\); всакому \(\Theta_i\) одповедаје суманд \(e_i\Omega\) од \(\mathfrak o_\Omega\) како модул представјеньа. Некды применьујемо \(\Theta_i\) само на \(\mathfrak A\) и тада говоримо о представјеньу од \(\mathfrak A\). Те хомоморфизмы од \(\mathfrak A\) на групы кореньев из јединице сут \emph{\(h\) карактеров од \(\mathfrak A\)}. Специалны карактер јест \emph{главны карактер}, кторы всакому елементу из \(\mathfrak A\) приряджаје \(1\). Нека приналежечи \(\Theta_i\) буде \(\Theta_1\).

Представјенье главным карактером јест можно за всаку групу, не само за Абелове групы. Понеже то представјенье уже лежи в \(\mathsf P\), в случају, когда \(\mathfrak o\) јест полно редуцибилно, \(\mathfrak o\) мора одцепити просты суманд \(E\cdot\mathfrak o\), кторы даваје то представјенье. \(E\) можно јест безпосреднье израчунати. Мора бити \(E=\sum x_i a_i\). Прво јест
\[
aE=E\cdot 1\quad\text{ili}\quad \sum x_i a a_i=\sum x_i a_i
\]
из чего за \(a=a_j a_i^{-1}\)
\[
x_1a_ja_i^{-1}a_1+\cdots+x_i a_j+\cdots+x_na_ja_i^{-1}a_n
=x_1a_1+\cdots+x_na_n
\]
теды
\[
x_i=x_j
\]
следује.

Друго јест:

\[
E=E^2=x^2\left(\sum a_i\right)^2=x^2h\sum a_i=xhE.
\]

В случају \(h=0(p)\) следује \(E=0\), т.~ј.: \(\mathfrak o\) не садржује модул представјеньа за главны карактер, зато не могло јест бити полно редуцибилно. Тым јест знову доказана једна половица теорема на стр.~13.

Ако але \(h\not\equiv0(p)\), теды \(\mathfrak o\) јест полно редуцибилно, непосреднье следује

\[
x=\frac{1}{h},\qquad E=\frac{1}{h}(a_1+\cdots+a_h).
\]
% END BOOK_S15

% BEGIN BOOK_S16
\subsection*{§ 16. Релације карактеров}\label{die-charakterenrelationen}

Факт, же \(E=(1/h)\cdot\sum a_i\) даваје главны карактер, непосреднье води к релацијам карактеров. Мора бити

\[a \cdot e_i = e_i \cdot \Theta_i a\]

теды

\[e_i \cdot \Theta_i e_j = \begin{cases} e_i & i = j \\ 0 & i \neq j \end{cases}\]

или

\[\Theta_i e_j = \begin{cases} 1 & i = j \\ 0 & i \neq j \end{cases}\]

В особливости јест

\[\Theta_i e_1 = \begin{cases} 1 & i = 1 \\ 0 & i \neq 1 \end{cases}\]

или, понеже \(e_1 = (1/h) \sum a_i\),

\[
\sum_\nu \Theta_i a_\nu=\begin{cases} h & i=1,\\ 0 & i\ne 1, \end{cases}
\]

То сут познате релације карактеров:

\emph{Сума једного карактера по всих елементах групы јест нула, осим главного карактера, за кторы она јест х.}
% END BOOK_S16

% BEGIN BOOK_S17
\subsection*{§ 17. Галоисова теорија Абеловых груп}\label{die-galoissche-theorie-abelscher-gruppen}

За групове колца Абеловых груп можно јест установити теоремы подобне оным, кторе сут доказиване в Галоисовој теорији комутативных поль. В Галоисовој теорији \(\Theta_i\Theta_1^{-1}\) творили сут групу \(\mathfrak{G}\). Главны теорем изражаје једнозначно взајемно приряđенье, основано на установјеных инвариантных својствах, меджу подгрупами од \(\mathfrak{G}\) и подпольами од \(Z\).

Хомоморфизмы \(\Theta_i\) од \(\mathfrak{o}[\mathfrak{A}]\) можно јест сјединити в групу \(\mathsf A\), чије подгрупы стојят в подобном одношеньу к подгрупам од \(\mathfrak{A}\).

\emph{Ако продукт \(\Theta_i \Theta_k = \Theta\) дефинујемо формулом}

\[\Theta a = \Theta_i a \cdot \Theta_k a,\]

\emph{то множство всих \(\Theta_i\) јест с \(\mathfrak A\) изоморфна група \(\mathsf A\).}

\emph{Доказ.} \(\Theta = \Theta_i \cdot \Theta_k\) јест заради

\[\begin{aligned}
\Theta a \cdot \Theta b &= \Theta_i a \cdot \Theta_k a \cdot \Theta_i b \cdot \Theta_k b \\
&= \Theta_i a \cdot \Theta_i b \cdot \Theta_k a \cdot \Theta_k b \\
&= \Theta_i a b \cdot \Theta_k a b = \Theta a b
\end{aligned}\]

хомоморфизм од \(\mathfrak A\) на \(\mathsf A\), а понеже осим \(\Theta_i\) не имаје других хомоморфизмов од \(\mathfrak A\), \(\Theta\) јест равны некому \(\Theta_j\); теды \(\Theta_i\) наистино творят групу из \(h\) елементов.

Ныне под \(\varepsilon_{h_\nu}\) разумејемо примитивны \(h_\nu\)-ты корень из јединице. Тада \(\Theta_i\) образује творечи елемент \(z_\nu\) цикличного фактора \(\mathfrak Z_\nu\) од \(\mathfrak A\) на потенцију од \(\varepsilon_{h_\nu}\)

\[
\Theta_i z_\nu=\varepsilon_{h_\nu}^{\varrho_\nu^{(i)}}.
\]

Хомоморфизму \(\Theta_i\) приряджајемо елемент \(\prod_\nu z_\nu^{\varrho_\nu^{(i)}}\) од \(\mathfrak A\). То јест једнозначно взајемно приряđенье од \(\mathsf A\) на \(\mathfrak A\), бо понеже \(\Theta_i\) јест определьен посредством \(\varrho_\nu^{(i)}\), различным \(\Theta_i\) одповедајут различне \(\prod_\nu z_\nu^{\varrho_\nu^{(i)}}\). Приряđенье јест изоморфно. Бо с једној страны

\[
\Theta_i\Theta_j z_\nu=\Theta_i z_\nu\cdot \Theta_j z_\nu
=\varepsilon_{h_\nu}^{\varrho_\nu^{(i)}+\varrho_\nu^{(j)}},
\]

а с другој страны

\[
\prod_\nu z_\nu^{\varrho_\nu^{(i)}}\cdot
\prod_\nu z_\nu^{\varrho_\nu^{(j)}}=
\prod_\nu z_\nu^{\varrho_\nu^{(i)}+\varrho_\nu^{(j)}},
\quad\text{čto bylo dokazati.}
\]

Ныне дефинујемо

\emph{Инвариантну област подгрупы \(\mathsf B\) од \(\mathsf A\):}

Подгрупа \(\mathfrak B\) елементов од \(\mathfrak A\), кторе всакы \(\Theta\in\mathsf B\) образује на 1.

\emph{Инвариантну групу подгрупы \(\mathfrak B\) од \(\mathfrak A\):}

Подгрупа \(\mathsf B\) елементов од \(\mathsf A\), кторе образујут всакы \(a\in\mathfrak B\) на 1.

и формулујемо главны теорем:

\emph{Подгрупы \(\mathfrak B\) од \(\mathfrak A\) и подгрупы \(\mathsf B\) од \(\mathsf A\) можно јест једнозначно взајемно прирядити тако, же ако \(\mathfrak B\) и \(\mathsf B\) одповедајут једна другој, \(\mathsf B\) јест инвариантна група од \(\mathfrak B\), а \(\mathfrak B\) јест инвариантна област од \(\mathsf B\).}

Треба јест теды показати:

\begin{enumerate}
\item \emph{Ако \(\mathsf B\) јест инвариантна група од \(\mathfrak B\), то \(\mathfrak B\) јест инвариантна област од \(\mathsf B\).}
\item \emph{Ако \(\mathfrak B\) јест инвариантна област од \(\mathsf B\), то \(\mathsf B\) јест инвариантна група од \(\mathfrak B\).}
\end{enumerate}

\paragraph{1.}
следује како на стр.~10 из обчего теорема о продолженьу. Предпокладамо, же \(\mathsf B\) јест инвариантна група од \(\mathfrak B\), \(\mathfrak B'\) инвариантна област од \(\mathsf B\), а \(\mathfrak B'\) обсахвата \(\mathfrak B\). Елементи од \(\mathsf B\) сут продолженьа на целу \(\mathfrak A\) једного хомоморфизма од \(\mathfrak B\), кторы образује \(\mathfrak B\) на 1. Ако те продолженьа применимо само на \(\mathfrak B'\), из обчего теорема о продолженьу следује, же она разпадајут се на \(t\) клас тако, же хомоморфизмы исте класы сут идентичне, а хомоморфизмы различных клас сут различне на \(\mathfrak B'\); при том \(t\) јест ступень од \(\mathfrak o[\mathfrak B']\) над \(\mathfrak o[\mathfrak B]\), теды индекс од \(\mathfrak B\) в \(\mathfrak B'\), бо \(\mathfrak o[\mathfrak B]\) и \(\mathfrak o[\mathfrak B']\) имајут тојже елемент јединице, јменно елемент јединице од целој \(\mathfrak A\). Але понеже \(\mathfrak B'\) јест инвариантна област од \(\mathsf B\), мора бити \(t=1\), теды \(\mathfrak B'=\mathfrak B\), что было доказати.

\textbf{2.} сводимо на 1.

Елементы \(a\) од \(\mathfrak A\) сматрамо хомоморфизмами од \(\mathsf A\). Установјујемо бо, же \(a\Theta\) јест равно кореньу из јединице \(\Theta a\), кратко \(a\Theta=\Theta a\). Тым \(a\) јест наистино определьен како хомоморфизм од \(\mathsf A\), бо
\[
a\Theta\cdot a\Theta'=\Theta a\cdot\Theta' a=\Theta\Theta' a=a\Theta\Theta'.
\]

Продукт двох хомоморфизмов \(a,b\) јест обычны груповы продукт:
\[
\begin{array}{cccccc}
ab\Theta&=&a\Theta\cdot b\Theta&=&\Theta a\cdot\Theta b=\Theta ab&=ab\Theta\\
&&\text{produkt homomorfizmov}&&\text{grupovy produkt}&
\end{array}
\]

Елементи \(a,b\), различне како групове елементи, такоже сут различне како хомоморфизмы: из \(a\Theta=b\Theta\) за все \(\Theta\) следује \(\Theta a=\Theta b\) за все \(\Theta\), или \(\Theta ab^{-1}=1\) за все \(\Theta\). Теды \(\Theta\) сут продолженьа на целу \(\mathfrak A\) идентичного хомоморфизма посредством \(ab^{-1}\) створјене цикличној подгрупы \(\mathfrak Z\) од \(\mathfrak A\). По теорему о продолженьу индекс од \(\mathfrak Z\) в \(\mathfrak A\) јест равны \(h\), т.~ј. \(\mathfrak Z=e,\ a=b\).

Ныне преводимо тврдженьа, учиньена в новом разуменьу,
\[
\begin{array}{ll}
\text{a)}& \mathfrak B\ \text{jest invariantna grupa od }\mathsf B.\\
\text{b)}& \mathsf B\ \text{jest invariantna oblast od }\mathfrak B.
\end{array}
\]
в старо разуменье:

К а). \(\mathfrak B\) состоји се из всих \(a\in\mathfrak A\), за кторе \(a\Theta=1\) за все \(\Theta\in\mathsf B\), т.~ј. \(\mathfrak B\) состоји се из всих \(a\in\mathfrak A\), за кторе \(\Theta\cdot a=1\) за все \(\Theta\in\mathsf B\). Теды превод од а) јест
\[
\text{a')} \quad \mathfrak B\ \text{jest invariantna oblast od }\mathsf B.
\]

Одповедајучи превод од б) јест
\[
\text{b')} \quad \mathsf B\ \text{jest invariantna grupa od }\mathfrak B.
\]

Теды превод тврдженьа уже доказаного под 1.:

Ако \(\mathfrak B\) јест инвариантна група од \(\mathsf B\), то \(\mathsf B\) јест инвариантна област од \(\mathfrak B\),
јест равно тврдженьу, кторо треба доказати:

Ако \(\mathfrak B\) јест инвариантна област од \(\mathsf B\), то \(\mathsf B\) јест инвариантна група од \(\mathfrak B\).

Ако релације карактеров из стр.~15 запишемо за \(\mathsf A\) вместо за \(\mathfrak A\), доставајут се једначеньа
\[
\sum_{i=1}^{h} \Theta_i a_\nu =
\begin{cases}
h & \nu=1,\\
0 & \nu\ne 1.
\end{cases}
\]

Сума всих карактеров над груповым елементом јест нула, осим елемента јединице, за кторы она јест \(h\).

\section*{Глава IV. Двустранно просте колца}\label{kapitel-iv.-zweiseitig-einfache-ringe}
% END BOOK_S17

% BEGIN BOOK_S18
\subsection*{§ 18. Помочна лема}\label{ein-hilfssatz}

Под \(\mathsf K\) разумејемо тело, а под \(\mathfrak G\) групу автоморфизмов (изоморфных образовань на себе) од \(\mathsf K\). Множство елементов од \(\mathsf K\), кторе оставајут неизменьене при всаком автоморфизму належечем к \(\mathfrak G\), јест подтело \(\mathsf P\) од \(\mathsf K\), \emph{инвариантно тело од} \(\mathfrak G\).

Ако \(\tau\ne 0\) јест либовольны елемент од \(\mathsf K\), автоморфизм од \(\mathsf K\) достава се тако, же елементу од \(\mathsf K\), \(\alpha\), прирядимо елемент \(\tau^{-1}\alpha\tau\); по одповедајучем названьу из теорије груп тако автоморфизм называмо \emph{внутрным автоморфизмом}. Инвариантно тело \(\mathsf P\) групы всих внутрных автоморфизмов тела \(\mathsf K\) јест ньегов центар.

Ныне нека буду задане тело \(\mathsf K\), група автоморфизмов \(\mathfrak G\) од \(\mathsf K\) и ньено инвариантно тело \(\mathsf P\).

\[
\mathfrak M=x_1\mathsf P+\cdots+x_n\mathsf P
\]

нека буде \(\mathsf P\)-модул ранга \(n\). Дејство елемента из \(\mathfrak G\), \(\gamma\), на елементы од \(\mathfrak M\) дефинујемо формулом

\[
\gamma\cdot x_i=x_i,\qquad\text{tedy}\qquad
\gamma\sum \varrho_i x_i=\sum \varrho_i x_i.
\]

Последица тога јест, же \(\mathfrak G\) јест постала групоју автоморфизмов разширјеного модула

\[
\mathfrak M_{\mathsf K}=x_1\mathsf K+\cdots+x_n\mathsf K
\]

а инвариантна област од \(\mathfrak G\) в \(\mathfrak M_{\mathsf K}\) јест тада \(\mathfrak M\). При тых предположеньах важи теорем:

\begingroup\itshape
Ако
\[
\mathfrak C=z_1\mathsf K+\cdots+z_r\mathsf K
\]

јест подмодул од \(\mathfrak M_{\mathsf K}\), кторы всакы елемент из \(\mathfrak G\), \(\gamma\), образује в ньега самого (не нужно по елементах), то \(\mathfrak C\) јест разширјены модул, образованы с \(\mathsf K\), некого подмодула \(\mathfrak c\) од \(\mathfrak M\).

\[
\mathfrak C=\mathfrak c_{\mathsf K}.
\]
\endgroup

\emph{Доказ.} 1. Предходна припомненка: Мора бити \(\mathfrak c=\mathfrak C\cap\mathfrak M\). Бо обче важи
\[
\mathfrak c=\mathfrak c_{\mathsf K}\cap\mathfrak M.
\]

\(\mathfrak c\subseteq\mathfrak c_{\mathsf K}\cap\mathfrak M\) јест јасно. За доказ од \(\mathfrak c_{\mathsf K}\cap\mathfrak M\subseteq\mathfrak c\) изберемо базу \(y_1,\ldots,y_r\) од \(\mathfrak c\),

\[
\mathfrak c=y_1\mathsf P+\cdots+y_r\mathsf P,
\]

тада јест

\[
\mathfrak c_{\mathsf K}=y_1\mathsf K+\cdots+y_r\mathsf K.
\]

Всакы елемент из \(\mathfrak c_{\mathsf K}\cap\mathfrak M\), \(c\), допушча базовы запис како елемент од \(\mathfrak c_{\mathsf K}\)
\[
c=\sum_{i=1}^{r}y_i\chi_i
\]
и базовы запис како елемент од \(\mathfrak M\)

\[
c=\sum_{i=1}^{r}y_i\varrho_i+\sum_{j=r+1}^{n}y_j\varrho_j
\]

ако под \(y_{r+1},\ldots,y_n\) разумејемо систем елементов од \(\mathfrak M\), кторы дополньаје \(y_1,\ldots,y_r\) до базы од \(\mathfrak M\) (M. Z. § 6).

Понеже оба записа можно сматрјати базовыми записами в \(\mathfrak M_{\mathsf K}\), они сут идентичне, \(\chi_i=\varrho_i\ (i=1,\ldots,r)\), \(\varrho_j=0\ (j=r+1,\ldots,n)\), зато

\[
c=\sum_{i=1}^{r}y_i\varrho_i
\qquad\text{čto bylo dokazati.}
\]

\textbf{2.} Да бы доказати разматраны теорем, конструујемо базу \(z_1,\ldots,z_r\) од \(\mathfrak C\), чији елементи \(z_i\) належет к \(\mathfrak M\). Тада, ако положимо

\[
\mathfrak c=z_1\mathsf P+\cdots+z_r\mathsf P
\]

наистино јест

\[
\mathfrak C=\mathfrak c_{\mathsf K}.
\]

Изберемо либовольну базу \(y_1,\ldots,y_r\) од \(\mathfrak C\) и дополнимо ју до базы од \(\mathfrak M_{\mathsf K}\) одповедајучими елементами \(x_{r+1},\ldots,x_n\) једној базы \(x_1,\ldots,x_n\) од \(\mathfrak M\), ктора јест такоже база од \(\mathfrak M_{\mathsf K}\). То јест можно по M. Z. § 5.

\[
\mathfrak M_{\mathsf K}=\mathfrak C+x_{r+1}\mathsf K+\cdots+x_n\mathsf K.
\]

Из тога за првих \(r\) елементов \(x_1,\ldots,x_r\) базы од \(\mathfrak M\) следујут једначеньа

\[
x_i=z_i-\sum_{j=r+1}^{n}x_j\chi_j^{(i)},\qquad
z_i\in\mathfrak C,\quad \chi_j^{(i)}\in\mathsf K.
\]

Понеже \(z_i\) заједно с \(x_{r+1},\ldots,x_n\) очивидно творят базу од \(\mathfrak M_{\mathsf K}\), они сут линеарно независне; зато \(z_1,\ldots,z_r\) јест база од \(\mathfrak C\). \(z_1,\ldots,z_r\) јест база исканој формы, чији \(z_i\) належет к \(\mathfrak M\):

Нека \(\gamma\) буде елемент из \(\mathfrak G\). Тада с једној страны

\[
\begin{array}{rcl}
(1)\qquad \gamma z_i
&=&\gamma x_i+\displaystyle\sum_{j=r+1}^{n}\gamma x_j\cdot \gamma\chi_j^{(i)}
=x_i+\displaystyle\sum_{j=r+1}^{n}x_j\gamma\chi_j^{(i)}.
\end{array}
\]

а с другој страны
\[
\begin{array}{rcl}
(2)\qquad \gamma z_i&=&\displaystyle\sum_{\mu=1}^{r}z_\mu\sigma_\mu^{(i)},
\qquad \sigma_\mu^{(i)}\in\mathsf K,
\end{array}
\]

бо \(\mathfrak C\) допушча все \(\gamma\in\mathfrak G\). (2) можно јест записати јешче в форме

\[
\begin{array}{rcl}
(3)\qquad \gamma z_i
&=&\displaystyle\sum_{\mu=1}^{r}x_\mu\sigma_\mu^{(i)}
+\displaystyle\sum_{j=r+1}^{n}x_j\sum_{\mu=1}^{r}\chi_j^{(\mu)}\sigma_\mu^{(i)}
\end{array}
\]

Сравньенье коефициентов в (1) и (3) даваје

\[
\sigma_j^{(i)}=
\begin{cases}
1 & i=j,\\
0 & i\ne j.
\end{cases}
\]

Зато јест

\[
\gamma\cdot z_i=z_i.
\]

Теды \(z_i\) належет к инвариантној области од \(\mathfrak G\), т.~ј. к \(\mathfrak M\), что было доказати.
% END BOOK_S18

% BEGIN BOOK_S19
\subsection*{§ 19. Представјеньа двустранно простых хиперкомплексных системов в некомутативных телах, кторе разширјајут јих польа коефициентов}\label{darstellungen-zweiseitig-einfacher-hyperkomplexer-systeme-in-nichtkommutativen-erweiterungskoerpern-ihrer-koeffizientenkoerper}

\[
\mathfrak S=x_1\mathsf P+\cdots+x_n\mathsf P
\]

нека буде двустранно просты, теды полно редуцибилны хиперкомплексны систем, а \(\mathsf K\) тело с центром \(\mathsf P\).

\textbf{Теорем 1.} \emph{\(\mathfrak S_{\mathsf K}\) јест двустранно просто (теды такоже полно редуцибилно).}

\emph{Доказ} 1. Всакы двустранны \(\mathfrak S_{\mathsf K}\)-идеал \(\mathfrak a\) јест инвариантны взходно к групе \(\mathfrak G\) внутрных автоморфизмов од \(\mathsf K\). Бо ако, под \(\chi\) разумејучи ненуловы елемент од \(\mathsf K\), положимо \(\chi^{-1}\mathfrak a\chi=\bar{\mathfrak a}\), то \(\bar{\mathfrak a}\subseteq\mathfrak a\) заради својства идеала од \(\mathfrak a\); але понеже \(\bar{\mathfrak a}\), разматраны како \(\mathsf K\)-модул, имаје над \(\mathsf K\) тојже ранг како \(\mathfrak a\), мора бити \(\bar{\mathfrak a}=\mathfrak a\).

\textbf{2.} По управо установјеној помочној леми двустранны \(\mathfrak S_{\mathsf K}\)-идеал \(\mathfrak a\) мора теды бити разширјены идеал некого подмодула \(\mathfrak A\) од \(\mathfrak S\), \(\mathfrak a=\mathfrak A_{\mathsf K}\). \(\mathfrak A\) јест двустранны идеал од \(\mathfrak S\), бо из \(\mathfrak A=\mathfrak S\cap\mathfrak a\) достава се
\[
\mathfrak S\mathfrak A\subseteq\mathfrak S\mathfrak S\cap\mathfrak S\mathfrak a\subseteq\mathfrak S\cap\mathfrak a=\mathfrak A
\]
\[
\mathfrak A\mathfrak S\subseteq\mathfrak S\mathfrak S\cap\mathfrak a\mathfrak S\subseteq\mathfrak S\cap\mathfrak a=\mathfrak A.
\]

Але понеже \(\mathfrak S\) садржује само двустранне идеалы 0 и \(\mathfrak S\), в \(\mathfrak S_{\mathsf K}\) такоже постојут само два двустранна идеала 0 и \(\mathfrak S_{\mathsf K}\), что было доказати.

Ныне разматрајемо иредуцибилна представјеньа од \(\mathfrak S\) в \(\mathsf K\), и хочемо установити реципрочна представјеньа, что имаје формалну предност. В конечном јест всеједно, избирајемо ли према или реципрочна представјеньа, бо она сут једнозначно взајемно приряđена.

Нека \(\mathfrak M\) буде реципрочны модул представјеньа од \(\mathfrak S\) в \(\mathsf K\), теды левы \(\mathfrak S\)-модул и левы \(\mathsf K\)-модул конечного ранга,
\[
\mathfrak M=\mathsf K\cdot y_1+\cdots+\mathsf K\cdot y_r
\]
с законом комутативности
\[
s(xm)=x(sm).
\]

Доказујемо, же \(\mathfrak M\) можно такоже разматрјати како левы \(\mathfrak S_{\mathsf K}\)-модул. Најпрво треба објаснити множенье елемента од \(\mathfrak S_{\mathsf K}\) елементом од \(\mathfrak M\). Ако \(s\in\mathfrak S,\ x\in\mathsf K,\ m\in\mathfrak M\), установјујемо
\[
sx\cdot m=s\cdot(x\cdot m)
\]
и соответно за обчи \(\mathfrak S_{\mathsf K}\)-елемент \(\sum s_ix_i\)
\[
\sum s_ix_i\cdot m=\sum s_i\cdot(x_i\cdot m).
\]

Можно јест уверити се, же та дефиниција јест једнозначна и не зависи од начина, на кторы елемент од \(\mathfrak S_{\mathsf K}\) записујемо како суму продуктов \(s_ix_i\).

Треба теды из
\[
\begin{array}{rcl}
(1)\qquad \displaystyle\sum_{i=1}^{p}s_ix_i&=&0
\end{array}
\]
закльучити
\[
\sum_{i=1}^{p}s_i(x_i m)=0
\]
Предпокладамо, же \(s_1,\ldots,s_q\) сут линеарно независне и же \(s_{q+1},\ldots,s_p\) можно изразити посредством \(s_1,\ldots,s_q\)
\[
s_j=\sum_{i=1}^{q}s_i\varrho_i^{(j)},\qquad j=q+1,\ldots,p.
\]

Коефициенты морајут лежати в \(\mathsf P\), бо \(s_j\) можно изразити само на једен начин, а запис јест уже можны в \(\mathfrak S\). Из (1) достава се

\[
x_i+\sum_{j=q+1}^{p}\varrho_i^{(j)}x_j=0,\qquad i=1,\ldots,q.
\]

Теды јест

\[
\begin{aligned}
\sum_{i=1}^{p}s_i(x_i m)
&=s_1(x_1m)+\cdots+s_q(x_qm)
  +\left(\sum_i s_i\varrho_i^{(q+1)}\right)(x_{q+1}m)\\
&\quad+\cdots+\left(\sum_i s_i\varrho_i^{(p)}\right)(x_pm)\\
&=s_1(x_1m)+\cdots+s_q(x_qm)\\
&\quad+s_1\varrho_1^{(q+1)}(x_{q+1}m)+\cdots+s_q\varrho_q^{(q+1)}(x_{q+1}m)
  \qquad\text{Zaradi svojstva modula}\\
&\quad+\cdots\\
&\quad\cdots\\
&\quad+s_1\varrho_1^{(p)}(x_pm)+\cdots+s_q\varrho_q^{(p)}(x_pm)\\
&=s_1(x_1m)+\cdots+s_q(x_qm)\\
&\quad+s_1(\varrho_1^{(q+1)}x_{q+1}m)+\cdots+s_q(\varrho_q^{(q+1)}x_{q+1}m)\\
&\quad+\cdots\\
&\quad\cdots\\
&\quad+s_1(\varrho_1^{(p)}x_pm)+\cdots+s_q(\varrho_q^{(p)}x_pm)
  \qquad\begin{gathered}\text{Zaradi zakona}\\ \text{asociativnosti }s'(sm)\\ =s's\cdot m\text{ i}\\ x'(xm)=x'x\cdot m\end{gathered}\\
&=s_1\left(x_1m+\varrho_1^{(q+1)}x_{q+1}\cdot m+\cdots+\varrho_1^{(p)}x_p\cdot m\right)+\cdots\\
&=s_1\left(\left(x_1+\varrho_1^{(q+1)}x_{q+1}+\cdots+\varrho_1^{(p)}x_p\right)m\right)+\cdots\\
&=0.
\end{aligned}
\]


Из својств \(\mathfrak S_{\mathsf K}\)-модула од \(\mathfrak M\) дистрибутивне законы \((S+S')m=Sm+S'm,\ S(m+m')=Sm+Sm'\) сут јасне. Јешче треба обратити вниманье на закон асоциативности \(S(S'm)=SS'\cdot m\); понеже дистрибутивне законы важе, можно јест ограничити се на случај

\[
sx\cdot(s'x' m)=sx(s'x'\cdot m)
\]

С једној страны, по дефиницији од \(sx\cdot m\), јест

\[
sx\cdot(s'x'\cdot m)=s\cdot(x\cdot(s'\cdot(x'\cdot m)))
\]

а с другој страны

\[
sx s'x'\cdot m=ss'xx'\cdot m=ss'\cdot(xx'\cdot m)
=s\cdot(s'\cdot(x\cdot(x'\cdot m)))
\]

заради двох релациј комутативности \(=s\cdot(x\cdot(s'\cdot(x'\cdot m)))\)

\[
x\cdot s'=s'\cdot x
\]
\[
s'\cdot(x\cdot m)=x\cdot(s'\cdot m)
\]

Теды \(\mathfrak M\) наистино јест левы \(\mathfrak S_{\mathsf K}\)-модул, очивидно конечны. Обратно, всакы конечны левы \(\mathfrak S_{\mathsf K}\)-модул јест реципрочны модул представјеньа од \(\mathfrak S\) в \(\mathsf K\), бо очивидно јест конечны \(\mathsf K\)-модул и заради

\[
s(x\cdot m)=sx\cdot m=xs\cdot m=x(s\cdot m)
\]

такоже задовольава услов комутативности. (Сравньај такоже простејши доказ в § 24).

Ныне по M.Z. § 18 всакы конечны \(\mathfrak S_{\mathsf K}\)-модул јест сума простых модулов, кторе сут операторно изоморфне с простыми левыми идеалами од \(\mathfrak S_{\mathsf K}\) взходно к операторам из \(\mathfrak S_{\mathsf K}\). Обратно, всакы левы идеал од \(\mathfrak S_{\mathsf K}\) јест модул представјеньа од \(\mathfrak S\) в \(\mathsf K\), а понеже вси леви идеалы од \(\mathfrak S_{\mathsf K}\) сут операторно изоморфне, важи

\textbf{Теорем 2.} \emph{Једина класа иредуцибилных реципрочных представјень од \(\mathfrak S\) в \(\mathsf K\) јест она, ктора јест створјена простым левым идеалом од \(\mathfrak S_{\mathsf K}\).}

Ако то иредуцибилно представјенье имаје ступень \(r\), то ступень либовольного представјеньа јест кратны од \(r\); очивидно јест можно надовезиваньем иредуцибилных представјень за заданы кратны \(rs\) од \(r\) конструовати представјенье \(rs\)-того ступеньа. Тому одповедаје образованье модула представјеньа, кторы јест према сума од \(s\) простых модулов.

Ныне уже знајемо једно представјенье од \(\mathfrak S\) в \(\mathsf K\), јменно оно, кторо јест створјено самым \(\mathfrak S\) како модулом представјеньа (главно представјенье, M. Z. § 20); оно уже лежи в \(\mathsf P\). Ньегов ступень јест \(n\), теды \(n\) јест деливы посредством \(r\):

\[
n=rt.
\]

Квоциент \(t\) јест легко определити: понеже \(r\) јест ранг једного левого идеала над \(\mathsf K\), а \(n\) ранг целокупного \(\mathfrak S_{\mathsf K}\), \(t\) јест число левых идеалов од \(\mathfrak S_{\mathsf K}\), теды \(t^2\) јест ранг од \(\mathfrak S_{\mathsf K}\) над телом автоморфизмов ньеговых простых идеалов.

Специализујемо наше разматранье на случај, когда \(\mathfrak S\) јест комутативно полье.

\textbf{Теорем 3.} \emph{Ако \(\mathfrak S\) јест комутативно полье, то \(\mathfrak S\) јест центар од \(\mathfrak S_{\mathsf K}\).}

\emph{Доказ.} Понеже всакы елемент од \(\mathfrak S\) јест комутативен с всакым елементом од \(\mathfrak S\) и с всакым елементом од \(\mathsf K\), теды такоже с всакым елементом од \(\mathfrak S_{\mathsf K}\), \(\mathfrak S\) належи к центру од \(\mathfrak S_{\mathsf K}\). Елемент од \(\mathfrak S_{\mathsf K}\) јест комутативен с всакым елементом од \(\mathfrak S_{\mathsf K}\) само тогда, ако јест комутативен хоти с всакым елементом од \(\mathsf K\), теды ако ньегови коефициенты належет к центру од \(\mathsf K\), кторы јест \(\mathsf P\), теды ако он лежи в \(\mathfrak S\), что было доказати.

Понеже при комутативных колцах не имаје разлике меджу реципрочным и премым изоморфизмом, всако реципрочно представјенье можно јест сматрјати такоже премым. Иредуцибилно представјенье од \(\mathfrak S\) в \(\mathsf K\) мора теды такоже бити дано премым модулом представјеньа. Тој премы модул представјеньа јест просты правы идеал од \(\mathfrak S_{\mathsf K}\). Бо \(\mathfrak r\) јест конечны правы \(\mathsf K\)-модул ранга \(r\) и заради \(\mathfrak r\mathfrak S=\mathfrak S\mathfrak r\) такоже левы \(\mathfrak S\)-модул (саморазумно с смјешаным законом асоциативности), теды премы модул представјеньа ранга \(r\); а понеже постојí само једна класа иредуцибилных представјень, он мора стварјати дано представјенье.
% END BOOK_S19

% BEGIN BOOK_S20
\subsection*{§ 20. Некомутативне тела}\label{nichtkommutative-koerper}

Ныне под

\[
\mathfrak K=y_1\mathsf P+\cdots+y_m\mathsf P
\]

разумејемо некомутативно тело с центром \(\mathsf P\). За \(\mathfrak K\) развијемо ред теоремов, аналогных теоремам из теорије комутативных конечных алгебраичных разширјень поль.

Ако \(Z=\xi_1\mathsf P+\cdots+\xi_n\mathsf P\) означаје хиперкомплексны систем над \(\mathsf P\), два разширјена система \(\mathfrak K_Z\) и \(Z_{\mathfrak K}\) сут идентичне, бо за \(\mathfrak K_Z\) можно јест записати

\[
\mathfrak K_Z=y_1\xi_1\mathsf P+\cdots+y_m\xi_n\mathsf P
\]

а за \(Z_{\mathfrak K}\)

\[
Z_{\mathfrak K}=\xi_1y_1\mathsf P+\cdots+\xi_ny_m\mathsf P
\]

але оба израза сут идентичне заради \(y_i\xi_j=\xi_jy_i\).

В следујучем \(Z\) всечасно буде комутативно полье, зато можно јест применити теоремы из § 19. По теорему 1 из § 19 \(\mathfrak K_Z\) јест всечасно двустранно просто.

\textbf{Теорем 1.} \emph{Нека \(\Omega\) буде алгебраично затворјено разширјено полье од \(\mathsf P\). Тада центар од \(\mathfrak K_\Omega\) јест равны \(\Omega\).}

\emph{Доказ.} Осланьамо се на одповедајучи теорем 3 из § 19. Јасно јест, же \(\Omega\) лежи в центру од \(\mathfrak K_\Omega\). Ако \(\alpha\) јест либовольны елемент центра од \(\mathfrak K_\Omega\), ньегови коефициенты належет к некому конечному алгебраичному разширјеному польу \(Z\) од \(\mathsf P\) (јменно к польу, кторо они стварјајут). Саморазумно \(\alpha\) такоже мора належати к центру од \(\mathfrak K_Z\). Але тој центар јест \(Z\), теды \(\alpha\) належи к \(Z\), а зато такоже к \(\Omega\), что было доказати.

На совсем подобны начин доказујемо

\textbf{Теорем 2.} \emph{\(\mathfrak K_\Omega\) јест двустранно просто.}

\emph{Доказ.} Нека \(\mathfrak a=z_1\Omega+\cdots+z_r\Omega\) буде двустранны идеал. Коефициенты од \(z_1,\ldots,z_r\) належет к некому конечному разширјеному польу \(Z\) од \(\mathsf P\); в \(\mathfrak K_Z\) јест \(\mathfrak a'=z_1Z+\cdots+z_rZ\) двустранны идеал, а понеже \(\mathfrak K_Z\) јест двустранно просто, \(\mathfrak a'\) јест или нула или цело \(\mathfrak K_Z\); теды \(\mathfrak a\) јест или 0 или цело \(\mathfrak K_\Omega\), т.~ј. \(\mathfrak K_\Omega\) јест двустранно просто, что было доказати.

По том теорему \(\mathfrak K_\Omega\) јест матрицно колцо над некым разширјеным польем својего центра \(\Omega\). То полье како подколцо целого \(\mathfrak K_\Omega\) мора имети конечны ступень над \(\Omega\), теды мора бити равно \(\Omega\):

\textbf{Теорем 3.} \emph{\(\mathfrak K_\Omega\) јест матрицно колцо над \(\Omega\),}

\[
\mathfrak K_\Omega=\sum\Omega c_{ik}
\]

\emph{зато ступень од \(\mathfrak K\) над ньеговым центром \(\mathsf P\) јест квадратно число \(m=t^2\).}

Число \(t\) правых (или левых) идеалов, на кторе разпада се \(\mathfrak K_\Omega\), называјмо кратко абсолутным числом компонентов од \(\mathfrak K\).

Својство од \(\Omega\), же \(\mathfrak K_\Omega\) разпада се на \(t\) простых идеалов, очивидно имајут уже конечна разширјеньа \(Z\) од \(\mathsf P\), напр. така \(Z\), кторе се доставајут приложеньем коефициентов од \(c_{ik}\) к \(\mathsf P\). Ныне разматрајемо та разпадна польа.

\textbf{Дефиниција.} Конечно алгебраично разширјено полье \(Z\) од \(\mathsf P\) называје се \emph{разпадным польем} од \(\mathfrak K\), ако число простых правых идеалов, на кторе разпада се \(\mathfrak K_Z\), уже јест равно абсолутному числу компонентов \(t\)\srcfn{1}{Сравньај појем разпадного польа на стр.~9}.

\textbf{Теорем 4.} \emph{Ако \(Z\) јест разпадно полье од \(\mathfrak K\), то \(Z\) јест тело автоморфизмов правых идеалов од \(\mathfrak K_Z\); обратно, ако тело автоморфизмов правых идеалов од \(\mathfrak K_Z\) јест равно \(Z\), то \(Z\) јест разпадно полье.}

\emph{Доказ.} 1. Нека \(Z\) буде разпадно полье, а \(\mathfrak T\) тело автоморфизмов правых идеалов од \(\mathfrak K_Z\), теды

\[
\mathfrak K_Z=\sum\mathfrak T c_{ik}
\]

Из тога следује

\[
\mathfrak K_\Omega=\sum\mathfrak T_\Omega c_{ik}
\]

ранг од \(\mathfrak K_\Omega\) над \(\Omega\) јест теды продукт од \(t^2\) с рангом од \(\mathfrak T_\Omega\); с другој страны он јест равны \(t^2\), зато ранг од \(\mathfrak T_\Omega\), а такоже од \(\mathfrak T\), јест равны 1, \(\mathfrak T=Z\), что было доказати.

\textbf{2.} Ако \(Z\) јест тело автоморфизмов простых идеалов од \(\mathfrak K_Z\), то јест

\[
\mathfrak K_Z=\sum Zc_{ik}
\]

теды

\[
\mathfrak K_\Omega=\sum\Omega c_{ik}
\]

число елементов \(c_{ik}\) мора теды бити равно \(t^2\), что было доказати.

\textbf{Теорем 5.} \emph{Иредуцибилно представјенье \(\mathfrak Z\) разпадного польа \(Z\) од \(\mathfrak K\) в \(\mathfrak K\) -- по теорему 2 на стр.~22 постојí само једна класа иредуцибилных представјень -- јест максимално комутативно подполье матрицного колца \(\mathfrak K_r\) (\(r\) јест ступень иредуцибилного представјеньа од \(Z\)).}

\emph{Ако обратно иредуцибилно представјенье \(\mathfrak Z\) комутативного разширјеного польа \(Z\) од \(\mathsf P\) јест максимално комутативно подполье од \(\mathfrak K_r\), то \(Z\) јест разпадно}

\emph{полье.}

\emph{Доказ.} 1. Ако \(n\) јест ступень разпадного польа \(Z\) над \(\mathsf P\), то ступень \(r\) иредуцибилного представјеньа \(\mathfrak Z\) од \(Z\) јест равны \(n/t\) (стр.~22). Нека ныне \(\mathfrak Z\subseteq\mathfrak Z^*\subseteq\mathfrak K_r\), а \(\mathfrak Z^*\) буде максимално комутативно полье ступеньа \(n^*\geq n\) над \(\mathsf P\). Треба доказати \(\mathfrak Z^*=\mathfrak Z\). Вместо тога достаточно јест доказати \(n^*=n\). Ако \(r^*\) означаје ступень иредуцибилного представјеньа с \(\mathfrak Z^*\) изоморфного разширјеного польа \(Z^*\) од \(Z\), знову јест \(n^*=r^*t\); с другој страны \(r\) јест деливы посредством \(r^*\), бо в \(\mathfrak Z^*\) јест дано представјенье \(r\)-того ступеньа од \(Z^*\), зато \(r^*\leq r\). Из

\[
n\leq n^*=r^*t\leq rt=n \qquad\text{slěduje}\qquad n^*=n.
\]

2. Нека иредуцибилно представјенье \(\mathfrak Z\) \(r\)-того ступеньа комутативного разширјеного польа \(Z\) од \(\mathsf P\) в \(\mathfrak K\) буде максимално комутативно подполье од \(\mathfrak K_r\). Нека тело автоморфизмов простых идеалов од \(\mathfrak K_Z\) буде \(\mathfrak T\), теды

\[
\mathfrak K_Z=\sum_{1}^{t'}\mathfrak T c_{ik}.
\]

\(\mathfrak T\) садржује \(Z\). Треба доказати \(\mathfrak T=Z\). Ныне просты правы идеал \(\mathfrak r\) од \(\mathfrak K_r\) не јест само модул представјеньа од \(Z\) в \(\mathfrak K\), але такоже модул представјеньа од \(\mathfrak T\), бо јест

\[
\mathfrak r=\mathfrak T c_{i1}+\cdots+\mathfrak T c_{ik},
\]

теды \(\mathfrak T\mathfrak r=\mathfrak r\); \(\mathfrak r\) јест левы \(\mathfrak T\)-модул и очивидно правы \(\mathfrak K\)-модул.

Постојí теды с \(\mathfrak T\) изоморфно подтело \(\overline{\mathfrak T}\) од \(\mathfrak K_r\), кторо садржује \(\mathfrak Z\). Ако бы \(\mathfrak T\) было различно од \(Z\), теды \(\overline{\mathfrak T}\) различно од \(\mathfrak Z\), приложенье не належечего к \(\mathfrak Z\) елемента од \(\overline{\mathfrak T}\), \(\alpha\), к \(\mathfrak Z\) дало бы право комутативно разширјено полье од \(\mathfrak Z\) в \(\mathfrak K_r\), что противречи максималности од \(\mathfrak Z\), что было доказати.

Из теорема 5 безпосреднье следује

\textbf{Теорем 6.} \emph{Максимална комутативна подпольа од \(\mathfrak K\) имајут ступень \(t\) над \(\mathsf P\) и сут разпадна польа\srcfn{1}{Предположенье, же всако разпадно полье такоже садржује неко максимално комутативно подполье од \(\mathfrak K\), не јест правилно; могут постојати разпадна польа без тога својства, а одповедајучи \(r\) сигурно јест выше од 1. Види: R. Брауер и E. Ноетхер, Üбер минимале Зерфäллунгскöрпер ирредузиблер Дарстеллунген; Ситзунгсберицхте д. преусс. Академие дер Висссцх. 1927 XXXII, стр.~221.}.}

Ныне развијемо теорем о изоморфизму подколц једного \(\mathfrak K_r\), на ктором основана јест подробнėјша знајемост структуры Галоисовој групы од \(\mathfrak K\) и с тым связаных фактов.

\textbf{Теорем 7.} \emph{Нека \(\mathfrak S_1\) и \(\mathfrak S_2\) буду двустранно просте, \(\mathsf P\) садржујуче подколца матрицного колца \(\mathfrak K_r\) над \(\mathfrak K\), и нека меджу ними постојí изоморфизм, оставјајучи \(\mathsf P\) по елементах неподвижным, \(\mathfrak S_1\cong\mathfrak S_2\). Тада тој изоморфизм јест трансформација регуларным \(\mathfrak K_r\)-елементом \(\tau\).}

\emph{Ако \(\sigma_1\in\mathfrak S_1\) и \(\sigma_2\in\mathfrak S_2\) сут взајемно приряđене, то јест}

\[
\tau^{-1}\sigma_1\tau=\sigma_2.
\]

\emph{Доказ.} Конструујемо с \(\mathfrak S_1\) и \(\mathfrak S_2\) реципрочно изоморфно колцо \(\Sigma\), чије подполье изоморфно с \(\mathsf P\) идентификујемо с \(\mathsf P\). \(\mathfrak S_1,\mathfrak S_2\) сут тада два реципрочна представјеньа једнакого ступеньа \(r\) од \(\Sigma\) в \(\mathfrak K\). Понеже по теорему 2 на стр.~22 постојí само једна класа иредуцибилных реципрочных представјень, теды такоже само једна класа представјень заданого ступеньа од \(\Sigma\) в \(\mathfrak K\) (кратны од иредуцибилној класы), \(\mathfrak S_1\) и \(\mathfrak S_2\) морајут належати к истој класе; теды једно настава из другого трансформацијоју, что было доказати.

Из теорема 7 непосреднье следује

\textbf{Теорем 8.} \emph{Група автоморфизмов, кторе оставјајут \(\mathsf P\) по елементах неподвижным, од \(\mathfrak K\) -- т.~ј. Галоисова група \(\mathfrak G\) од \(\mathfrak K\) взходно к ньеговому центру -- јест група внутрных автоморфизмов од \(\mathfrak K\).}

В теорему 7 треба само положити \(r=1\) и \(\mathfrak S_1=\mathfrak S_2=\mathfrak K\).

Следујуча два теорема показывајут употребльивост досега развитој обчеј теорије за разматранье специалных пытань.

\textbf{Теорем 9.} \emph{Једино некомутативно тело конечного ступеньа, чији центар јест полье \(R\) реалных числ, јест тело кватернионов.}

\emph{Доказ.} Нека \(\mathfrak K\) буде некомутативно тело с центром \(R\) (конечного ранга над \(R\)). По теорему 6 ступень од \(\mathfrak K\) над \(R\) јест квадрат ступеньа \(t\) некого максималного комутативного подпольа од \(\mathfrak K\); а понеже оно може бити само полье изоморфно с польем комплексных числ, јест \(t=2\), а \(\mathfrak K\) имаје ступень 4 над \(R\). Нека \(\mathfrak Z_1\) буде максимално комутативно подполье од \(\mathfrak K\); тада можно јест положити \(\mathfrak Z_1=R(i)\), гдје \(i^2=-1\). Але тада такоже \(\mathfrak Z_2=R(-i)\) јест -- како множство совпадајуче с \(\mathfrak Z_1\) -- полье изоморфно с \(\mathfrak Z_1\), а изо-

морфизм јест заданы посредством \(i\mapsto -i\). По теорему 7 постојí теды елемент \(j^*\in\mathfrak K\) с

\[
j^{*-1}ij^*=-i.
\]

Ранг од \(\mathfrak Z_1(j^*)\) над \(\mathfrak Z_1\) јест ныне најменеј 2, бо \(j^*\) не може належати к \(\mathfrak Z_1\). Теды \(\mathfrak Z_1(j^*)\) имаје ранг 4 над \(R\), јест равно \(\mathfrak K\), и можно јест положити:

\[
\mathfrak K=R+Ri+Rj^*+Rij^*.
\]

Заради

\[
\begin{aligned}
j^{*-2}1j^{*2}&=1,\\
j^{*-2}ij^{*2}&=i,\\
j^{*-2}j^*j^{*2}&=j^*,\\
j^{*-2}ij^*j^{*2}&=ij^*
\end{aligned}
\]

\(j^{*2}\) јест комутативен с всими елементами од \(\mathfrak K\), належи к центру \(R\), а \(j^{*2}\) јест реално число \(r\). \(r\) не може бити позитивно, бо иначе полином \(x^2-r\) в комутативном польу \(R(j^*)\) имел бы четыри кореньа \(+\sqrt r,-\sqrt r,+j^*\) и \(-j^*\). Теды \(j^{*2}=-m^2,\ m\in R\). Положимо \(j=j^*m^{-1}\); тада \(j^2=-1\), а \(1,i,j,k=ij\) сут јединице кватернионов, что было доказати.

Доказ очивидно важи за всако алгебраично реално затворјено полье вместо \(R\).

\textbf{Теорем 10 (теорем од Веддербурна).} \emph{Всако тело из конечного числа елементов јест комутативно.}

\emph{Доказ.} Нека \(\mathfrak K\) буде тело с конечным числом елементов, а \(\mathsf P\) јего центар. Нека \(t^2\) буде ступень од \(\mathfrak K\) над \(\mathsf P\). Вса максимална комутативна подпольа од \(\mathfrak K\) имајут ступень \(t\) над \(\mathsf P\), теды имајут једнако много елементов; по познатом теорему из теорије Галоисовых поль (теорем од Моореа) она сут изоморфне, а те изоморфизмы, понеже всако Галоисово полье јест Галоисово разширјенье својего простого польа, можно јест избрати тако, же всакы елемент од \(\mathsf P\) одповедаје самому себе. По теорему 7 вса максимална комутативна подпольа од \(\mathfrak K\) наставајут теды из једного \(\mathfrak Z\) трансформацијоју. Але \(\mathfrak K\) јест множство сјединеньа всих својих максималных комутативных подполь. (Приложенье \emph{једного} елемента к \(\mathsf P\) стварјаје комутативно полье). Ако ныне разматрамо само мултипликативну групу \(\mathfrak K^*\) од \(\mathfrak K\) (\(\mathfrak K\) без нулы), достали јесмо, же \(\mathfrak K^*\) јест множство сјединеньа подгруп коньугованых подгрупе \(\mathfrak Z^*\). Але в случају \(t>1\) то даваје противречност. Бо ако \(\mathfrak K^*\) садржује \(M\) елементов, \(\mathfrak Z^*\) \(N\) елементов, а \(L\) јест индекс од \(\mathfrak Z^*\) в \(\mathfrak K^*\), то важи \(M=NL\), а число различных коньугованых подгруп јест највыше \(L\). Чак ако то число было бы равно \(L\), да бы все коньуговане подгрупы заједно изчрпале целу групу, никаке две из ньих не смели бы имети обчи елемент. Але все имајут обчи елемент јединице.
% END BOOK_S20

% BEGIN BOOK_S21
\subsection*{§ 21. Галоисова теорија некомутативных тел}\label{die-galoissche-theorie-der-nichtkommutativen-koerper}

Под \(G\) разумејемо Галоисову групу од \(\mathfrak K\) взходно к ньеговому центру \(\mathsf P\). По теорему 8 \(G\) јест изоморфна с \(\overline G=\mathfrak K^*/\mathsf P^*\); сматрамо тој изоморфизм установјеным. Подгрупу \(H\) од \(G\) называмо \emph{затворјеној}, ако посредством изоморфизма \(G\simeq\mathfrak K^*/\mathsf P^*\) јеј приряđена подгрупа \(\mathfrak H^*\) од \(\mathfrak K^*\) по доданьу \emph{нулы} поставаје телом.


\textbf{Теорем 11. Главны теорем Галоисовој теорије.} \emph{Затворјене подгрупы \(H\) од \(G\) и тела \(\mathfrak S\) меджу \(\mathsf P\) и \(\mathfrak K\) можно јест једнозначно взајемно прирядити тако, же ако \(H\) и \(\mathfrak S\) одповедајут једно другому, \(H\) јест инвариантна група од \(\mathfrak S\), а \(\mathfrak S\) инвариантно тело од \(H\).}

\emph{Доказ.} Раздельамо наш теорем на три честне теоремы
\begin{enumerate}
\item Инвариантна група \(H\) тела \(\mathfrak S\) јест затворјена.
\item Ако \(\mathfrak S\) имаје инвариантну групу \(H\), то \(\mathfrak S\) јест инвариантно тело од \(H\).
\item Ако \(H\) имаје инвариантно тело \(\mathfrak S\), то \(H\) јест инвариантна група од \(\mathfrak S\).
\end{enumerate}

\begin{enumerate}
\item Инвариантној групе \(H\) тела \(\mathfrak S\) посредством изоморфизма \(G\simeq \mathfrak K^*/\mathsf P^*\) приряđено јест множство \(\mathfrak H^*\) всих ненуловых елементов од \(\mathfrak K\), кторе сут по елементах комутативне с \(\mathfrak S\). Заједно с 0 она очивидно творят тело, т.~ј. \(H\) јест затворјена.

\item Трети честны теорем сводимо на вторы. Телу \(\mathfrak S\) приряđена јест подгрупа \(S\) од \(G\), јменно подгрупа, приналежеча к \(\mathfrak S^*/\mathsf P^*\), од \(G\). Једнако групе \(H\) приналежи подтело \(\mathfrak H\) од \(\mathfrak K\), и јест \(H\simeq \mathfrak H^*/\mathsf P^*\). Ако \(\mathfrak S\) јест инвариантно тело од \(H\), то \(\mathfrak S\) состоји се из всих с \(\mathfrak H\) по елементах комутативных елементов од \(\mathfrak K\); можно јест теды такоже казати: \(S\) јест инвариантна група од \(\mathfrak H\). Соответно: Ако \(H\) јест инвариантна група од \(\mathfrak S\), то \(\mathfrak H\) јест инвариантно тело од \(S\). Зато трети честны теорем можно јест формуловати такоже овако:

\item[\(3'\).] Ако \(\mathfrak H\) имаје инвариантну групу \(S\), то \(\mathfrak H\) јест инвариантно тело од \(S\), а то јест равно второму честному теорему, кторы јешче треба доказати.

\item[3.] Тој доказ тече, осим неколико малых разлика, како в комутативном случају:
\end{enumerate}

Конструујемо с \(\mathfrak K\) реципрочно изоморфно тело \(K\), чији центар идентификујемо с \(\mathsf P\). Вместо автоморфизмов од \(\mathfrak K\) можно јест разматрјати реципрочне изоморфизмы од \(\mathfrak S\) с \(K\).

Знову хочемо доказати помочну лему, аналогну теорему о продолженьу, употребјеному в комутативном случају:

Нека \(\mathfrak S\) и \(\mathfrak T\) буду тела меджу \(\mathsf P\) и \(\mathfrak K\),
\[
\mathsf P \subseteq \mathfrak S \subseteq \mathfrak T \subseteq \mathfrak K.
\]
Тврдимо, же всакы реципрочны изоморфизм од \(\mathfrak S\) на подтело од \(K\) (реципрочно представјенье првого ступеньа од \(\mathfrak S\) в \(K\)) можно јест на најменеј толико начинов продолжити до реципрочных изоморфизмов од \(\mathfrak T\), колико указује ступень \(s\) од \(\mathfrak T\) над \(\mathfrak S\).

Тела \(\mathsf P\), \(\mathfrak S\), \(\mathfrak T\), \(\mathfrak K\) разширјајут се с \(K\)
\[
K\subseteq \mathfrak S_K\subseteq \mathfrak T_K\subseteq \mathfrak K_K.
\]
По теорему 1 на стр.~20 всако \(\mathfrak S_K\) јест двустранно просто и полно редуцибилно, а прости леви идеалы -- меджу собоју операторно изоморфне -- сут модулы представјеньа за једину класу иредуцибилных представјень од \(\mathfrak S\) в \(K\). Понеже та класа представјень имаје ступень 1 -- бо по конструкцији од \(K\) в \(K\) постојут тела реципрочно изоморфне с \(\mathfrak S\) -- вси прости леви идеалы имајут ранг 1, т.~ј. систем јест сума толико левых идеалов, колико указује ньегов ступень над \(\mathsf P\)
\[
\mathfrak S_K=K E_1+\cdots+K E_p,\qquad K E_\nu=\mathfrak S_K E_\nu.
\]

Тада јест
\[
\begin{aligned}
\mathfrak T_K&=\mathfrak T\cdot K=\mathfrak T\cdot \mathfrak S_K,\\
\mathfrak T_K&=\mathfrak T\mathfrak S_K E_1+\cdots+\mathfrak T\mathfrak S_K E_p
              =\mathfrak T_K E_1+\cdots+\mathfrak T_K E_p.
\end{aligned}
\]
Всакы идеал \(\mathfrak T_K\cdot E_\nu\) разпада се на \(s\) простых идеалов, бо
\[
[\mathfrak T_K E_\nu:K]=[\mathfrak T_K E_\nu:K E_\nu]=[\mathfrak T_K E_\nu:\mathfrak S_K E_\nu]\leqq[\mathfrak T_K:\mathfrak S_K]=[\mathfrak T:\mathfrak S]=s
\]
и
\[
\sum [\mathfrak T_K E_\nu:K]=sp
\]
теды
\[
[\mathfrak T_K E_\nu:K]=s,
\qquad
\mathfrak T_K E_\nu=K e_1^{(\nu)}+\cdots+K e_s^{(\nu)}.
\]

Представјенье, дано посредством \(K e_i^{(\nu)}\), од \(\mathfrak T\) јест продолженье представјеньа, заданого посредством \(K E_\nu\), од \(\mathfrak S\), бо ако \(\alpha\in \mathfrak S\), то јест
\[
\alpha E_\nu=\sigma E_\nu,\qquad
\sum \alpha e_i^{(\nu)}=\sigma\sum e_i^{(\nu)},\qquad
\alpha e_i^{(\nu)}=\sigma e_i^{(\nu)}.
\]

Ныне треба доказати само, же представјеньа дане посредством \(K e_i^{(\nu)}\) сут все различне. В комутативном случају то было јасно, бо она належали к различным класам представјень. Тут она, насупрот, все належет к истој класе представјень, зато треба разматрјати саме представјеньа. За то разматрајемо дане посредством \(K e_i^{(\nu)}\) \(K\)-хомоморфизмы целого \(\mathfrak T_K\); они сут уже целком задане, ако знајемо представјеньа од \(\mathfrak T\), бо за конструкцију целого хомоморфизма достаточно јест знати, како сут представльене базове елементи. Можно јест теды ограничити се на доказ, же посредством \(K e_i^{(\nu)}\) дане хомоморфизмы од \(\mathfrak T_K\) сут все различне. Ако \(e_i^{(\nu)}\) једен раз представимо посредством модула представјеньа \(K e_i^{(\nu)}\), а други раз посредством \(K e_j^{(\nu)}\), тада заради
\[
e_i^{(\nu)2}=e_i^{(\nu)}
\]
првы раз најдемо 1, а заради
\[
e_i^{(\nu)}e_j^{(\nu)}=0
\]
други раз 0 како представльајучи елемент од \(e_i^{(\nu)}\); оба представјеньа сут теды наистино различне. (Точно теже закльучкы можно јест применити такоже в комутативном случају, але там сут непотребне).

Од ныне закльучујемо точно како на стр.~10 и стр.~16. Нека \(H\) буде инвариантна група од \(\mathfrak S\), а \(\mathfrak T\) инвариантно тело од \(H\). Елементи од \(H\) сут вса продолженьа идентичного автоморфизма од \(\mathfrak S\) на цело \(\mathfrak K\). Понеже по управо доказаној помочној леми меджу ними постојут најменеј \(s\) различне на \(\mathfrak T\), мора бити \(s=1\), \(\mathfrak T=\mathfrak S\), что было доказати.

\section*{Глава V. Факторне системы}\label{kapitel-v.-faktorensysteme}
% END BOOK_S21

% BEGIN BOOK_S22
\subsection*{§ 22. Група тел с даным центром}\label{die-gruppe-der-koerper-mit-gegebenem-zentrum}

За основу нашего разматраньа изберемо установјено комутативно полье \(\mathsf P\). Множство всих матрицных колц \(\mathfrak K_r\), чија тела автоморфизмов \(\mathfrak K\) сут тела конечного ранга над \(\mathsf P\) с самым \(\mathsf P\) како центром, јест мултипликативно затворјены систем взходно к ,,образованьу премого продукта``. Премы продукт двох такых колц \(\mathfrak K_r\) и \(\mathfrak L_s\) јест просто колцо \(\mathfrak K_r\cdot\mathfrak L_s\), кторо настава из \(\mathfrak K_r\) разширјеньем основного польа до \(\mathfrak L_s\):
\[
\mathfrak K_r\times\mathfrak L_s=\mathfrak K_r\cdot\mathfrak L_s.
\]
По припомненке на стр.~22 то образованье премого продукта јест комутативно:
\[
\mathfrak K_r\times\mathfrak L_s=\mathfrak L_s\times\mathfrak K_r.
\]
Тврдженье, же систем всих \(\mathfrak K_r\) јест мултипликативно затворјены, не значи ничто друго неж, же \(\mathfrak K_r\times\mathfrak L_s=\mathfrak K_r\mathfrak L_s\) јест всечасно знову матрицно колцо, чије тело автоморфизмов \(\mathfrak M\) имаје центар \(\mathsf P\) -- конечны ранг взходно к \(\mathsf P\) разумеје се сам собоју. По теорему 1 на стр.~20 \(\mathfrak K_r\times\mathfrak L_s\) јест двустранно просто и полно редуцибилно, теды матрицно колцо над некым разширјеным телом \(\mathfrak M\) од \(\mathsf P\). Факт, же центар од \(\mathfrak M\) јест равны \(\mathsf P\), непосреднье следује из тога, же прво \(\mathsf P\) сигурно јест садржано в центру, а друго оно мора бити садржано в центру \(\mathsf P\) од \(\mathfrak K\) (или од \(\mathfrak L\)). Множенье \(\mathfrak K_r\times\mathfrak L_s\) очивидно јест асоциативно. \(\mathfrak K_r\) творят мултипликативно затворјены систем, але не групу, бо ако \(\mathfrak K_r\times\mathfrak L_s=\mathfrak M_m\), то сигурно \(m\geq rs\); теды не буде можно најдти \(\mathfrak L_s\), кторо задовольава једначенье \(\mathfrak K_r\times\mathfrak L_s=\mathfrak M_m\), ако \(r>m\).

Ныне сдружујемо все \(\mathfrak K_r\) с једнакым \(\mathfrak K\) в класу \(\{\mathfrak K\}\). Тада важи: Ако \(\mathfrak K_r\) и \(\mathfrak K_{r'}\) належет к истој класе \(\{\mathfrak K\}\), а \(\mathfrak L_s\) и \(\mathfrak L_{s'}\) к истој класе \(\{\mathfrak L\}\), то такоже \(\mathfrak K_r\times\mathfrak L_s\) и \(\mathfrak K_{r'}\times\mathfrak L_{s'}\) належет к истој класе. Достаточно јест показати, же \(\mathfrak K_r\times\mathfrak L_s\) належи к истој класе како \(\mathfrak K\times\mathfrak L\), теды же ако \(\mathfrak K\times\mathfrak L=\mathfrak M_m\), то такоже \(\mathfrak K_r\times\mathfrak L_s\) јест матрицно колцо над \(\mathfrak M\). То јест легко видети: Јест
\[
\mathfrak K_r\times\mathfrak L
=\left(\sum_1^r\mathfrak K c_{ik}\right)\mathfrak L
=\sum_1^r\mathfrak K_{\mathfrak L}c_{ik}
=\sum_1^r\mathfrak M_m c_{ik}.
\]
Матрицно колцо \(r\)-того ступеньа над \(\mathfrak M_m\), а то очивидно јест матрицно колцо \(mr\)-того ступеньа над \(\mathfrak M\). Тада продолжујемо точно тако
\[
\mathfrak K_r\times\mathfrak L_s
=\left(\sum_1^s\mathfrak L d_{ik}\right)_{\mathfrak K_r}
=\sum_1^s\mathfrak L_{\mathfrak K_r}d_{ik}
=\sum_1^s\mathfrak M_{mr}d_{ik}
=\mathfrak M_{mrs}
\]
чим наше тврдженье уже јест доказано. Запомнимо, же из \(\mathfrak L\times\mathfrak K=\mathfrak M_m\) следује \(\mathfrak K_r\times\mathfrak L_s=\mathfrak M_{mrs}\).

Из доказаного следује: ако продукт двох клас \(\{\mathfrak K\}\) и \(\{\mathfrak L\}\) дефинујемо како класу продукта једного репрезентанта од \(\{\mathfrak K\}\) с једним репрезентантом од \(\{\mathfrak L\}\), та дефиниција не зависи од избора репрезентантов.

Приряđенье \(\mathfrak K_r\mapsto\{\mathfrak K\}\) јест теды хомоморфизм система всих \(\mathfrak K_r\) на множство всих клас \(\{\mathfrak K\}\); множство \(\mathscr K\) клас \(\{\mathfrak K\}\) јест теды мултипликативно затворјены систем. Доказујемо, же \(\mathscr K\) јест група. Како се непосреднье види, елемент јединице јест класа \(\{\mathsf P\}\), \(\mathfrak K_r\times\mathsf P=\mathfrak K_r\).

Остаје само доказати за всаку класу \(\{\mathfrak K\}\) постојанье класы \(\{\mathfrak K\}^{-1}\). Показује се, же треба положити \(\{\mathfrak K\}^{-1}=\{\overline{\mathfrak K}\}\), гдје \(\{\overline{\mathfrak K}\}\) означаје тело реципрочно изоморфно с \(\mathfrak K\). Другыми словами: \(\mathfrak K\times\overline{\mathfrak K}\) мора бити матрицно колцо над \(\mathsf P\), или ако \(\mathfrak K\) имаје ранг \(t^2\), теды \(\mathfrak K\times\overline{\mathfrak K}\) ранг \(t^4\) над \(\mathsf P\), то \(\mathfrak K\times\overline{\mathfrak K}\) мора разпадати се на \(t^2\) простых левых идеалов. То следује из тога, же просты левы идеал од \(\mathfrak K\times\overline{\mathfrak K}\), како модул представјеньа некого иредуцибилного реципрочного представјеньа од \(\overline{\mathfrak K}\) в \(\mathfrak K\) -- кторо јест првого ступеньа -- имаје ранг 1 над \(\mathfrak K\), теды ранг \(t^2\) над \(\mathsf P\).

\(\mathfrak K\) хочемо называти \emph{простым телом}, ако меджу \(\mathsf P\) и \(\mathfrak K\) не имаје тела, чији центар такоже јест \(\mathsf P\).

\textbf{Теорем.} \emph{Всако тело \(\mathfrak K\) јест премы продукт простых подтел \(\mathfrak K^{(i)}\),}
\[
\mathfrak K=\mathfrak K^{(1)}\times\cdots\times\mathfrak K^{(p)}.
\]

\emph{Доказ.} Достаточно јест показати, же ако \(\mathfrak K^{(1)}\) јест подтело од \(\mathfrak K\) с центром \(\mathsf P\), то постојí подтело \(\mathfrak S\) од \(\mathfrak K\) с центром \(\mathsf P\), кторо задовольава једначенье
\[
\mathfrak K^{(1)}\times\mathfrak S=\mathfrak K
\]
В всаком случају постојí класа \(\{\mathfrak S\}\) с својством
\[
\{\mathfrak K^{(1)}\}\times\{\mathfrak S\}=\{\mathfrak K\}
\]
јменно
\[
\{\mathfrak S\}=\{\overline{\mathfrak K^{(1)}}\}\times\{\mathfrak K\}.
\qquad
\text{Neka bude}\qquad
\overline{\mathfrak K^{(1)}}\times\mathfrak K=\mathfrak S_\lambda.
\]
Тада јест
\[
\mathfrak K^{(1)}\times\mathfrak S_\lambda
=\mathfrak K^{(1)}\times\overline{\mathfrak K^{(1)}}\times\mathfrak K
=\mathsf P_{t^2}\times\mathfrak K
=\mathfrak K_{t^2},
\qquad
t^2=\text{rang od }\mathfrak K^{(1)}.
\]
Але \(\mathfrak K^{(1)}\times\mathfrak S_\lambda\) јест матрицно колцо ступеньа најменеј равного \(\lambda\), теды \(\lambda\leq t^2\). \(\mathfrak S_\lambda=\overline{\mathfrak K^{(1)}}\times\mathfrak K\) мора разпадати се на просте идеалы ранга 1 над \(\mathfrak K\), бо \(\overline{\mathfrak K^{(1)}}\) допушча реципрочна представјеньа првого ступеньа в \(\mathfrak K\); зато \(\lambda\geq t^2\), \(\lambda=t^2\),
\[
\mathfrak K_{t^2}=\mathfrak K^{(1)}\times\mathfrak S_{t^2}.
\]
Теды само \(\mathfrak K^{(1)}\times\mathfrak S\) мора бити тело, и
\[
\mathfrak K^{(1)}\times\mathfrak S_{t^2}
=\bigl(\mathfrak K^{(1)}\times\mathfrak S\bigr)_{t^2}
=\mathfrak K_{t^2}.
\]
\(\mathfrak K^{(1)}\times\mathfrak S\) јест тело автоморфизмов простых левых идеалов в једном разкладу од \(\mathfrak K_{t^2}\), а \(\mathfrak K\) тело автоморфизмов простых левых идеалов в другом разкладу од \(\mathfrak K_{t^2}\). Але понеже всака два различна лева идеала од \(\mathfrak K_{t^2}\) можно јест вместити в једен разклад, \(\mathfrak K\) и \(\mathfrak K^{(1)}\times\mathfrak S\) сут изоморфне, а тој изоморфизм можно јест избрати тако, же \(\mathsf P\) оставаје по елементах неподвижно. Зато по теорему 7 на стр.~25 тој изоморфизм јест створјен внутрным автоморфизмом од \(\mathfrak K_{t^2}\).
\[
\mathfrak K
=\lambda^{-1}\bigl(\mathfrak K^{(1)}\times\mathfrak S\bigr)\lambda
=\lambda^{-1}\mathfrak K^{(1)}\lambda\times\lambda^{-1}\mathfrak S\lambda,
\qquad
\lambda\in\mathfrak K_{t^2}.
\]
\(\lambda^{-1}\mathfrak K^{(1)}\lambda\) јест с \(\mathfrak K^{(1)}\) изоморфно подтело од \(\mathfrak K\); зато по теорему 7 на стр.~25 с одповедајучим \(\mu\) из \(\mathfrak K\):
\[
\mu^{-1}\lambda^{-1}\mathfrak K^{(1)}\lambda\mu=\mathfrak K^{(1)}
\]
теды
\[
\mu^{-1}\mathfrak K\mu
=\mathfrak K
=\mathfrak K^{(1)}\times\mu^{-1}\lambda^{-1}\mathfrak S\lambda\mu,
\]
чим јест тврдженье доказано.
% END BOOK_S22

% BEGIN BOOK_S23
\subsection*{§ 23. Факторне системы}\label{faktorensysteme}

В следујучих разматраньах предпокладамо, же основно полье \(\mathsf P\) јест совршено; вместо тога можно јест учинити обчејше предположенье, же разматрајут се само така некомутативна \(\mathfrak K\) с центром \(\mathsf P\), кторе имајут својство, же всако \(\mathfrak K_r\) садржује максимално комутативно подполье првого рода над \(\mathsf P\). (Како меджутым показал Кöтхе, то важи за всако \(\mathfrak K\).)

Под \(Z\) разумејемо разпадно полье некомутативного тела \(\mathfrak K\) с центром \(\mathsf P\); нека иредуцибилно представјенье \(\mathfrak Z\) од \(Z\) в \(\mathfrak K\) буде ступеньа \(r\), тада оно јест максимално комутативно подполье од \(\mathfrak K_r\). Нека \(\Gamma\) буде Галоисово полье од \(Z\), теды
\[
\mathfrak Z_\Gamma=e_1\Gamma+\cdots+e_n\Gamma
\qquad\text{(str.~9)}
\]
\(n\) јест ступень од \(Z\) над \(\mathsf P\). За с \(\Gamma\) разширјено колцо \(\mathfrak K_r\) достава се такоже разклад
\[
\mathfrak K_{r\Gamma}=\mathfrak K_{r\Gamma}e_1+\cdots+\mathfrak K_{r\Gamma}e_n
\]
\[
\mathfrak K_{r\Gamma}=\mathfrak l_1+\cdots+\mathfrak l_n,\qquad
\mathfrak l_i=\mathfrak K_{r\Gamma}e_i.
\]
\(\mathfrak l_i\) очивидно сут леви идеалы. Они сут чак прости леви идеалы, бо понеже \(Z\) јест разпадно полье, \(\mathfrak K_{r\Gamma}\) јест матрицно колцо ступеньа \(r\cdot t=n\) над \(\Gamma\). (Теорем 4, стр.~23, \(n=rt\) на стр.~22.) Тако достаны разклад од \(\mathfrak K_{r\Gamma}\) на просте леве идеалы особливо јест прилагоджены алгебраичным својствам од \(\mathfrak K\) и \(Z\). Ако, како на стр.~18, дејство субституције \(S\) из Галоисовој групы \(\mathfrak G\) од \(\Gamma/\mathsf P\) на все елементы од \(\mathfrak K_{r\Gamma}\) дефинујемо тако, же за всакы \(\mathfrak K_r\)-елемент \(z\) положимо
\[
S(z)=z
\]
то идеалы \(\mathfrak l_i\) сут коньуговане, т.~ј. применьенье једного \(S\) на неко \(\mathfrak l_j\) даваје друго
\[
S(\mathfrak l_j)=\mathfrak l_i,
\]
бо всака неразложива компонента јединице \(e_j\) мора преходити в неку другу \(e_i\). (Види такоже стр.~11.)

Модул представјеньа \(e_i\Gamma\) образује \(\mathfrak Z\) на подполье \(Z_i\) од \(\Gamma\). \(e_i\) уже лежи в \(\mathfrak Z_{Z_i}\), бо представјенье \(Z_i\) од \(\mathfrak Z\) уже јест садржано в \(Z_i\), теды \(\mathfrak Z_{Z_i}\) мора одцепити модул представјеньа \(e_iZ_i\). Нека \(\{Z_i,Z_k\}\) буде композитум од \(Z_i\) и \(Z_k\). \(\mathfrak K_{r\{Z_i,Z_k\}}\) одцепја два лева идеала
\[
\bar{\mathfrak l}_i=\mathfrak K_{r\{Z_i,Z_k\}}e_i,\qquad
\bar{\mathfrak l}_k=\mathfrak K_{r\{Z_i,Z_k\}}e_k
\]
кторе сут просте, бо важи \(\mathfrak l_i=\mathfrak K_{r\Gamma}\bar{\mathfrak l}_i\). Теды јест
\[
\bar{\mathfrak l}_i=\sum_{\lambda=1}^m\{Z_i,Z_k\}c_{\lambda i}^{(ik)},\qquad
\bar{\mathfrak l}_k=\sum_{\lambda}\{Z_i,Z_k\}c_{\lambda k}^{(ik)}.
\]
\(c_{\lambda\rho}^{(ik)}\) јест систем матрицных јединиц од \(\mathfrak K_{r\{Z_i,Z_k\}}\).

Всакы изоморфизм \(a_i\mapsto a_k\) од \(\bar{\mathfrak l}_i\) на \(\bar{\mathfrak l}_k\) с \(\mathfrak K_{r\{Z_i,Z_k\}}\) како областју операторов определьен јест елементом \(p_{ik}\) с \(e_i\mapsto p_{ik}\), бо тада јест \(a_i=a_ie_i\mapsto a_ip_{ik}\), теды \(\bar{\mathfrak l}_i\cdot p_{ik}=\bar{\mathfrak l}_k\). Вси можне таки \(p_{ik}\) очивидно сут
\[
p_{ik}=c_{ik}^{(ik)}\gamma_{ik},\qquad \gamma_{ik}\ne0\quad\text{iz}\quad \{Z_i,Z_k\}.
\]
Посредством \(a_i\mapsto a_ip_{ik}\) тада јест заданы такоже операторны изоморфизм од \(\mathfrak l_i\) на \(\mathfrak l_k\). Ныне избирајемо установјене системы такых \(p_{ik}\), налагајучи јим услов коньугованости.

Дејство једного \(S\) на индекс \(i\) нека буде дефиновано тако
\[
S(i)=k,\qquad\text{ako}\qquad S(Z_i)=Z_k.
\]
За всакы пар индексов \((\nu,\mu)\) постојí најменеј једен коньугованы пар \(S(\nu,\mu)\) формы \(S(\nu,\mu)=(1,k)\). Из всакој класы коньугованых паров индексов избираје се установјены пар \((1,k)\), а тада за тој \((1,k)\):
\[
p_{1k}=c_{1k}^{(1k)}\gamma_{1k},\qquad \gamma_{1k}\ne0\quad\text{iz}\quad \{Z_1,Z_k\},
\]
иначе се положи либовольно, а тада, ако \(S(1,k)=(\nu,\mu)\),
\[
p_{\nu\mu}=S(p_{1k}).
\]
\(p_{\nu\mu}\) тада наистино имаје форму
\[
p_{\nu\mu}=c_{\nu\mu}^{(\nu\mu)}\gamma_{\nu\mu},\qquad
\gamma_{\nu\mu}\ne0\quad\text{iz}\quad \{Z_\nu,Z_\mu\}.
\]
Бо дејство од \(S\) на изоморфизм
\[
a_1\mapsto a_1p_{1k},\qquad\text{specialno}\qquad e_1\mapsto p_{1k}
\quad\text{od }\bar{\mathfrak l}_1\text{ na }\bar{\mathfrak l}_k
\]
даваје изоморфизм
\[
a_\nu\mapsto a_\nu p_{\nu\mu},\qquad\text{specialno}\qquad e_\nu\mapsto p_{\nu\mu}
\quad\text{od }\bar{\mathfrak l}_\nu\text{ na }\bar{\mathfrak l}_\mu.
\]

Ако под \(c_{\lambda\rho}\) разумејемо систем матрицных јединиц од \(\mathfrak K_{r\Gamma}\), јасно јест, же \(p_{\nu\mu}\) можно такоже записати в форме
\[
p_{\nu\mu}=c_{\nu\mu}\bar\gamma_{\nu\mu}
\]
релације
\[
c_{ij}c_{1k}=0,\quad\text{ako}\quad j\ne1,
\qquad
c_{ij}c_{jk}=c_{ik}
\]
имајут теды како последицу релације
\[
\begin{aligned}
p_{ij}p_{1k}&=0,\quad &&\text{ako }j\ne1,\\
p_{ij}p_{jk}&=\alpha_{ik}^{(j)}p_{ik},\quad
&&\alpha_{ik}^{(j)}\text{ iz }\Gamma\text{ (točnėje iz }\{Z_i,Z_j,Z_k\}\text{)}
\end{aligned}
\]
Те \(n^3\) величины \(\alpha\) творят \emph{факторны систем од \(\mathfrak K\) при избраньу од \(Z\)}, кторы приналежи к систему \(p_{ik}\) ,,\emph{псеудоматрицных јединиц}``.

\(\alpha\) сут коньуговане: Из \(S(i,k,j)=(\nu,\mu,\tau)\) следује
\[
S\bigl(\alpha_{ik}^{(j)}\bigr)=\alpha_{\nu\mu}^{\tau}.
\]

Всакы систем \(p_{ik}^*\) псеудоматрицных јединиц настава из некого система \(p_{ik}\) тако, же почетне \(p_{ik}\) множимо либовольными \(\delta_{ik}\ne0\) из \(\{Z_i,Z_k\}\):
\[
p_{1k}^*=p_{1k}\delta_{1k}
\]
и тада знову, ако \(S(1,k)=(\nu,\mu)\), положимо
\[
p_{\nu\mu}^*=S(p_{1k}^*)
\]
Теды јест
\[
p_{\nu\mu}^*=p_{\nu\mu}\delta_{\nu\mu}
\]
с \(\delta_{\nu\mu}\) из \(\{Z_\nu,Z_\mu\}\); \(\delta_{\nu\mu}\) сут коньуговане, из \(S(i,k)=(\nu,\mu)\) следује \(S(\delta_{ik})=\delta_{\nu\mu}\). Ако \(\alpha^*\) јест факторны систем приналежечи к \(p_{ik}^*\), то јест
\[
\alpha_{ik}^{(j)*}=\frac{\delta_{ij}\delta_{jk}}{\delta_{ik}}\alpha_{ik}^{(j)}.
\]

Факторне системы \(\alpha^*\) и \(\alpha\) называјут се \emph{асоциоваными} (једнако \(p_{ik}^*\) и \(p_{ik}\)); вси факторне системы од \(\mathfrak K\) при избраньу од \(Z\) творят класу \(\{\alpha\}\) асоциованых факторных системов.

Класа \(\{\alpha\}\) не зависи од тога, кторо представјенье \(\mathfrak Z\) од \(Z\) в \(\mathfrak K_r\) изберемо за основу јеј дефиниције. Бо преход од једного \(\mathfrak Z\) к другому можно јест учинити трансформацијоју елементом \(\lambda\) из \(\mathfrak K_r\), кторы јест теды \(\mathfrak G\)-инвариантны; зато из система \(p_{ik}\), приналежечего к \(\mathfrak Z\), настава систем
\[
p'_{ik}=\lambda^{-1}p_{ik}\lambda,
\]
кторы приналежи к \(\mathfrak Z'\); факторне системы од \(p_{ik}\) и \(p'_{ik}\) але сут идентичне.

Факторны систем не јест ничто друго неж таблица множеньа од \(\mathfrak K_{n\Gamma}\), прилагоджена особливым алгебраичным својствам од \(\mathfrak K\).
% END BOOK_S23

% BEGIN BOOK_S24
\subsection*{§ 24. Множенье факторных системов}\label{multiplikation-von-faktorensystemen}

Вместо о разпадном польу некомутативного тела \(\mathfrak K\) говоримо такоже о разпадном польу од \(\{\mathfrak K\}\); то имаје смысл, бо сполу с \(\mathfrak K_Z\) такоже всако колцо \(\mathfrak K_{rZ}\) разпада се на абсолутно иредуцибилне компоненты, и обратно.

Всака два тела \(\mathfrak K\), \(\mathfrak S\) с центром \(\mathsf P\) имајут обча разпадна польа, напр. композитум разпадного польа од \(\mathfrak K\) и разпадного польа од \(\mathfrak S\).

\textbf{Теорем 1.} \emph{Класы \(\{\mathfrak K\}\), кторе имајут полье \(Z\) како разпадно полье, творят подгрупу \(\mathscr K_Z\) групы \(\mathscr K\) всих \(\{\mathfrak K\}\).}

\emph{Доказ.} Ако \(\mathfrak K\) и \(\mathfrak S\) имајут \(Z\) како разпадно полье, то \(\mathfrak K_Z\) и \(\mathfrak S_Z\), теды такоже \((\mathfrak K\times\mathfrak S)_Z\), разпадајут се на абсолутно иредуцибилне компоненты, т.~ј. \(Z\) јест разпадно полье од \(\{\mathfrak K\times\mathfrak S\}\) (теорем 4 на стр.~23).

Реципрочно изоморфне тела имајут вса разпадна польа обче.

Помочно разматранье.
\[
\mathfrak K_r=\mathfrak L_1+\cdots+\mathfrak L_r,\qquad
\mathsf P_m=\mathfrak z_1+\cdots+\mathfrak z_m
\]
нека буду разклады од \(\mathfrak K_r\), \(\mathsf P_m\) на просте леве идеалы; тада
\[
\mathfrak K_{rm}=\mathfrak K_r\times\mathsf P_m
=\mathfrak L_1\mathfrak z_1+\cdots+\mathfrak L_i\mathfrak z_j+\cdots+\mathfrak L_r\mathfrak z_m
\]
јест разклад од \(\mathfrak K_{rm}\) на \(rm\) ненуловых, теды простых левых идеалов. Зато всакы \(r\)-членны левы идеал \(A\) од \(\mathfrak K_{rm}\) јест связаны с \(r\)-членным левым идеалом
\[
\mathfrak L_1\mathfrak z_1+\cdots+\mathfrak L_r\mathfrak z_1=\mathfrak K_r\mathfrak z_1
\]
посредством внутрного автоморфизма од \(\mathfrak K_r\)
\[
A=\tau^{-1}\mathfrak K_r\tau\cdot\tau^{-1}\mathfrak z_1\tau
=\mathfrak K'_r\cdot\mathfrak z'_1.
\]

Ако \(\mathfrak K_{r\Gamma}=\mathfrak l_1+\cdots+\mathfrak l_n\) јест разклад од \(\mathfrak K_{r\Gamma}\) на коньуговане леве идеалы, приналежечи к некому \(\mathfrak Z\) в \(\mathfrak K_r\), \(\mathfrak l_i=\mathfrak K_{r\Gamma}e_i\), то заради
\[
S(\tau^{-1}\mathfrak l_i\tau\cdot\mathfrak z'_1)
=S(\tau^{-1}\mathfrak l_i\tau)\,\mathfrak z'_1
=\tau^{-1}(S\mathfrak l_i)\tau\cdot\mathfrak z'_1
=\tau^{-1}\mathfrak l_k\tau\cdot\mathfrak z'_1
\]
в
\[
\tag{1}
A_\Gamma=\mathfrak l'_1\mathfrak z'_1+\cdots+\mathfrak l'_n\mathfrak z'_1
\]
имајемо разклад од \(A_\Gamma\) на коньуговане леве идеалы. Јест
\[
\mathfrak l'_i\mathfrak z'_1=\mathfrak K_{rm\Gamma}\cdot e'_iE',
\]
ако положимо \(\tau^{-1}e_i\tau=e'_i\), а \(E'\) означаје јединицу од \(\mathfrak z'_1\).

Всакы другы разклад
\[
\tag{2}
A_\Gamma=\mathfrak q_1+\cdots+\mathfrak q_n,\qquad
S\mathfrak q_i=\mathfrak q_k\quad\text{ako}\quad
S\mathfrak l'_i\mathfrak z'_1=\mathfrak l'_k\mathfrak z'_1
\]
на коньуговане леве идеалы настава из (1) трансформацијоју с \(\mathfrak K_{rm}\)-елементом \(\lambda\)
\[
\mathfrak q_i=\lambda^{-1}\mathfrak l'_i\mathfrak z'_1\lambda,
\qquad
\text{tedy, ako \(e_i^*\) jest jedinica od \(\mathfrak q_i\),}\quad
e_i^*=\lambda^{-1}e'_iE'\lambda.
\]

Доказ јест непосредньа последица следујучеј помочној лемы:

\textbf{Помочна лема.} \emph{Нека \(A\) буде левы идеал из \(\mathfrak K_n\), и нека}
\[
A=\mathfrak l_1+\cdots+\mathfrak l_n
=\mathfrak K_{n\Gamma}e_1+\cdots+\mathfrak K_{n\Gamma}e_n
\]
\emph{буде -- по предположеньу постојачи -- разклад на коньуговане леве идеалы; при том можно јест -- без ограниченьа обчости -- предпокладати, же \(\mathfrak l_i\) сут коньуговане.}

\emph{Нека \(\mathfrak z\) буде операторно изоморфны с \(A\); и нека важе теже предположеньа:}
\[
\mathfrak z=\mathfrak q_1+\cdots+\mathfrak q_n
=\mathfrak K_{n\Gamma}e_1^*+\cdots+\mathfrak K_{n\Gamma}e_n^*
\]
\emph{(теды \(\mathfrak q_i\) и \(e_i^*\) сут коньуговане). Тада постојí регуларны елемент \(\lambda\) из \(\mathfrak K_n\), тако же \(\mathfrak q_i=\lambda^{-1}\mathfrak l_i\lambda\).}

\emph{Доказ. 1. Факт, же \(e_i\) можно предпокладати коньуговаными, следује тако:} Ако \(e_1\) јест либовольны творечи идемпотент од \(\mathfrak l_1\), теды \(\mathfrak l_1=\mathfrak K_{n\Gamma}e_1\), применьеньем автоморфизма на \(\Gamma\) следује \(\mathfrak l_i=\mathfrak K_{n\Gamma}e_i\), гдје \(e_i\) јест коньугованы с \(e_1\) и знову творечи идемпотент.

2. \emph{Всако \(\mathfrak l\) јест изоморфно с всакым \(\mathfrak q\);} бо \(\mathfrak l\), како коньуговане, все имајут једнаку (абсолутну) должину, числа од \(\mathfrak l\) и \(\mathfrak q\) сут једнака, а \(\mathfrak z\) и \(A\) сут по предположеньу изоморфне.

3. Приряđеньем
\[
\tag{3}
e_1\mapsto s_1=e_1s_1;\qquad
e_1^*\mapsto t_1=e_1^*t_1;\qquad
s_1t_1=e_1;\qquad
t_1s_1=e_1^*
\]
нека буду установјене изоморфизм меджу \(\mathfrak l_1\) и \(\mathfrak q_1\) и ньегов обрат. Применьеньем всих автоморфизмов (3) преходи в
\[
e_i\mapsto s_i=e_is_i;\qquad
e_i^*\mapsto t_i=e_i^*t_i;\qquad
s_it_i=e_i;\qquad
t_is_i=e_i^*,
\]
чим сут теды установјене изоморфизмы \(\mathfrak q_i=\mathfrak l_is_i\) и \(\mathfrak l_i=\mathfrak q_it_i\).

4. Ныне положимо:
\[
s=s_1+\cdots+s_n;\qquad t=t_1+\cdots+t_n;
\]
тада достава се:
\[
\mathfrak q_i=\mathfrak l_is;\qquad
\mathfrak l_i=\mathfrak q_it;\qquad
st=e_1+\cdots+e_n=\xi;\qquad
ts=e_1^*+\cdots+e_n^*=\xi^*,
\]
причем \(\xi\) и \(\xi^*\) лежет в \(\mathfrak K_n\); достава се
\[
A\xi=A,\qquad \mathfrak z\xi^*=\mathfrak z\quad
(\text{zaradi }\mathfrak l_i\xi=\mathfrak l_i\mathfrak l_i=\mathfrak l_i),
\]
\[
\xi^2=\xi,\qquad \xi^{*2}=\xi^*.
\]
Ако ныне положимо: \(e=\overline E+E=\overline E^*+E^*\) и с тым:
\[
\mathfrak K_n=\overline A+A
=\mathfrak K_n\overline\xi+\mathfrak K_n\xi
=\overline{\mathfrak z}+\mathfrak z
=\mathfrak K_n\overline\xi^*+\mathfrak K_n\xi^*,
\]
то \(\overline A\) и \(\overline{\mathfrak z}\) сут такоже изоморфне; ако теды
\[
\overline E\mapsto\overline s=\overline E\,\overline s
\qquad\text{i}\qquad
\overline E^*\mapsto\overline t=\overline E^*\overline t
\]
сут обратне изоморфизмы, и ако положимо:
\[
\lambda=s+\overline s,\qquad \mu=t+\overline t,
\]
достава се:
\[
\lambda\mu=\mu\lambda=e,\qquad
\mathfrak l_i\lambda=\mathfrak q_i,\qquad
\lambda^{-1}\mathfrak q_i=\mathfrak q_i
\quad
(\lambda^{-1}\mathfrak q_i\subseteq\mathfrak q_i
=\lambda\cdot\lambda^{-1}\mathfrak q_i\subseteq\lambda^{-1}\mathfrak q_i)
\]
и зато
\[
\lambda^{-1}\mathfrak l_i\lambda=\mathfrak q_i,\qquad \text{čto bylo dokazati.}
\]

К разкладу (1) можно јест знову определити елементы
\[
r_{ik}=c_{ik}^{\prime(ik)}\gamma_{ik},\qquad \gamma_{ik}\in\{Z_i,Z_k\},
\]
кторе посредством \(e'_iE'\mapsto r_{ik}\) давајут операторне изоморфизмы меджу идеалами \(\mathfrak l'_i\mathfrak z'_1\) и \(\mathfrak l'_k\mathfrak z'_1\), и такоже задовольавајут услов коньугованости \(S(r_{ik})=r_{\nu\mu}\), ако \(S(i,k)=(\nu,\mu)\). Они посредством
\[
r_{ij}r_{jk}=\alpha_{ik}^{(j)}r_{ik}
\]
определьајут факторны систем \(\alpha\) од \(A_\Gamma\). Ныне јест важно, же за определьенье факторных системов од \(A_\Gamma\) можно једнако употребити либовольны разклад (2). Бо ако систем \(r_{ik}\), приналежечи к (1), заменимо с \(r'_{ik}=\lambda^{-1}r_{ik}\lambda\), то \(e_i^*\mapsto r'_{ik}\) даваје операторны изоморфизм од \(\mathfrak q_i\) на \(\mathfrak q_k\); \(r'_{ik}\) задовольавајут услов коньугованости, бо \(\lambda\) јест \(\mathfrak q\)-инвариантны, а \(r'_{ik}\) имајут тојже факторны систем како \(r_{ik}\).

Ако изберемо систем псеудоматрицных јединиц \(p_{ik}\) од \(\mathfrak K\), то
\[
r_{ik}=\tau^{-1}p_{ik}\tau\cdot E'
\]
јест \(r_{ik}\)-систем од \(A_\Gamma\) с једнакым факторным системом. Теды имајемо помочну лему:

\textbf{Теорем.} \emph{Факторне системы од \(A_\Gamma\), образоване с помочју либовольного разклада}
\[
A_\Gamma=\mathfrak q_1+\cdots+\mathfrak q_n
\]
\emph{од \(A_\Gamma\) на коньуговане леве идеалы, сут једнака факторным системам од \(\mathfrak K_r\).}

Следујучи теорем показује, же ако \(\alpha_{ik}^{(j)}\), \(\beta_{ik}^{(j)}\) сут факторне системы од \(\{\mathfrak K\}\), \(\{\mathfrak S\}\) при избраньу од \(Z\), то \(\varepsilon_{ik}^{(j)}=\alpha_{ik}^{(j)}\cdot\beta_{ik}^{(j)}\) јест факторны систем од \(\{\mathfrak K\times\mathfrak S\}\) при избраньу од \(Z\). Зато имаје смысл дефинирати \emph{продукт \(\{\alpha\}\times\{\beta\}\) двох клас асоциованых факторных системов} \(\{\alpha\}\), \(\{\beta\}\) како класу \(\{\alpha_{ik}^{(j)}\beta_{ik}^{(j)}\}\).

\textbf{Теорем 2.} \emph{Нека \(\{\mathfrak K\}\) при избраньу од \(Z\) имаје класу \(\{\alpha\}\) асоциованых факторных системов. Класы \(\{\alpha\}\) всих \(\{\mathfrak K\}\) с разпадным польем \(Z\) творят групу, ктора јест посредством \(\{\mathfrak K\}\mapsto\{\alpha\}\) изоморфно однесена на \(\mathscr K_Z\).}

\textbf{За доказ} треба показати двоје:

1. Ако \(\{\mathfrak K\}\) имаје факторны систем \(\alpha_{ik}^{(j)}\), а \(\{\mathfrak S\}\) \(\beta_{ik}^{(j)}\), то \(\varepsilon_{ik}^{(j)}=\alpha_{ik}^{(j)}\beta_{ik}^{(j)}\) јест факторны систем од \(\{\mathfrak K\times\mathfrak S\}\).

2. Ако \(\{\mathfrak K\}\) имаје факторны систем \(\alpha_{ik}^{(j)}=1\), то \(\mathfrak K=\mathsf P\).

1. Нека ступеньи иредуцибилных представјень од \(Z\) в \(\mathfrak K\), \(\mathfrak S\) буду \(r,s\)
\[
\mathfrak K_{r\Gamma}=\mathfrak l_1+\cdots+\mathfrak l_n,\qquad
\mathfrak S_{s\Gamma}=\mathfrak m_1+\cdots+\mathfrak m_n
\]
разклады на коньуговане идеалы, причем \(\mathfrak l_i\) и \(\mathfrak m_i\) треба да буду взајемно одповедајучи идеалы (изведене из одповедајучих компонентов од \(\mathfrak Z_r\), соответно од \(\mathfrak Z'_r\)). Тада јест
\[
(\mathfrak K_r\times\mathfrak S_s)_\Gamma
=\mathfrak l_1\mathfrak m_1+\cdots+\mathfrak l_i\mathfrak m_j+\cdots+\mathfrak l_n\mathfrak m_n
\]

разклад на \(n^2\) ненуловых, теды простых идеалов. Зато јест
\[
A_1=\mathfrak l_1\mathfrak m_1+\mathfrak l_2\mathfrak m_2+\cdots+\mathfrak l_n\mathfrak m_n
\]
идеал абсолутној должины \(n\). \(A_1\) јест инвариантны при всих \(S\), бо \(\mathfrak m_i\) трансформујут се тако како \(\mathfrak l_i\); зато постојí \(\mathfrak K_r\times\mathfrak S_s\)-идеал \(A\) с
\[
A_1=A_\Gamma
\]
(помочна лема, стр.~18).

Ако \(u\) јест ступень иредуцибилного представјеньа од \(\mathfrak Z\) в телу автоморфизмов \(\mathfrak M\) простых идеалов од \(\mathfrak K\times\mathfrak S\), то \(A\) разпада се на \(u\) простых идеалов (бо абсолутна должина јест \(n\)). Наша помочна лема даваје теды, же факторне системы од \(A\) сут такоже факторне системы од \(\mathfrak M\).

Але разклад од \(A_\Gamma\) на коньуговане идеалы, употребјены за дефиницију, јест:
\[
\tag{3}
A_\Gamma=\mathfrak l_1\mathfrak m_1+\mathfrak l_2\mathfrak m_2+\cdots+\mathfrak l_n\mathfrak m_n.
\]
Ако \(\alpha_{ik}^{(j)}\), \(\beta_{ik}^{(j)}\) приналежет к псеудоматрицным јединицам \(p_{ik}\), \(q_{ik}\) од \(\mathfrak K\), \(\mathfrak S\), то очивидно \(r_{ik}=p_{ik}q_{ik}\) јест систем псеудоматрицных јединиц од \(A_\Gamma\) взходно к разкладу (3); к тому приналежечи факторны систем \(\varepsilon_{ik}^{(j)}=\alpha_{ik}^{(j)}\beta_{ik}^{(j)}\) јест тада такоже једен из факторных системов од \(\mathfrak M\).

2. Нека \(\{\mathfrak K\}\) имаје факторны систем \(\alpha_{ik}^{(j)}=1\) с \(p_{ik}\) како приналежечими псеудоматрицными јединицами и разкладом
\[
\mathfrak K_{r\Gamma}=\mathfrak l_1+\cdots+\mathfrak l_n
\]
на коньуговане идеалы. Треба доказати \(\mathfrak K=\mathsf P\), или \(\mathfrak K_r=\mathsf P_r\), или постојанье \(n\)-членного левого идеала \(\mathfrak l\) од \(\mathfrak K_r\) -- бо просты левы идеал од \(\mathfrak K_r\) имаје ранг \(t^2r\), теды \(n=t^2rh\) с целым \(h\), или мора бити \(t=1\).

Идеал од \(\mathfrak K_{r\Gamma}\)
\[
\Omega=\mathfrak l_1(p_{11}+\cdots+p_{1n})
\]
не јест нула -- \(e_1\sum_i p_{1i}\ne0\) -- и јест просты, теды ранга \(n\). \(\Omega\) јест \(\mathfrak G\)-инвариантны,
\[
S\Omega
=\mathfrak l_i\sum_\nu S p_{1\nu}
=\mathfrak l_i\sum_\mu p_{i\mu}
=\mathfrak l_i\sum_\mu p_{i1}p_{1\mu}
\]
бо \(\alpha_{ik}^{(j)}=1\). Зато
\[
S\Omega=\mathfrak l_ip_{i1}\sum_\mu p_{1\mu}
=\mathfrak l_1\sum_\mu p_{1\mu}
=\Omega.
\]
Теды \(\Omega=\mathfrak l_\Gamma\), \emph{гдје \(\mathfrak l\) означаје \(\mathfrak K_r\)-идеал \(n\)-того ранга.}

\textbf{Теорем 3.} \emph{Всакы елемент \(\{\mathfrak K\}\) групы \(\mathscr K\) имаје конечны ред \(\lambda\), кторы јест делитель абсолутного числа компонентов \(t\) од \(\mathfrak K\).}

\emph{Првы доказ (Брауер).} Изберемо два разклады од \(\mathfrak K_{r\Gamma}\) на коньуговане леве идеалы
\[
\mathfrak K_{r\Gamma}=\mathfrak l_1+\cdots+\mathfrak l_n=\mathfrak z_1+\cdots+\mathfrak z_n,\qquad
\mathfrak l_i=\mathfrak K_{r\Gamma}e_i,\quad
\mathfrak z_i=\mathfrak K_{r\Gamma}e_i^*.
\]
Нека посредством \(a_i\mapsto a_ir_{ik}\), специално \(e_i\mapsto r_{ik}\), буде заданы операторны изоморфизм од \(\mathfrak l_i\) на \(\mathfrak z_k\). При ньеговом обрату \(e_k^*\) преходи в неко \(s_{ki}\) из \(\mathfrak l_i\), \(s_{ki}\cdot r_{ik}=e_k^*\). Ако изберемо \(r_{ik}\) коньуговаными
\[
S(r_{ik})=r_{\nu\mu}\quad\text{ako}\quad S(i,k)=(\nu,\mu),
\]
то \(s_{ik}\) такоже сут коньуговане:
\[
S(s_{ki})S(r_{ik})=S(s_{ki}r_{ik})=S(e_k^*)=e_\mu^*
=s_{\mu\nu}r_{\nu\mu}=s_{\mu\nu}S(r_{ik}),
\]
теды \(S(s_{ki})=s_{\mu\nu}\). Понеже посредством \(e_i\mapsto e_ir_{ij}s_{jk}\) јест заданы операторны изоморфизм од \(\mathfrak l_i\) на \(\mathfrak l_k\), мора бити
\[
r_{ij}s_{jk}=p_{ik}\xi_{ik}^{(j)},\qquad
\xi_{ik}^{(j)}\ne0\in\Gamma
\]

Понеже \(r\)-елементи, \(s\)-елементи и \(p\)-елементи сут коньуговане, \(\xi\) такоже сут коньуговане: Из \(S(i,j,k)=(\nu,\tau,\mu)\) следује
\[
S(\xi_{ik}^{(j)})=\xi_{\nu\mu}^{\tau}.
\]

Из
\[
r_{i\lambda}s_{\lambda j}r_{j\lambda}s_{\lambda k}
=r_{i\lambda}e_\lambda^*s_{\lambda k}
=r_{i\lambda}s_{\lambda k}
\]
следује
\[
p_{ij}\xi_{ij}^{(\lambda)}\cdot p_{jk}\xi_{jk}^{(\lambda)}
=p_{ik}\xi_{ik}^{(\lambda)}
\]
или, понеже \(p_{ij}p_{jk}=p_{ik}\alpha_{ik}^{(j)}\):
\[
\alpha_{ik}^{(j)}
=\frac{\xi_{ik}^{(\lambda)}}{\xi_{ij}^{(\lambda)}\xi_{jk}^{(\lambda)}}.
\]
Ако положимо \(\delta_{\nu\mu}=\prod_{\lambda}\xi_{\nu\mu}^{(\lambda)}\), то \(\delta_{\nu\mu}\) јест инвариантны при всих \(S\), кторе оставјајут \(\nu\) и \(\mu\) неподвижными, лежи теды в \(\{Z_\nu,Z_\mu\}\); надто \(\delta_{\nu\mu}\) сут коньуговане, бо \(\xi_{\nu\mu}^{(j)}\) сут коньуговане. Понеже ныне
\[
\alpha_{ik}^{(j)n}
=\frac{\prod_{\lambda}\xi_{ik}^{(\lambda)}}
       {\prod_{\lambda}\xi_{ij}^{(\lambda)}\prod_{\lambda}\xi_{jk}^{(\lambda)}}
=\frac{\delta_{ik}}{\delta_{ij}\delta_{jk}}
\]
из формулы за асоциоване факторне системы (стр.~32) следује, же \(\alpha_{ik}^{(j)n}\) јест асоциованы с системом \(\beta_{ik}^{(j)}=1\); зато, понеже ред од \(\{\mathfrak K\}\) јест равны реду од \(\{\alpha\}\), ред од \(\{\mathfrak K\}\) јест делитель од \(n\). Ако особливо изберемо \(\mathfrak Z\) како максимално комутативно подполье од \(\mathfrak K\), то јест \(n=t\), теды \(\lambda\) јест делитель од \(t\).

\emph{Другы доказ (Сцхур).} Тој доказ основан јест на представјеньу \(p_{ik}\) посредством регуларных матриц реда \(n\), \(P_{ik}\), в \(\Gamma\), кторе задовольавајут једнакы условы коньугованости како \(p_{ik}\). \(P_{ik}\) лежи в \(\{Z_i,Z_k\}\); ако \(p_{ik}\) приналежи к \(P_{ik}\), то \(S(p_{ik})\) приналежи к \(S(P_{ik})\). То представјенье јест хомоморфно; т.~ј. ако \(P_{ik}\) приналежи к \(p_{ik}\), а \(P_{kj}\) к \(p_{kj}\), то \(P_{ik}P_{kj}\) приналежи к \(p_{ik}p_{kj}\), или
\[
P_{ik}P_{kj}=\alpha_{ij}^{(k)}P_{ij}
\]
\((P_{ik_1}P_{k_2j}=0\) за \(k_1\neq k_2\) не може важити). Ако имајемо то представјенье од \(p_{ik}\), закльучујемо тако: Преход к детерминантам даваје:
\[
|P_{ik}|\cdot |P_{kj}|=\alpha_{ij}^{(k)n}\cdot |P_{ij}|,
\]
теды, понеже матрице \(P_{ik}\) сут регуларне, \(|P_{\nu\mu}|\ne0\),
\[
\alpha_{ij}^{(k)n}=\frac{\delta_{ik}\delta_{kj}}{\delta_{ik}},
\]
гдје јест положено \(|P_{\nu\mu}|=\delta_{\nu\mu}\). И потом далье како в Брауеровом доказу.

\emph{Представјенье од \(p_{ik}\) посредством матриц \(P_{ik}\):}

Идеалы \(\mathfrak l_i\) разклада
\[
\mathfrak K_{r\Gamma}=\mathfrak l_1+\cdots+\mathfrak l_n
\]
сут реципрочне модулы представјеньа од \(\mathfrak K_r\) в \(\Gamma\) (теорем 2 на стр.~32). Ако за основу изберемо установјену базу \(k_1,\ldots,k_n\) од \(\mathfrak K_r\) взходно к \(\mathfrak Z\), то јест
\[
\mathfrak K_{r\Gamma}=k_1\mathfrak Z_r+\cdots+k_n\mathfrak Z_r,
\]
теды
\[
\mathfrak l_i=\mathfrak K_{r\Gamma}e_i
=k_1\mathfrak Z_r e_i+\cdots+k_n\mathfrak Z_r e_i
=k_1e_i\Gamma+\cdots+k_ne_i\Gamma.
\]
Из тога следује
\[
\mathfrak l_k=\mathfrak l_i p_{ik}
=k_1p_{ik}\Gamma+\cdots+k_np_{ik}\Gamma.
\]
Две базы \((k_1e_i,\ldots,k_ne_i)\) и \((k_1p_{ik},\ldots,k_np_{ik})\) од \(\mathfrak l_k\) можно јест превести једну в другу множеньем регуларној матрицоју \(P_{ik}\):
\[
P_{ik}\cdot
\begin{pmatrix}
k_1e_i\\ \vdots\\ k_ne_i
\end{pmatrix}
=
\begin{pmatrix}
k_1p_{ik}\\ \vdots\\ k_np_{ik}
\end{pmatrix},
\qquad
\text{kratko}\quad
P_{ik}(k_\varrho e_i)=(k_\varrho p_{ik}).
\]
Из тога следује коньугованост од \(P_{ik}\):
\[
S(P_{ik})(k_\varrho S(e_k))=(k_\varrho S(p_{ik}))
\quad\text{ili}\quad
S(P_{ik})(k_\varrho e_\mu)=(k_\varrho p_{\nu\mu})
=P_{\nu\mu}(k_\varrho e_\mu),
\]
т.~ј.
\[
S(P_{ik})=P_{\nu\mu}\quad\text{ako}\quad S(i,k)=(\nu,\mu).
\]
Из тога знову следује, же \(P_{ik}\) лежи в \(\{Z_i,Z_k\}\) -- јменно в инвариантном польу тых \(S\), кторе оставјајут \(i\) и \(k\) неподвижными. Хомоморфност такоже јест легко видети:
\[
\begin{split}
P_{ik}P_{kj}(k_\varrho e_j)
&=P_{ik}(k_\varrho p_{kj})
=(k_\varrho p_{ik}p_{kj})\\
&=\alpha_{ij}^{(k)}(k_\varrho p_{ij})
=\alpha_{ij}P_{ij}(k_\varrho e_j),
\quad\text{tedy}\\
&\cdot\, P_{ik}P_{kj}=\alpha_{ij}^{(k)}P_{ij},
\quad\text{čto bylo dokazati.}
\end{split}
\]

Теорем 3 можно јест посилити:

\textbf{Теорем 4.} \emph{Ако \(p\) јест просты фактор абсолутного числа компонентов \(t\) од \(\mathfrak K\), то \(p\) входи такоже в ред \(\lambda\) од \(\{\mathfrak K\}\). (Брауер)}

\emph{Доказ.} Нека \(p^\sigma\) буде највышша потенција од \(p\), ктора се налази в ступеньу \(n\) разпадного польа \(Z\), \(\mathfrak H\) подгрупа реда \(p^\sigma\) од \(\mathfrak G\) (она постојí по \emph{теорему од Сылова}, Спеисер 1. изд. § 19, стр.~42, теорем 43), а \(T\) инвариантно полье од \(\mathfrak H\).

Тада \([\Gamma:T]=p^\sigma\) и \([T:\mathsf P]=n/p^\sigma\) не јест деливы посредством \(p\), а тым выше не посредством \(t\). \(T\) теды не јест разпадно полье; тело автоморфизмов \(\mathfrak S\) од \(\mathfrak K_T=\mathfrak S_m\) јест различно од \(\mathsf P\). Але понеже \(\mathfrak S\) имаје разпадно полье \(\Gamma\), чији ступень над центром \(T\) од \(\mathfrak S\) јест равны \(p^\sigma\), абсолутно число компонентов од \(\mathfrak S\) јест потенција \(p^\beta>1\), зато по теорему 3 ред од \(\{\mathfrak S\}\) јест потенција \(p^\alpha>1\) од \(p\). Понеже из
\[
\{\mathfrak K\}^{\lambda}=\{\mathsf P\},
\qquad
\{\mathfrak K_T\}^{\lambda}=\{T\}
\]
следује, то \(\lambda\equiv0\pmod {p^\alpha}\), что было доказати. (Брауер дал примеры, же не мора бити \(\lambda=t\)\srcfn{1}{Брауер: Унтерсуцхунген üбер дие аритхметисцхен Еигенсцхафтен вон Группен линеарер Субститутионен II. Матх. Зеитсцхр. 31, стр.~733 § 5, (1930).}).

\textbf{Теорем 5.} \emph{Нека \(t\), разложено на просте факторы, буде \(t=\prod_i p_i^{\varrho_i}\), \(\varrho_i\ge1\), \(p_i\ne p_j\), ако \(i\ne j\). Тада јест \(\mathfrak K=\mathfrak K^{(1)}\times\mathfrak K^{(2)}\times\cdots\), гдје \(\mathfrak K^{(i)}\) јест тело с абсолутным числом компонентов \(p_i^{\varrho_i}\).}

\emph{Доказ.} По теоремах 3 и 4 ред \(\lambda\) од \(\{\mathfrak K\}\) јест
\[
\lambda=\prod_i p_i^{\alpha_i},\qquad 1\le\alpha_i\le\varrho_i.
\]
Посредством \(\{\mathfrak K\}\) створјена циклична подгрупа \(\mathfrak Z\) од \(\mathscr K\) јест премы продукт подгруп \(\mathfrak Z_i\) редов \(p_i^{\alpha_i}\), теды
\[
\{\mathfrak K\}=\prod_i\{\mathfrak K^{(i)}\},
\]

из чего достава се, же всако разпадно полье \(Z\) од \(\mathfrak K\) јест такоже разпадно полье од \(\mathfrak K^{(i)}\), теды абсолутно число компонентов \(p_i^{\sigma_i}\) од \(\mathfrak K^{(i)}\) дели \(t\), \(\sigma_i\le\varrho_i\). Нека \(\prod_i \mathfrak K^{(i)}=\mathfrak K_m\), теды
\[
\prod_i p_i^{\varrho_i}=\prod_i p_i^{\sigma_i}\cdot m,
\qquad
\sigma_i\ge\varrho_i,
\]
зато
\[
\sigma_i=\varrho_i,\qquad m=1,\qquad
\mathfrak K=\prod_i \mathfrak K^{(i)},
\quad\text{čto bylo dokazati.}
\]
% END BOOK_S24

% BEGIN BOOK_S25
\subsection*{§ 25. Нормално представјенье од \(\mathfrak K_r\) с Галоисовым максималным комутативным подпольем\protect\footnotemark}
\label{normaldarstellung-von-r_r-mit-galoisschem-maximalem-kommutativem-teilkoerper-1}
\footnotetext{Простејше основанье од § 27 далье.}

Нека максимално комутативно подполье \(\mathfrak Z\) од \(\mathfrak K_r\) буде Галоисово. \(n\) автоморфизмов \(H_i\) од \(\mathfrak Z/\mathsf P\) сут по теорему 7 на стр.~25 внутрне автоморфизмы од \(\mathfrak K_r\), т.~ј. постојут елементи \(u_i\) с
\[
H_i z=u_i^{-1}zu_i\quad\text{za vse }z\text{ iz }\mathfrak Z.
\]
Всакы \(u_i\) јест определьен до ненулового фактора \(y\) из \(\mathfrak Z\). Бо ако \(u^{-1}zu=u'^{-1}zu'\) за все \(z\) из \(\mathfrak Z\), то следује \((u^{-1}u')^{-1}z(u^{-1}u')=z\) за все \(z\in\mathfrak Z\). Приложенье од \(u^{-1}u'\) к \(\mathfrak Z\) даваје теды комутативно подполье од \(\mathfrak K_r\), кторо заради максималности од \(\mathfrak Z\) совпада с \(\mathfrak Z\). Теды \(u^{-1}u'\) лежи в \(\mathfrak Z\).

Доказујемо следујучи теорем:
\[
\mathfrak K_r=u_1\mathfrak Z+\cdots+u_n\mathfrak Z.
\]
За то очивидно јест потребна само линеарна независност од \(u_i\) взходно к \(\mathfrak Z\). Доказ се изводи непосредним заданьем од \(u_i\). Нека буде
\[
\mathfrak Z_Z=e_1Z+\cdots+e_nZ,\qquad
\mathfrak K_{rZ}=\mathfrak l_1+\cdots+\mathfrak l_n,\quad
\mathfrak l_i=\mathfrak K_{rZ}e_i
\]
и нека \(p_{ik}\) буде систем псеудоматрицных јединиц взходно к тому разкладу. На стр.~8 означили јесмо с \(\Theta_i\) изоморфизм од \(\mathfrak Z\) на \(Z\), даны посредством модула представјеньа \(e_iZ\), тако же
\[
S_i=\Theta_i\Theta_1^{-1}
\]
сут \(n\) автоморфизмов од \(Z\) (але само колико се они применьујут на само \(Z\), не выше за тут всечасно предпокладано дејство од \(S_i\) на цело \(\mathfrak K_{rZ}\)!). Заради простоты хочемо вместо \(\Theta_i\) в \(S_i=\Theta_i\Theta_1^{-1}\) писати \(\Theta_{S_i}\), а тада такоже при \(p_{ik}\) заменити индексы \(i,k\) приналежечими автоморфизмами \(S\)
\[
p_{ik}=p_{S_iS_k}.
\]

\(n\) автоморфизмов од \(\mathfrak Z\) сут (\(E\) јест идентична субституција)
\[
H=\Theta_E^{-1}\Theta.
\]
Нумерованье од \(H\) ныне избирајемо тако:
\[
H_S=\Theta_E^{-1}\Theta_{S^{-1}}.
\]

При тых условјах можно јест положити

\[
u_T=\sum_S p_{S,ST}=\sum_S S(p_{E,T})
\]

Бо

1. \(u_T^{-1}\) постојí:
\[
u_T^{-1}=\sum_S p_{ST,S}
\frac{1}{\alpha_{S,S}^{ST}\overline{\gamma}_{S,ST}\overline{\gamma}_{ST,S}}
\]

(Значенье од \(\bar{\gamma}\) види на стр.~32).

2. \(u_T\) јест елемент од \(\mathfrak K_r\), бо јест \(\mathfrak G\)-инвариантны:
\[
R(u_T)=\sum_S RS(p_{E,T})=\sum_{S'} S'(p_{E,T})=u_T.
\]

3. \(u_T\) стварјаје \(H_T\), т.~ј. ако \(z\in\mathfrak Z\), то \(u_T^{-1}zu_T=H_T(z)\), или
\[
z u_T=u_T H_T(z).
\]

Јест бо

\[
z=\sum_S\Theta_S(z)e_S,\qquad
u_T=\sum_S S(e_Ep_{E,T})=\sum_S e_S S(p_{E,T})
\]

теды

\[
zu_T=\sum_{S,S'}\Theta_S(z)e_Se_{S'}S(p_{E,T})
=\sum_S\Theta_S(z)S(p_{E,T}).
\]

С другој страны имајемо

\[
u_T=\sum_S S(p_{E,T}e_T)=\sum_S S(p_{E,T})e_{ST}
\]

и важи следујуче једначенье за дејство на елементы од \(\mathfrak Z\)

\[
\Theta_S H_T=\Theta_S\Theta_E^{-1}\Theta_{T^{-1}}
=S\Theta_{T^{-1}}=ST^{-1}\Theta_E=\Theta_{ST^{-1}}
\]

теды

\[
H_T(z)=\sum_S e_S\Theta_S H_T(z)=\sum_S e_S\Theta_{ST^{-1}}(z)
\]

т.~ј.

\[
\begin{aligned}
u_TH_T(z)
&=\sum_{S,S'} S(p_{E,T})e_{ST}\cdot e_{S'}\Theta_{S'T^{-1}}(z)\\
&=\sum_S S(p_{E,T})\Theta_S(z)
\end{aligned}
\]

теды

\[
z\cdot u_T=u_T\cdot H_T(z).
\]

4. \(u_S\) сут линеарно независне взходно к \(\mathfrak Z\).

Из
\[
\sum_T z_Tu_T=0
\]
заради \(z=\sum_S\Theta_S(z)e_S\) следује

\[
\sum_{T,S,L}\Theta_S(z_T)e_Sp_{L,LT}
=\sum_{S,T}\Theta_S(z_T)p_{S,ST}\quad\text{tedy}\quad
\Theta_S(z_T)=0,
\]

теды: \(z_T=0\).

Од ныне вместо \(H_S(z)\) пишемо \(z^S\), ,,симболична потенција``. За \(u_S\) морајут важити релације следујучего вида

\[
u_Ru_T=u_{RT}a_{R,T},\qquad a_{R,T}\ne0\quad\text{iz }\mathfrak Z.
\]

Понеже

\[
\begin{aligned}
u_Ru_T
&=\sum_{S,S'} S(p_{E,R})S'(p_{E,T})
=\sum_S p_{S,SR}p_{SR,SRT}\\
&=\sum_S p_{S,SRT}\alpha_{S,SRT}^{(SR)}
=\sum_S p_{S,SRT}\alpha_{S,SRT}^{(SR)}e_{SRT}
\end{aligned}
\]

и

\[a_{R,\,T} = \sum\limits_{S} e_{SRT}\, \Theta_{SRT} \left(a_{R,\,T}\right),\]

теды

\[u_{RT} a_{R,T} = \sum\limits_{S} p_{S,SRT} \Theta_{SRT} (a_{R,T})\]

следује

\[\Theta_{SRT}(a_{R,\,T}) = \alpha_{S,SRT}^{(SR)}, ~~ a_{R,\,T} = \sum\limits_{S} e_{SRT} \alpha_{S,SRT}^{(SR)} = \sum\limits_{L} e_{L} \alpha_{LT^{-1}R^{-1},L}^{(LT^{-1})}.\]

Систем од \(a_{R,\,T}\) хочемо называти малым факторным системом од \(\mathfrak K_r\); велики факторны систем \(\alpha_{BC}^{(A)}\) и мали факторны систем \(a_{R,\,T}\) определьајут се взајемно. К продукту \(\varepsilon_{BC}^{(A)} = \alpha_{BC}^{(A)} \beta_{BC}^{(A)}\) двох великых факторных системов приналежи продукт \(f_{R,\,T} = a_{R,\,T}b_{R,\,T}\) малых факторных системов \(a_{R,\,T}\), \(b_{R,\,T}\), приряđеных к \(\alpha_{BC}^{(A)}\), \(\beta_{BC}^{(A)}\); то непосреднье следује из горньих формул. К великому јединчному факторному систему \(\alpha_{BC}^{(A)} = 1\) приналежи мали јединчны факторны систем \(a_{R,\,T} = 1\), и обратно.

Можно јест заменити \(u_S\) и \(u_S'=u_Sr_S\), гдје \(r_S\neq 0\) јест избрано либовольно из \(\mathfrak Z\). К \(u_S'\) приналежи мали факторны систем \(a_{R,T}'\)

\[
\text{(1)}\qquad u_R'u_T'=u_{RT}'a_{R,T}'
\]

\(a_{R,T}^{\prime}\) называјемо асоциованым с \(a_{R,T}\). Јест

\[u_R r_R u_T r_T = u_{RT} r_{RT} a'_{R,T} = u_R u_T r_R^T r_T = u_{RT} a_{R,T} r_R^T r_T\]

из чего

\[
\text{(2)}\qquad a'_{R,T}=a_{R,T}\frac{r_R^T r_T}{r_{RT}}
\]

следује.

К \(u_S' = u_S r_S\) приналежи следујучи систем псеудоматрицных јединиц

\[p_{S,ST}^{\prime}=p_{S,ST}r_{T}\]

мали факторны систем \(a'_{R,\,T}\) приналежи теды к великому факторному систему \(\alpha^{(A)\prime}_{BC}\), кторы јест асоциованы с факторным системом \(\alpha^{(A)}_{BC}\), приналежечим к \(a_{R,\,T}\); асоциоване мале факторне системы одповедајут асоциованым великым факторным системам и обратно. Приряđенье клас \(\{\alpha\}\) асоциованых великых факторных системов к приналежечим класам асоциованых малых факторных системов, \(\{a\}\), јест изоморфизм груп.

Ныне преходимо к представјеньу од \(u_S\) посредством матриц реда н в \(\mathfrak{Z}\), кторе се сучствено разликује од досега разматраных представјень.

Ако \(\mathfrak Z\)-базу

\[
\text{(I)}\qquad k_{S_1}, \ldots, k_{S_n}
\]

од \(\mathfrak K_r\) множимо с права посредством \(u_S\), настава

\[
\text{(II)}\qquad k_{S_1}u_S, \ldots, k_{S_n}u_S.
\]

(Регуларну) матрицу \(A_S\), ктора преводи (И) в (II),

\[
\text{(III)}\qquad (k_{S_1}u_S, \ldots, k_{S_n}u_S) = (k_{S_1}, \ldots, k_{S_n})A_S
\]

называмо матрицоју, ктора представльа \(u_S\). Представјенье се односи на все елементы од \(\mathfrak K_{r}\), кторе како \(u_S\) стварјајут автоморфизм од \(\mathfrak Z/P\), т.~ј. на все елементы \(u_Sr_S\), гдје \(r_S \neq 0\) јест избрано либовольно из \(\mathfrak Z\). Вси ти елементи творят групу \(\mathfrak G^*\), ктора како нормалну подгрупу садржује мултипликативну групу \(\mathfrak Z^*\) од \(\mathfrak Z\); факторна група \(\mathfrak G^*/\mathfrak Z^*\) јест изоморфна с Галоисовој групоју \(\mathfrak G\) од \(\mathfrak Z/P\) тако, же всакы елемент једного побочного класа од \(\mathfrak G^*\) взходно к \(\mathfrak Z^*\) стварјаје автоморфизм од \(\mathfrak Z\), приряđены тому класу. Та група \(\mathfrak G^*\) представльа се управо указаным начином

\(u_S \mapsto A_S\)

Але то представјенье не јест хомоморфно, ни реципрочно хомоморфно; из

\(u_S \mapsto A_S\), \(u_T \mapsto A_T\)

не следује \(u_Su_T \mapsto A_SA_T\), але заради

\[
\begin{aligned}
(k_{S_1}, \ldots, k_{S_n})u_Su_T
&=(k_{S_1}, \ldots, k_{S_n})A_Su_T
=(k_{S_1}, \ldots, k_{S_n})u_TA_S^T\\
&=(k_{S_1}, \ldots, k_{S_n})A_TA_S^T:\ u_Su_T\mapsto A_TA_S^T.
\end{aligned}
\]

То представјенье од \(\mathfrak{G}^*\) называмо прекриженым представјеньем. Ако \(h_{S_1}, \ldots, h_{S_n}\) јест друга база од \(\mathfrak{K}_r\) взходно к \(\mathfrak{Z}\), и

\[
\text{(IV)}\qquad (h_{S_1}, \ldots, h_{S_n}) = (k_{S_1}, \ldots, k_{S_n})M
\]

то из (III) и (IV) достава се:

\[(h_{S_1},\ldots,h_{S_n})\,u_S=(h_{S_1},\ldots,h_{S_n})\,M^{-1}A_S\,M^S,\]

т.~ј.: Ако \(u_S\) јест при \(k\)-базе представльено посредством \(A_S\), то при \(h\)-базе јест представльено посредством \(M^{-1}A_SM^S\), ако \(M\) јест матрица, ктора трансформује \(k\)-базу в \(h\)-базу.

Формула (III), ктора дефинујеме прекрижено представјенье, показује, же прекрижено представјенье стоји в подобном одношеньу к колцу \(\mathfrak K_r\), како обычно реципрочно представјенье (але записано с права вместо с лева) к својему модулу представјеньа. Разлика настава, бо вместо релације комутативности

\[(m^*\tau^*)c = (m^*c)\tau^*\]
(види стр.~5)

при реципрочном модулу представјеньа тут

\[
(kz)u_S=(ku_S)z^S\qquad k\text{ v }\mathfrak K_r,\quad z\text{ v }\mathfrak Z
\]

наступаје.

Прекрижено представјенье од \(\mathfrak{G}^*\) можно јест такоже образовати посредством базы некого правого идеала \(\mathfrak r\) од \(\mathfrak{K}_r\). Всакы правы идеал од \(\mathfrak{K}_r\) стварјаје на тој начин класу еквивалентных представјень; при том представјеньа \(u_S \mapsto A_S, u_S \mapsto B_S\) называјут се еквивалентными, ако постојí матрица \(M\) с \(B_S = M^{-1}A_SM^S\).

Вместо ограниченьа на представјеньа од \(\mathfrak G^*\) в \(\mathfrak Z\), створјене правыми идеалами од \(\mathfrak K_r\), ныне можно јест обче дефинирати прекрижено представјенье \(m\)-того ступеньа од \(\mathfrak G^*\) в \(\mathfrak Z\) тако:

\begin{enumerate}
\item Всакому \(u_S\) одповедаје матрица \(A_S\) реда \(m\) над \(\mathfrak Z\).
\item Из \(u_S \mapsto A_S\), \(u_T \mapsto A_T\) следује \(u_Su_T \mapsto A_TA_S^T\).
\end{enumerate}

Изберемо правы \(\mathfrak Z\)-модул ранга \(m\)

\[\mathfrak{M}=x_1\mathfrak{Z}+\cdots+x_m\mathfrak{Z}\]

и установјеньем

\[(x_1u_S,\ldots,x_mu_S)=(x_1,\ldots,x_m)A_S\]

чинимо го правым \(\mathfrak G^*\)-модулом с релацијоју комутативности

\[(mz) u_S = (mu_S) z^S\]

целком аналогно конструкцији модула представјеньа к заданому представјеньу на стр.~5. Ныне показујемо, како на стр.~20, же \(\mathfrak{M}\) поставаје конечным правым \(\mathfrak K_r\)-модулом посредством установјеньа

\[m \cdot z u_S = mz \cdot u_S = mu_S \cdot z^S\]

обче

\[m \cdot \sum z_S u_S = \sum m z_S \cdot u_S = \sum m u_S \cdot z_S^S.\]

При том никако не јест потребно рачунанье аналогно рачунаньу на стр.~20, бо при нашем установјеньу употребили јесмо установјену базу.

Але по M.Z. § 18 всакы конечны просты \(\mathfrak K_r\)-модул \(\mathfrak{M}\) јест операторно изоморфны с простым правым идеалом од \(\mathfrak K_r\); из тога аналогно теорему 2 на стр.~22 закльучујемо, же постојí само \emph{једна} класа иредуцибилных прекриженых представјень од \(\mathfrak{G}^*\) в \(\mathfrak Z\) --- јменно она, ктора приналежи к простым правым идеалам.

Ступень тога иредуцибилного представјеньа јест \(t\), гдје \(t\) јест индекс (абсолутно число експонента) од \(\mathfrak K\). То се достава простым считаньем рангов. Бо ступень јест равны рангу \((\mathfrak r:\mathfrak Z)\) простого правого идеала над \(\mathfrak Z\). Ныне на стр.~22 важи:

\[
(\mathfrak Z:P)=n=rt;\qquad (\mathfrak K_r:P)=n^2;\qquad (\mathfrak r:P)=n\cdot t
\]

последнье заради \(\mathfrak K_r = \mathfrak{r}_1 + \cdots + \mathfrak{r}_r\).

Из
\[(\mathfrak{r}: P) = (\mathfrak{r}: \mathfrak Z) \cdot (\mathfrak Z: P)\]
или \(n \cdot t = (\mathfrak{r}: \mathfrak Z) \cdot n\) следује теды \((\mathfrak{r}: \mathfrak Z) = t\).
% END BOOK_S25

% BEGIN BOOK_S26
\subsection*{§ 26. Множенье прекриженых представјень}\label{multiplikation-der-verschraenkten-darstellungen}

Досели јесмо всегда разматрали прекрижено представјенье како представјенье од \(\mathfrak{G}^*\). Але јего можно јест разматрати такоже како представјенье Галоисовој групы \(\mathfrak{G}\) од \(\mathfrak Z/P\), и то тако, же всакому елементу \(S\) од \(\mathfrak{G}\) приряđајемо матрицу \(A_S\), ктора представльа елемент \(u_S\) из \(\mathfrak K_r\), стварјајучи автоморфизм \(S\) посредством трансформације -- \(u_S^{-1}zu_S=z^S\) --.

За означенье такого представјеньа \(S\mapsto A_S\)\srcfn{1}{Тако представјенье (независно од хиперкомплексных системов)
\[
S\mapsto A_S;\quad T\mapsto A_T;\quad ST\mapsto A_TA_S^T=A_{ST}\cdot a_{S,T}^{-1}
\]
(или транспоновано представјенье) было изходишчем Спеисера (Матх. Зсцхр. 5) при определьеньу факторных системов \(a_{S,T}\). Те факторне системы показывајут се теды идентичными с тут из другој основы уведеными.} треба сказати, в ктором \(\mathfrak K_r\) оно јест взято, кторе \(u_S\) и кторе \(k_{S_1},\ldots,k_{S_n}\) од \(\mathfrak K\) над \(\mathfrak Z\) сут употребјене. Замена базы даваје еквивалентно представјенье, а замена од \(u\) даваје асоциоване факторне системы \(a_{S,\,T}\); сучствене сут теды \(\mathfrak K_r\), в кторых представјенье јест взято, т.~ј. класа \(\{a_{S,\,T}\}\) асоциованых факторных системов.

Нека ныне будут дане два така представјеньа од \(\mathfrak G\).

\[
\begin{array}{lll}
S\mapsto u_S\mapsto A_S & \text{naležeči mali faktorny sistem} & a_{R,T},\\
S\mapsto v_S\mapsto B_S & \text{naležeči mali faktorny sistem} & b_{R,T}.
\end{array}
\]

Под Кронецкеровым продуктом \(A\times B\) двох матриц \(A=(a_{ik})\), \(B=(b_{\mu\nu})\) разумејемо матрицу \(A\times B=(a_{ik}b_{\mu\nu})\).

\noindent\textbf{Тврдженье.} \emph{Такоже \(S\mapsto A_S\times B_S\) јест прекрижено представјенье од \(\mathfrak G\), и оно приналежи к продукту \(\mathfrak K_r\times\mathfrak S_s\), ако \(A_S\) произходи из \(\mathfrak K_r\), а \(B_S\) из \(\mathfrak S_s\).}

\noindent\emph{Доказ.} Представјенье \(S\mapsto A_S\times B_S\) јест прекрижено представјенье. Нека буде
\[
S\mapsto A_S\times B_S,\qquad T\mapsto A_T\times B_T
\]
и
\[
ST\mapsto A_{ST}\times B_{ST}
=(A_TA_S^T\times B_TB_S^T)a_{S,T}^{-1}b_{S,T}^{-1},
\]
теды треба доказати:
\[
(A_T\times B_T)(A_S\times B_S)^T a_{S,T}^{-1}b_{S,T}^{-1}
=(A_TA_S^T\times B_TB_S^T)a_{S,T}^{-1}b_{S,T}^{-1}.
\]
То але следује из следујучего разматраньа:

Ако матрице \(A\) и \(B\) представимо посредством матричных јединиц \(c_{ik}\) и \(d_{\mu\nu}\) (кторе не имајут ничто обче с матричными јединицами од \(\mathfrak K_r\) или \(\mathfrak S_s\))
\[
A=\sum c_{ik}\alpha_{ik};\qquad
B=\sum d_{\mu\nu}\beta_{\mu\nu},\qquad\text{tada}
\]
\[
A\times B=B\times A
=\sum c_{ik}d_{\mu\nu}\cdot\alpha_{ik}\cdot\beta_{\mu\nu},
\]
т.~ј. оно јест равно матрици в нових матричных јединицах
\(c_{ik}d_{\mu\nu}=d_{\mu\nu}c_{ik}\). Множенье \(A\)-матриц и \(B\)-матриц теды једнозначно одповедаје множеньу елементов в премом продукту двох матричных колц:
\[
\sum c_{ik}\mathfrak Z\times\sum d_{\mu\nu}\mathfrak Z.
\]
Зато всака \(A\)-матрица јест комутативна с всакој \(B\)-матрицој (але \(A\)-матрице ни \(B\)-матрице не сут меджу собоју комутативне); и та комутативност всакој \(A\)-матрице с всакој \(B\)-матрицоју даваје управо --- по множеньу с \(a_{S,T}^{-1}b_{S,T}^{-1}\) --- релацију, ктору треба доказати.

\section*{Глава VI. Теорија прекриженых продуктов}\label{kapitel-vi.-theorie-der-verschraenkten-produkte}
% END BOOK_S26

% BEGIN BOOK_S27
\subsection*{§ 27. Представјенье од \(\mathfrak K_r\) како прекриженых продуктов}

Теорија прекриженых продуктов, развијена в § 24 из факторных системов, ныне буде знову независно основана. Значенье тых прекриженых продуктов јест в том, же Галоисово полье и јего групы сут \emph{заједно} описане в \(\mathfrak K_r\).

\noindent\textbf{Теорем 1.} \emph{Ако \(\mathfrak Z\) јест Галоисово разпадно полье од \(\mathfrak K_r\), теды максимално комутативно подполье, и ако \(\mathfrak G\) јест јего Галоисова група, то \(\mathfrak K_r\) јест равно прекриженому продукту \(\mathfrak Z\cdot\mathfrak G\). Обратно, всакы прекрижены продукт \(\mathfrak Z\cdot\mathfrak G\) стварјаје \(\mathfrak K_r\) такого вида.}

\noindent\textbf{Доказ и дефиниције. Дефиниција 1.} Нека \(\mathfrak G=\{E,S,\ldots,T\}\) буде Галоисова група од \(\mathfrak Z\); равност од \(\mathfrak K_r\) прекриженому продукту значи:
\[
\mathfrak K_r=\mathfrak Zu_E+\cdots+\mathfrak Zu_T
\]
с \(zu_S=u_Sz^S\) и \(u_Su_T=a_{S,T}u_{ST}\) (\(a_{S,T}\) в \(\mathfrak Z\) сут мале факторне системы).

\noindent\textbf{Дефиниција 2.} Стварјајуча група \(\mathfrak G^*\) јест определьена како множство регуларных елементов \(g\) из \(\mathfrak K_r\), кторе посредством трансформације преводет \(\mathfrak Z\) како целост в себе. \(g^{-1}\mathfrak Zg=\mathfrak Z\), теды \(g^{-1}zg=z^S\), с \(S\) в \(\mathfrak G\). Теды \(\mathfrak G\) јест групово-хомоморфны образ од \(\mathfrak G^*\), бо всакы автоморфизм од \(\mathfrak Z\) јест створјен некым \(g\), \(\mathfrak G^*\to\mathfrak G\), и далье
\[
\mathfrak G\cong\mathfrak G^*/\mathfrak T^*
\quad\text{s}\quad
\mathfrak T^*\supseteq\mathfrak Z^*,
\]
(же \(\mathfrak T^*=\mathfrak Z^*\), буде доказано поздеје --- још не было доказано, же \(\mathfrak Z\) было такоже максимално комутативно подколцо. То можно јест доказати и независно).

\noindent\textbf{1.} \emph{Доказ представјеньа како прекриженого продукта.} Нека \(u_E,u_S,\ldots,u_T\) будут елементи из \(\mathfrak K_r\), кторе стварјајут автоморфизмы \(E,\ldots,T\), теды \(zu_S=u_Sz^S\), особливо \(u_E=e\) заради 2. Треба показати, же сума
\[
\mathfrak M=\{\mathfrak Zu_E,\ldots,\mathfrak Zu_T\}
\]
јест према и равна \(\mathfrak Zu_E+\cdots+\mathfrak Zu_T\); тада она јест равна \(\mathfrak K_r\) заради једнакости рангов. То але следује тако (Сцхварз):

\(\mathfrak M\) јест левы и правы \(\mathfrak Z\)-модул, теды модул представјеньа од \(\mathfrak Z\) в \(\mathfrak Z\), и посредством \(\mathfrak M\) стварјајут се \(n\) различне представјеньа; теды --- понеже \(\mathfrak Z\) јест комутативно полье првого вида --- \(\mathfrak M\) имаје ранг најманье \(n\); иначе доказ тече управо тако.

\noindent\textbf{2.} \(\mathfrak G\cong\mathfrak G^*/\mathfrak Z^*\), т.~ј. \(\mathfrak T^*=\mathfrak Z^*\); само елементи из \(\mathfrak Z\) индуцирајут идентитет на \(\mathfrak Z\) посредством трансформације.

Нека именно буде \(tz=zt\) за всако \(z\) из \(\mathfrak Z\); положимо
\[
t=z_0+z_1u_S+\cdots+z_{n-1}u_T,
\]
теды
\[
zt=zz_0+\cdots+zz_{n-1}u_T
\]
и
\[
tz=z_0z+\cdots+z_{n-1}u_Tz
=z_0z+z^{S^{-1}}z_1u_S+\cdots+z^{T^{-1}}z_{n-1}u_T,
\]
теды заради премој сумы:
\[
(z-z^{S^{-1}})z_1u_S=0,\ldots,
(z-z^{T^{-1}})z_{n-1}u_T=0
\]
и из тога, понеже елементи \(u_S,\ldots\) сут регуларне:
\[
(z-z^{S^{-1}})z_1=0,\ldots,
(z-z^{T^{-1}})z_{n-1}=0,
\qquad\text{za vsako }z.
\]
(или понеже ранг јест равны \(n\)). Але по дефиницији за \(S,\ldots\), и тако далье (все неравне \(E\)), постојí најманье једно \(z\), тако же \(z-z^{S^{-1}}\neq0\); и неко \(\bar z\), тако же \(\bar z-\bar z^{T^{-1}}\neq0\); зато достава се \(z_i=0\), \(i=1,\ldots,(n-1)\), и теды \(t=z_0\) в \(\mathfrak Z^*\), что было доказати.

\noindent\textbf{3.} Елементи \(u\) задовольајут законы множеньа: \(u_Su_T=a_{S,T}u_{ST}\), при чем \(a_{S,T}\) задовольајут прекрижены асоциативны закон
\[
a_{R,S}\cdot a_{RS,T}=a_{S,T}^{R^{-1}}\cdot a_{R,ST}
\]
(\(a_{S,T}\) сут мали факторны систем).

Бо по 2. достава се:
\[
\mathfrak Z^*u_S\cdot\mathfrak Z^*u_T=\mathfrak Z^*u_{ST};
\]
теды
\[
u_Su_T=a_{S,T}u_{ST};
\]
с \(a_{S,T}\) из \(\mathfrak Z^*\), теды из \(\mathfrak Z\) и \(\neq0\). Заради \(u_R\cdot u_Su_T=u_Ru_S\cdot u_T\) далье достава се
\[
\begin{aligned}
u_Ra_{S,T}u_{ST}
&=a_{S,T}^{R^{-1}}\cdot u_Ru_{ST}
 =a_{S,T}^{R^{-1}}a_{R,ST}u_{RST}\\
&=a_{R,S}u_{RS}\cdot u_T
 =a_{R,S}a_{RS,T}u_{RST},
\end{aligned}
\]
что сут подане законы множеньа.

\noindent\textbf{4. Обрат:} прекрижены продукт \(=\mathfrak K_r\). Према сума:
\[
\mathfrak M=\mathfrak Zu_E+\cdots+\mathfrak Zu_T
\]
буде определьена како колцо посредством релациј: \(zu_S=u_Sz^S\), \(u_Su_T=a_{S,T}u_{ST}\), (кде \(a_{S,T}\) задовольајут прекрижене асоциативне релације; теды \(u_Ru_Su_T\) јест једнозначно определьено):
\[
\begin{array}{ll}
1.&z_1\cdot z_2u_S=z_1z_2\cdot u_S;\\[2pt]
2.&zu_S\cdot u_T=z\cdot u_Su_T
\end{array}
\]
тым \(\mathfrak M\) поставаје асоциативным колцом.

\noindent\textbf{Обратны теорем 2.} \emph{Прекрижены продукт \(\mathfrak Z\cdot\mathfrak G\) јест двустранно просты, с \(P\) како центром (по приступах R. Брауера).} Доказ основываје се на следујучих помочных лемах:

\noindent\textbf{Помочна лема 1.} \emph{Всакы елемент, кторы јест по елементах комутативны с \(\mathfrak Z\), приналежи к \(\mathfrak Z\).}

\noindent\emph{Доказ.} Нека \(w=\sum_S c_{\lambda_S}u_S\), (\(c_{\lambda_S}\) в \(\mathfrak Z\)) и нека буде предположено: \(wz=zw\) за всако \(z\) из \(\mathfrak Z\). Теды
\[
\sum_S zc_{\lambda_S}u_S
 =\sum_S c_{\lambda_S}u_Sz
 =\sum_S c_{\lambda_S}z^{S^{-1}}u_S
\]
теды \(zc_{\lambda_S}=c_{\lambda_S}z^{S^{-1}}\), или \(c_{\lambda_S}(z-z^{S^{-1}})=0\), за всако \(z\) и \(S\) (линеарна независност од \(u\)!). Из тога але следује: \(c_{\lambda_S}=0\) за \(S\neq E\), бо за всако тако \(S\) постојí најманье \emph{једно} \(z\), тако же \(z-z^{S^{-1}}\neq0\); теды \(w\) јест в \(\mathfrak Z\).

\noindent\textbf{Помочна лема 2.} \emph{Всакы \(\mathfrak Z\)-двојны модул \(\mathfrak A\), кторы јест двојно изоморфны с \(\mathfrak Zu_S=u_S\mathfrak Z\) (за \(\mathfrak Z\) како левы и правы оператор), јест идентичны с \(\mathfrak Zu_S\).}

\noindent\emph{Доказ 2.} \(\mathfrak A\) јест изоморфны с \(\mathfrak Zu_S\) како левы модул; ако теды \(u_S\mapsto a\) из \(\mathfrak A\), то \(\mathfrak Zu_S\mapsto\mathfrak Za\), теды \(\mathfrak A=\mathfrak Za\). Из двојного изоморфизма следује (1) \(za=az^S\) (заради \(u_Sz\mapsto az\)); и из тога \(\mathfrak A=\mathfrak Zu_S\). Ако именно положимо: \(a=u_Su_S^{-1}a=u_Sw\), достава се: \(za=zu_Sw=u_Sz^Sw\), але заради (1): \(za=u_Swz^S\), теды (по множеньу с \(u_S\)): \(z^Sw=wz^S\) за всако \(z\) из \(\mathfrak Z\). Зато такоже \(z^S\) пробегаје цело \(\mathfrak Z\), \(w\) јест по елементах комутативны с \(\mathfrak Z\), теды по 1. приналежи к \(\mathfrak Z\).

\noindent\textbf{Помочна лема 3.} \emph{Всакы \(\mathfrak Z\)-двојны модул из \(\mathfrak Z\cdot\mathfrak G\) имаје форму: \(\mathfrak Zu_{S_1}+\cdots+\mathfrak Zu_{S_k}\), кде \(u_{S_1}\) означаје ныкоје \(k\) различне из \(u_S,\ldots,u_T\).}

\noindent\emph{Доказ.} Прекрижены продукт \(\mathfrak Z\cdot\mathfrak G\) јест полно редуцибилны како \(\mathfrak Z\)-двојны модул, бо всакы суманд \(\mathfrak Zu_S\) имаје ранг 1. Како двојне модулы, вси суманды такоже приналежет к различным класам, теды не сут двојно изоморфне. Заради полној редуцибилности всакы \(\mathfrak Z\)-двојны модул \(\mathfrak M\) јест премы суманд и сам полно редуцибилны. \(\mathfrak M=\sum\mathfrak A_{S_i}\), кде \(\mathfrak A_{S_i}\) јест изоморфны с \(\mathfrak Zu_{S_i}\); и зато по 2. јест идентичны с \(\mathfrak Zu_{S_i}\). (Вси \(\mathfrak A_{S_i}\) приналежет к различным класам.)

\noindent\textbf{Помочна лема 4.} \emph{Всако колцо из \(\mathfrak Z\cdot\mathfrak G\), кторо садржује \(\mathfrak Z\), особливо само \(\mathfrak Z\cdot\mathfrak G\), јест двустранно просто. --- Подколца, кторе садржујут \(\mathfrak Z\), једнозначно одповедајут подгрупам од \(\mathfrak G\).}

\noindent\emph{Доказ.} Всако колцо \(\mathfrak o\), кторо садржује \(\mathfrak Z\), јест \(\mathfrak Z\)-двојны модул, теды по 3. имаје форму \(\sum\mathfrak Z\cdot u_{S_i}\); понеже \(\mathfrak o\) јест колцо, продукт двох \(u_{S_i}\) знову приналежи к \(\mathfrak o\), теды такоже приналежечи \(\mathfrak Z\)-модул. Зато \(u_{S_i}\) творят групу до факторных системов, бо јих јест конечно много. Теды \(\mathfrak o=\mathfrak Z\cdot\mathfrak H\), кде \(\mathfrak H\) јест подгрупа од \(\mathfrak G\). (Точније, \(\mathfrak Z\cdot u_{S_i}\) творят мултипликативны систем.) Ако ныне \(\mathfrak a\) јест либовольны двустранны идеал \(\neq0\) из \(\mathfrak o\), то он јест \(\mathfrak o\)-двојны модул, и тым \(\mathfrak Z\)-двојны модул и подмодул од \(\mathfrak o\). Теды вместе с \(\mathfrak Zu_{S_i}\) садржује целу групу \(\mathfrak H\), и зато поставаје \(\mathfrak o=\mathfrak Z\cdot\mathfrak H\).

\noindent\textbf{Помочна лема 5.} \emph{\(P\) јест центар од \(\mathfrak Z\cdot\mathfrak G\).}

\noindent\emph{Доказ.} Ако \(w\) јест по елементах комутативны с \(\mathfrak Z\cdot\mathfrak G\), то особливо с \(\mathfrak Z\), теды по 1. приналежи к \(\mathfrak Z\). Зато из \(wu_S=u_Sw^S\) следује \(w=w^S\), ако јест комутативны. Теды \(w\) приналежи к \(P\).

\noindent\emph{Приметка 1.} Доказ показује, же теорем важи такоже, когда \(\mathfrak Z\) јест бесконечно алгебраично Галоисово полье над \(P\). Бо теоремы о полној редуцибилности оставајут важити (Крулл Захлринге, Матх. Анн. 99). При том \(\mathfrak Z\cdot\mathfrak G\) јест систем всих конечных сум \(\sum_{S_\lambda}u_{S_\lambda}\cdot c_{S_\lambda}\).

Само при 4. треба за \(\mathfrak o\) предположити \(\mathfrak o=\mathfrak Z\cdot\mathfrak H\), что всегда важи за полны продукт \(\mathfrak Z\cdot\mathfrak G\), бо из належеньа од \(u_{S_i}u_{S_k}\) ныне не можно јест закльучити групово својство система, ако \(\mathfrak o\) јест бесконечно над \(\mathfrak Z\).

\noindent\emph{Приметка 2.} При премом основаньу тут треба присојединити теорију прекриженого матричного представјеньа, развијену в §§ 25/26. В следујучем, на концу § 28, употребја се, же иредуцибилно прекрижено представјенье имаје ступень \(t\), кде \(t\) јест абсолутно число компонентов (индекс) од \(\mathfrak K\). Знову треба указати, же та теорија представјеньа даваје идентитет факторных системов, уведеных тут како константы множеньа, с појмом обычным в литературе.
% END BOOK_S27

% BEGIN BOOK_S28
\subsection*{§ 28. Продуктны теорем за факторне системы}\label{produktsatz-fuer-faktorensysteme}

Нека буде
\[
\mathfrak K_r=\mathfrak Z\cdot\mathfrak G
=\mathfrak Zu_E+\cdots+\mathfrak Zu_T,
\qquad u_Su_T=u_{ST}a_{S,T},
\]
\[
\overline{\mathfrak K}_r
=\overline{\mathfrak Z}\,\overline{\mathfrak G}
=\overline{\mathfrak Z}\bar u_E+\cdots+\overline{\mathfrak Z}\bar u_T,
\qquad
\bar u_S\bar u_T=\bar u_{ST}\bar b_{S,T},
\]
(\(\mathfrak Z\) и \(\overline{\mathfrak Z}\) сут изоморфне Галоисове польа, а \(\mathfrak G\), \(\overline{\mathfrak G}\) ньих групы), теды достава се
\[
\mathfrak K_r\times\overline{\mathfrak K}_r
=\widetilde{\mathfrak Z}\,\widetilde{\mathfrak G}\times P_n
=\mathfrak A_f\times P_n
=\mathfrak A_j
\]
с \(\tilde u_S\tilde u_T=\tilde u_{ST}\tilde c_{S,T}\) (\(\widetilde{\mathfrak Z}\) јест изоморфно с \(\mathfrak Z\) и \(\overline{\mathfrak Z}\)). То следује из рачунаньа рангов и из тога, же \(\mathfrak Z\) јест такоже разпадно полье).

\noindent\textbf{Тврдженье.} \emph{Достава се \(\tilde c_{S,T}=\tilde a_{S,T}\tilde b_{S,T}\), гдје \(\tilde a,\tilde b\) сут с \(a\) и \(\bar b\) изоморфне елементи в \(\widetilde{\mathfrak Z}\).}

Доказ основываје се на том, же колцо автоморфизмов либовольного левого идеала из \(\mathfrak A_j\), кторы имаје абсолутну должину \(n\), јест изоморфно с \(\mathfrak A_f=\widetilde{\mathfrak Z}\widetilde{\mathfrak G}\), и же то колцо автоморфизмов можно јест премо определити при погодном левом идеалу из \(\mathfrak A_j\). Обче важи следујучи теорем, именно с \(s\) вместо \(n\), але заради удобности изречены за \(\mathfrak K_r\):

\noindent\textbf{1. Теорем 1.} \emph{Колцо автоморфизмов левого идеала \(\mathfrak l\) абсолутној должины \(n\) из \(\mathfrak K_{rs}=\mathfrak K_r\times P_s\) јест изоморфно с \(\mathfrak K_r\).}

Доказ следује из следујучих помочных лем.

\noindent\textbf{Помочна лема 1.} \emph{Вси леве идеалы абсолутној должины \(n\) из \(\mathfrak K_{rs}\) сут операторно изоморфне. Ньих колца автоморфизмов сут теды колцово изоморфне. Особливо \(\mathfrak l=\mathfrak K_r\times\mathfrak b\) имаје абсолутну должину \(n\), при чем под \(\mathfrak b\) разумејемо просты идеал из \(P_s\).}

\noindent\emph{Доказ.} То, же леве идеалы имајут абсолутну должину \(n\), т.~ј. разпадајут се на \(n\) абсолутно неразложимых простых идеалов, имаје за последицу једнаку должину в \(\mathfrak K_{rs}\), т.~ј. в \(\mathfrak K_{rs}\) наступаје једнако много простых идеалных сумандов. Из тога следује ньих операторны изоморфизм, а тым и колцовы изоморфизм колц автоморфизмов. Далье \(\mathfrak K_r=\mathfrak Z\cdot\mathfrak G\) имаје абсолутну должину \(n\), т.~ј. ранг јест \(n^2\), теды \(\mathfrak K_r\times P_s\) имаје абсолутну должину \(n\cdot s\), и зато \(\mathfrak K_r\times\mathfrak b\) имаје абсолутну должину \(n\).

\noindent\textbf{Помочна лема 2.} \emph{Ако \(\mathfrak o\) јест колцо с јединицоју, \(\mathfrak l=\mathfrak o e_1\) јест левы идеал и премы суманд, теды \(e_1\) јест идемпотент, то колцо автоморфизмов од \(\mathfrak l\) јест изоморфно с \(e_1\mathfrak o e_1\).}

\noindent\emph{Доказ.} Множенье с елементом \(e_1re_1\) даваје операторны хомоморфизм од \(\mathfrak l\) в себе; различне елементи \(e_1re_1\) давајут различне хомоморфизмы, бо при том к \(e_1\) сут приряđене различне елементи, и всакы хомоморфизм јест створјен тако: \(e_1\mapsto re_1\) и \(re_1=e_1re_1\) заради \(e_1e_1=e_1\). Сума и продукт в \(e_1\mathfrak o e_1\) але одповедајут суми и продукту приряđеных хомоморфизмов.

\noindent\textbf{Помочна лема 3.} \emph{Ако положимо \(P_s=\sum_1^s c_{ik}P\) с матричными јединицами \(c_{ik}\) и \(\mathfrak b=c_{11}P+\cdots+c_{s1}P\), то како колцо автоморфизмов од \(\mathfrak K_r\times\mathfrak b\) достава се колцо \(\mathfrak K_rc_{11}\), кторо јест изоморфно с \(\mathfrak K_r\).}

\noindent\emph{Доказ.} Бо \(c_{11}\) јест по елементах комутативны с \(\mathfrak K_r\) и ниједин елемент из \(\mathfrak K_r\) не јест анулованы, бо \(c_{11}\) јест стварјајучи идемпотент од \(\mathfrak K_r\times P_s\), и јест
\[
c_{11}\cdot\mathfrak K_r\times P_s\cdot c_{11}
=\mathfrak K_r\times c_{11}P_sc_{11}
=\mathfrak K_rc_{11}.
\]
Из тых помочных лем премо следује теорем 1.

\noindent\textbf{2. Определьенье и разклад од \(\mathfrak K_r\times\overline{\mathfrak K}_r\) на леве идеалы абсолутној должины \(n\).} По дефиницији достава се
\[
\mathfrak A_f=\mathfrak K_r\times\overline{\mathfrak K}_r
=\mathfrak Z\cdot\mathfrak G\times\overline{\mathfrak Z}\cdot\overline{\mathfrak G},
\]
при чем елемент с чртоју јест по елементах комутативны с елементом без чрты, теды
\[
\mathfrak A_f
=(\mathfrak G\times\overline{\mathfrak G})\cdot(\mathfrak Z\times\overline{\mathfrak Z})
=\mathfrak G\times\overline{\mathfrak G}\cdot(\mathfrak Ze_1+\cdots+\mathfrak Ze_n)
\]
с \(\mathfrak Ze_i=\overline{\mathfrak Z}e_i\), (множенье Галоисового польа самого с собоју). Понеже вси суманды имајут једнакы ранг взходно к \(\mathfrak Z\), а теды такоже једнаку абсолутну должину --- абсолутна должина од \(\mathfrak A_f\) але јест \(=n^2\) --- всакы суманд имаје абсолутну должину \(n\). В следујучем за основу буде взято
\[
\mathfrak l=(\mathfrak G\times\overline{\mathfrak G})\mathfrak Ze_1,
\]
при чем релација \(ze_1=e_1\bar z\) установјаје определьены изоморфизм меджу \(\mathfrak Z\) и \(\overline{\mathfrak Z}\).

\noindent\textbf{3. Прекрижено представјенье од \(e_1\mathfrak A_fe_1\), и тым доказ продуктного теорема.} --- По 1. и 2. факторне системы продукта \(\mathfrak K_r\times\overline{\mathfrak K}_r\) сут изоморфне с факторными системами прекриженого представјеньа од \(e_1\mathfrak A_fe_1=\mathfrak A\).

\noindent\textbf{Теорем 2.} \emph{\(\mathfrak A=\widetilde{\mathfrak Z}\cdot\widetilde{\mathfrak G}\) с \(\widetilde{\mathfrak Z}=\mathfrak Ze_1=\overline{\mathfrak Z}e_1\), \(\widetilde{\mathfrak G}=\{\ldots,e_1u_S\bar u_Se_1,\ldots\}\), при чем
\[
e_1u_S\bar u_Se_1\cdot e_1u_T\bar u_Te_1
=e_1u_{ST}\bar u_{ST}e_1\cdot a_{S,T}e_1\bar b_{S,T}e_1,
\]
тако же с \(\tilde a_{S,T}=a_{S,T}e_1\) и \(\tilde b_{S,T}=\bar b_{S,T}e_1\) продуктны теорем јест доказан; (\(\tilde a\) и \(\tilde b\) сут в том самом польу \(\widetilde{\mathfrak Z}\)).}

\noindent\emph{Доказ.} \(e_1\) поставаје јединица колца; елементи коньуговане с \(\tilde z=ze_1\) сут теды дане посредством \(\tilde z^Se_1\); треба показати, же елементи \(u_S\bar u_Se_1\) индуцирајут те субституције, чиме буде доказано прекрижено представјенье.

\noindent\textbf{Помочна лема 1.} \(e_1u_S\bar u_S=u_S\bar u_Se_1\), теды \(=e_1u_S\bar u_Se_1\).

\noindent\emph{Доказ.} Положимо \(u_S^{-1}e_1u_S=e^S\)\srcfn{1}{Тым установјеньем достава се \(e_1=e^E\).} или \(e_1u_S=u_Se^S\); тада \(e^S\) пробегајут коньуговане \(e_1,e_2,\ldots,e_n\). Тада важи (1)\srcfn{2}{заради \(\sum a_i\bar A_i\bar u_R=\sum a_i\bar u_R\bar A_i^R=\bar u_R\sum(a_i^R)^{R^{-1}}\bar A_i^{-R}\).}
\[
e_1\bar u_R=\bar u_Re^{R^{-1}}.
\]
Бо \(e_1=\sum a_i\bar A_i=\sum a_i^R\bar A_i^R\), кде \(a_i\) пробегаје либовольну \(\mathfrak Z\)-базу, а \(\bar A_i\) приналежечу комплементарну \(\overline{\mathfrak Z}\)-базу, теды то само важи за \(a_i^R\) и \(\bar A_i^R\). Теды \(e_1\bar u_R=\bar u_Re^{R^{-1}}\), чиме (1) јест доказано. Из (1) следује
\[
e_1u_S\bar u_S=u_Se^S\bar u_S=u_S\bar u_S\cdot e^{SS^{-1}}=u_S\bar u_Se_1,
\]
чиме јест доказана помочна лема 1.

\noindent\textbf{Помочна лема 2.} \(ze_1\cdot e_1u_S\bar u_Se_1=e_1u_S\bar u_Se_1\cdot z^Se_1\).

\noindent\emph{Доказ.} Бо
\[
ze_1\cdot e_1u_S\bar u_Se_1
=zu_S\bar u_Se_1\quad(\text{Pomoćna lema 1})
=u_Sz^S\bar u_Se_1
=u_S\bar u_Se_1\cdot z^Se_1
=b=e_1b,
\]
(\(z\) јест комутативны с \(\bar u\) и \(e_1\)). --- Нека \(u'=u_S\bar u_Se_1\). Помочна лема 2 изрека, же \(u'\) индицираје субституције од \(\mathfrak Ze_1\), тако же важи прекрижено представјенье подано в теорему.

\noindent\textbf{Помочна лема 3.} \emph{Факторны систем поставаје равны \(a_{S,T}e_1\cdot\bar b_{S,T}e_1\).}

\noindent\emph{Доказ.}
\[
\begin{aligned}
e_1u_S\bar u_Se_1\cdot e_1u_T\bar u_Te_1
&=u_S\bar u_S\cdot u_T\bar u_Te_1
 =u_{ST}a_{S,T}\cdot\bar u_{ST}\bar b_{S,T}e_1\\
&=u_{ST}\bar u_{ST}\cdot a_{S,T}\bar b_{S,T}e_1
 =e_1u_{ST}\bar u_{ST}e_1\cdot a_{S,T}e_1\cdot\bar b_{S,T}e_1.
\end{aligned}
\]
Тым јест доказан теорем из 2., а теды и продуктны теорем за факторне системы.

\noindent\textbf{4. Продуктны теорем за класы асоциованых факторных системов:}
\[
\{a_{S,T}\}\cdot\{b_{S,T}\}=\{a_{S,T}\cdot b_{S,T}\}.
\]
То премо следује из 3., бо преход к новым стварјајучим елементам \(v\) и \(\bar v\) означаје такоже такы преход к \(v\bar ve_1\), понеже \(u=vc\) и \(\bar u=\bar v\bar d\) давајут
\[
u\bar ue_1=vc\cdot\bar v\bar de_1=v\bar v\cdot c\bar de_1=v\bar ve_1\cdot c\bar de_1.
\]

\noindent\textbf{5.} Ако \(\{a\}\) јест факторны систем од \(\{\mathfrak K\}\), а \(t\) јест индекс од \(\{\mathfrak K\}\), то \(\{a\}^t\) јест равны јединчному класу. То следује из тога, же иредуцибилно прекрижено представјенье од \(\mathfrak G\) имаје ступень \(t\), управо како Сцхуров доказ по представјеньу од \(p_{ik}\) (§ 24, стр. 37).

\noindent\textbf{6.} Изоморфизм меджу групоју \(\mathscr K_{\mathfrak Z}\) од \(\{\mathfrak K\}\) к \(\mathfrak Z\), т.~ј. с \(\mathfrak Z\) како разпадным польем, и групоју клас \(\{a_{S,T}\}\) асоциованых факторных системов. По конструкцији в § 27 постојí једнозначно одповеданье меджу \(\{\mathfrak K\}\) и \(\{a_{S,T}\}\); по 4. приряđенье јест хомоморфно, теды имамо изоморфизм. Из тога особливо следује, же јединчне класы одповедајут једне другим:

\noindent\emph{Прекрижены продукт јест тада и само тада равны \(P\), теды матричному колцу в центру, когда факторны систем јест асоциованы с јединчным системом.}
% END BOOK_S28

% BEGIN BOOK_S29
\subsection*{§ 29. Теорем о главном роду в минималном}\label{hauptgeschlechtssatz-im-minimalen}

Нека \(\mathfrak K_r=\mathfrak Z\cdot\mathfrak G\) буде прекрижены продукт, а \(\mathfrak G^*\) група всих елементов, кторе посредством внутрных автоморфизмов преводет \(\mathfrak Z\) како целост в себе (\(\mathfrak Z^*u_E,\mathfrak Z^*u_S\), --- записано в побочных класах).

\noindent\textbf{Дефиниција.} „\emph{Главны род в минималном}“ \(\mathscr H\) состоји се из всих груповых автоморфизмов од \(\mathfrak G^*\), кторе сут продолженье идентитета на \(\mathfrak Z^*\). Разлог за то названье следује из специализације на цикличне польа. Види § 30.

\noindent\textbf{Теорем.} \emph{Всакы автоморфизм главного рода имаје форму \(u_S\mapsto u_S\cdot b^Sb^{-1}=u_Sb^{S-1}\), кде \(b\) јест в \(\mathfrak Z\), и обратно, всако тако приряđенье стварјаје автоморфизм главного рода.}

\noindent\emph{Доказ.} Нека при груповом автоморфизму буде \(u_S\mapsto v_S\), \(z\mapsto z\), за всако \(z\) из \(\mathfrak Z\), и нека при том продукты одповедајут једни другим; тада автоморфизм можно јест дополнити до автоморфизма колца \(\mathfrak K_r\) посредством приряđеньа
\[
\sum u_Sc_{\lambda_S}\mapsto\sum v_Sc_{\lambda_S},
\]
при чем \(c_{\lambda_S}\) сут в \(\mathfrak Z\). Але всакы такы автоморфизм јест, како знано, внутрны.

Нека \(b\) буде стварјајучи елемент, теды \(v_S=bu_Sb^{-1}\), \(z=bzb^{-1}\); понеже \(b\) јест комутативны со всими \(z\), он лежи в \(\mathfrak Z\), и зато достава се
\[
v_S=bu_Sb^{-1}=u_Sb^Sb^{-1}.
\]
Обратно, ако положимо \(v_S=u_Sb^Sb^{-1}\) с \(b\) в \(\mathfrak Z\), то достава се \(v_S=bu_Sb^{-1}\); иде теды о автоморфизм од \(\mathfrak K_r\), кторы преводи \(\mathfrak Z\) по елементах в себе, а тым по дефиницији преводи такоже \(\mathfrak G^*\) како целост в себе; он јест теды автоморфизм главного рода.
% END BOOK_S29

% BEGIN BOOK_S30
\subsection*{§ 30. Специализација на цикличне разпадне польа}\label{spezialisierung-auf-zyklische-zerfaellungskoerper}

Ако прекрижены продукт образујемо особливо с цикличным польем \(\mathfrak Z\), приходимо к толкованьу знаных нормных теоремов. Најпрво важи

\noindent\textbf{Теорем 1.} \emph{Група \(\mathscr K_{\mathfrak Z}\) од \(\{\mathfrak K\}\) с установјеным цикличным разпадным польем \(\mathfrak Z\) над \(P\) јест изоморфна с групоју \(P^*/N\mathfrak Z^*\), кде \(N\mathfrak Z^*\) пробегаје групу всих норм \(Nz\) елементов \(z\neq0\) из \(\mathfrak Z\).}

\noindent\emph{Доказ.} Доказ основываје се на нормованьу прекриженого представјеньа и на неколико помочных лемах о том.

\noindent\textbf{Дефиниција.} Циклично прекрижено представјенье называје се \emph{нормованым}, ако имаје форму
\[
\mathfrak K_r=\mathfrak Z+\mathfrak Zu+\cdots+\mathfrak Zu^{n-1},
\]
с \(zu=uz^S\), \(u^n=a\), кде \(S\) јест стварјајучи елемент из \(\mathfrak G\) (ако теды положимо \(u_{S^r}=u^r\), \(r=0,\ldots,n-1\), вместо обчего \(u_{S^r}=b_ru^r\) с \(b_r\) в \(\mathfrak Z\)).

\noindent\textbf{Помочна лема 1.} \emph{Циклично прекрижено представјенье можно јест всегда взяти нормованым. При том \(a\) лежи в \(P\), и всако \(a\) из \(P\) даваје тако представјенье.}

\noindent\emph{Доказ.} Прва чест јест јасна, бо из \(u^{-1}zu=z^S\) следује \(u^{-r}zu^r=z^{S^r}\), теды \(u^r\) стварјаје субституцију \(z\mapsto z^{S^r}\). Далье, понеже \(u^n\) индуцираје идентитет \(z\mapsto z^{S^n}\), достава се \(u^n=a\) с \(a\) в \(\mathfrak Z\). Але \(u^n\) јест комутативны со всими \(u^r\), то само важи за \(a\), и теды \(a^{S^r}=a\); зато \(a\) лежи в \(P\).

Теды \(a\) јест једина сучствена константа множеньа. Зато условја асоциативности не имајут содржанье; всако \(a\) из \(P\) даваје асоциативны систем, а тым и неко \(\mathfrak K_r\).

\noindent\textbf{Помочна лема 2.} \emph{Вместо полного класа \(\{a_{S,T}\}\) асоциованых факторных системов достаточно јест разматрјати нормоване факторне системы, кторе сут садржане в класу --- т.~ј. кторе произходет из нормованых представјень. Ту подмножину класа называмо нормованым класом. Група полных клас \(\{a_{S,T}\}\) јест изоморфна с групоју нормованых клас.}

\noindent\emph{Доказ.} Бо по помочној леми 1 ниједин нормованы клас не јест праздны, приряđенье јест теды једнозначно. Операција оставаје иста, бо јест определьена посредством либовольного представительа.

\noindent\textbf{Помочна лема 3.} \emph{Група нормованых факторных системов јест изоморфна с мултипликативној групоју \(P^*\), а група нормованых клас јест изоморфна с \(P^*/N\mathfrak Z^*\).}

\noindent\emph{Доказ.} Из помочној лемы 1 следује, же всакы нормованы факторны систем имаје вид \([1,1,\ldots,a,a,\ldots]\), бо
\[
u^\nu\cdot u^\mu=u^{\nu+\mu}\quad(\nu+\mu<n),
\qquad
u^\nu\cdot u^\mu=au^\varrho\quad(\nu+\mu=n+\varrho,\ 0\leq\varrho\leq n-2).
\]
Але приряđенье
\[
[1,\ldots,a,a,\ldots]\mapsto a,
\qquad
[1,1,\ldots,b,b,\ldots]\mapsto b,
\]
\[
[1,\ldots,ab,\ldots]\mapsto ab,
\]
јест груповы изоморфизм на \(P^*\), бо, знову по помочној леми 1, \(a\) пробегаје все елементы из \(P^*\). --- Два асоциоване нормоване факторне системы произходет једен из другого посредством трансформације \(v=uc\) (\(c\) в \(\mathfrak Z\)). То але даваје
\[
v=uc,\qquad v^2=u^2\cdot c^S\cdot c;
\]
\[
v^r=u^r c^{S^{r-1}}c^{S^{r-2}}\cdots c;
\qquad
v^n=u^nN(c).
\]
Из \(u^n=a\) теды следује \(v^n=aN(c)\). Нормованы клас состоји се теды из факторных системов \([1,\ldots,aN(c),aN(c),\ldots]\), кде \(c\) пробегаје все елементы \(\neq0\) из \(\mathfrak Z\). Тым јест доказана друга чест помочној лемы.

Из помочных лем 2 и 3 следује горньи теорем 1, с обзиром на изоморфизм тамошньих \(\{\mathfrak K\}\) с класами \(\{a_{S,T}\}\).

\noindent\textbf{Теорем 2.} \emph{Група главного рода в минималном јест изоморфна с групоју тых елементов \(c\) из \(\mathfrak Z\), за кторе \(N(c)=1\). Из \(N(c)=1\) теды следује постојанье некого \(b\neq0\) из \(\mathfrak Z\), тако же \(c=b^S\cdot b^{-1}=b^{S-1}\) (симболична \(S-1\)-та потенција).}

\noindent\emph{Доказ.} При доказу помочној лемы 3 достало се: \(N(c)=1\) јест идентично с тым, же трансформација \(v=uc\) представльа автоморфизм главного рода; различным \(c\) одповедајут различне автоморфизмы и обратно, а продукту одповедаје продукт.

\noindent\textbf{Приметка.} Теорем 2 оправдава названье главны род. По смыслу факторгрупу
\[
\mathfrak Z^*/\mathscr H
\quad(\mathscr H=\text{glavny rod})
\cong\mathfrak Z^*/\mathfrak Z^{*S-1}
\]
можно јест назвати групоју родов в минималном; за ту групу важи, же она јест изоморфна с \(N(\mathfrak Z^*)\), теды в свези с теоремом 1,
\[
\begin{aligned}
\text{grupa rodov v minimalnom}
&=\mathfrak Z^*/\mathscr H
 \cong\mathfrak Z^*/\mathfrak Z^{*S-1}
 \cong N(\mathfrak Z^*)\\
&=\text{grupa od }\{\mathfrak K\}\text{ k }\mathfrak Z
 =\mathscr K_{\mathfrak Z}
 \cong P^*/N(\mathfrak Z^*).
\end{aligned}
\]
% END BOOK_S30

% BEGIN BOOK_S31
\subsection*{§ 31. Применьеньа цикличного специалного случаја}\label{anwendungen-des-zyklischen-spezialfalles}

Циклично прекрижено представјенье даваје нове доказы за теоремы доказане в § 20: же всако конечно тело јест комутативно и же кроме кватернионов не постојí никако право некомутативно разширјенье польа реалных чисел. --- Те доказы можно бы было назвати доказами теорије класовых поль в минималном.

\noindent\textbf{1. Всако тело из конечного числа елементов јест комутативно (теорем 10 § 20).}

Доказ основываје се на двох знаных фактах о комутативных конечных польах (Галоисовых польах).

\noindent\textbf{1.} \emph{Галоисове польа сут цикличне над всакым својим подпольем.}

\noindent\textbf{2.} \emph{Ако \(\mathfrak Z\) јест Галоисово полье, теды \(\mathfrak Z\) јест циклично над \(P\), то всакы елемент из \(P\) јест норма некого елемента од \(\mathfrak Z\). \(P^*=N(\mathfrak Z^*)\).}

Нека ныне \(\mathfrak K\) буде конечно, \(P\) јего центар, а \(\mathfrak Z\) неко циклично разпадно полье; тада (§ 30 конец), меджу другым, група \(\mathscr K_{\mathfrak Z}\) од \(\{\mathfrak K\}\) к \(\mathfrak Z\) јест изоморфна с \(P^*/N\mathfrak Z^*\); теды по 2. она јест равна јединици \(\{P\}\). То але значи, же не постојí од \(P\) различно тело с \(P\) како центром, \(\mathfrak K=P\).

\noindent\textbf{2. Једино право некомутативно разширјенье над польем \(P\) реалных чисел сут кватернионы (теорем 9 § 20).}

\noindent\emph{Доказ.} \(P\) мора бити центар, бо \(P\) имаје како комутативно разширјенье само полье \(\mathfrak Z\) комплексных чисел, над кторым како центром, понеже \(\mathfrak Z\) јест алгебраично затворјено, не постојí некомутативно разширјенье тела. Једночасно само \(\mathfrak Z\) приходи в обзир како разпадно полье. И \(\mathfrak Z\) јест циклично над \(P\). Далье, всако позитивно число јест норма, теды група \(P^*/N\mathfrak Z^*\) имаје ступень 2; зато кроме \(\{P\}\) постојí само једен клас \(\{\mathfrak K\}\), а то сут кватернионы. Једночасно достава се нормална форма, ако взмемо \(-1\) како нормованы факторны систем в класу асоциованых системов. Бо тада достава се \(\mathfrak K=\mathfrak Z+\mathfrak Zu\), с \(u^2=-1\). Теды \(\mathfrak K=P+Pi+(P+Pi)u\) с \(i^2=-1\), \(u^2=-1\).

\noindent\textbf{Приметка к 1.} Такоже доказ знаного факта 2 можно јест на основе конца § 30 водити тако: Достава се \(\mathscr H=\mathfrak Z^{*S-1}\cong\mathfrak Z^*/P^*\), бо все и само елементи од \(P^*\) стварјајут идентичны автоморфизм главного рода, или такоже при применьеньу од \(S-1\) преходет в јединицу. Теды при конечном \(\mathfrak Z\) число елементов од \(\mathscr H\) јест равно индексу од \(P^*\) в \(\mathfrak Z^*\). Зато число елементов од \(\mathfrak Z^*/\mathscr H\) јест равно числу елементов од \(P^*\); то само теды важи за \(N(\mathfrak Z^*)=\mathfrak Z^*/\mathscr H\), и зато \(P^*=N(\mathfrak Z^*)\). (Сравни: число родов = число амбигвных клас. Тому числово-теоретичному начину изражаньа але всуде треба додати „в минималном“.)

\noindent\textbf{3.} \emph{Ако \(P\) јест \(p\)-адично основно полье, а \(\mathfrak Z\) јест циклично \(n\)-того ступеньа над \(P\), то група \(\mathscr K_{\mathfrak Z}\) од \(\{\mathfrak K\}\) к \(\mathfrak Z\) јест изоморфна с Галоисовој групоју од \(\mathfrak Z\).}

То јест по концу § 30 према последица „теорије класовых поль в малом“, ктора изрека
\[
P^*/N\mathfrak Z^*\cong\mathfrak G.
\]

\noindent\textbf{Приметка к 3.} Следује, же груповы ред елементов из \(\mathscr K_{\mathfrak Z}\), теды експонент од \(\{\mathfrak K\}\), поставаје равны јего индексу. Бо експонент стварјајучего елемента од \(\mathscr K_{\mathfrak Z}\) поставаје \(=n\), теды такоже индекс како делитель од \(n\) и многократник експонента; група \(\mathscr K_{\mathfrak Z}\) садржује теды само класы, кторых индекс јест делитель од \(n\), и за кторе \(\mathfrak Z\) јест разпадно полье. Надто H. Хассе показал, Матх. Анн. 104, же \(\mathscr K_{\mathfrak Z}\) состоји се из всих клас \(\{\mathfrak K\}\), за кторе индекс јест делитель од \(n\). Всако циклично разпадно полье \(\mathfrak Z\) \(n\)-того ступеньа над \(p\)-адичным \(P\) води теды к тој самој групе \(\mathscr K_{\mathfrak Z}\) клас тел \(\{\mathfrak K\}\).

Два факта:

\noindent\textbf{1.} група всих \(\{\mathfrak K\}\) над \(P\), кторых индекс јест делитель од \(n\), јест циклична \(n\)-того ступеньа.

\noindent\textbf{2.} Всако циклично разпадно полье \(n\)-того ступеньа над \(P\) стварјаје полну групу од \(\{\mathfrak K\}\).

можно јест сматрјати сучственым содржаньем обратного теорема теорије класовых поль в малом.


\clearpage
% END BOOK_S31

\end{document}
